% =============================================================================================
% main.tex -- paper skeleton. Owner of the preamble and notation macros: outline agent.
% Notation reference: paper/notation.md. Section plan and label registry: paper/OUTLINE.md.
% Compile from paper/:  latexmk -pdf -interaction=nonstopmode main.tex
% Section writers edit only their own file under sections/. New macros go here (ask the
% outline owner) so that every section uses the same notation.
% =============================================================================================
\documentclass[11pt,a4paper]{article}

% --- Fonts and page --------------------------------------------------------------------------
\usepackage[T1]{fontenc}
\usepackage{lmodern}          % scalable T1 fonts (the bitmap cm-super T1 font is not installed)
\usepackage{microtype}
\usepackage[margin=2.5cm]{geometry}

% --- Mathematics -----------------------------------------------------------------------------
\usepackage{amsmath,amssymb,amsthm,mathtools,bm}

% --- Bibliography ----------------------------------------------------------------------------
\usepackage[authoryear,round]{natbib}

% --- Tables and figures ----------------------------------------------------------------------
\usepackage{booktabs}
\usepackage{threeparttable}
\usepackage{siunitx}
\sisetup{detect-all, group-separator={,}, group-minimum-digits=4}
\usepackage{graphicx}
\graphicspath{{figures/}}
\usepackage{adjustbox}        % \fittable: shrink a generated tabular to the text width if needed
\usepackage[font=small,labelfont=bf,skip=4pt]{caption}
\usepackage{xcolor}

% --- Cross-references (cleveref must be loaded last) -----------------------------------------
\usepackage{hyperref}
\hypersetup{colorlinks=true, linkcolor=blue!50!black, citecolor=green!35!black,
            urlcolor=blue!60!black, pdftitle={Is Synthetic Financial Data Good Enough?}}
\usepackage{aliascnt}         % one alias counter per theorem type (see below)
\usepackage[capitalise,noabbrev]{cleveref}

% =============================================================================================
% Theorem environments. One counter for all numbered statements, reset in each section
% (Proposition 4.1, Remark 4.2, ...; Lemma A.1 in the appendix). Each environment gets an alias
% of the shared counter (aliascnt), so that \cref prints the right name ("Proposition 4.1",
% not "Theorem 4.1"; cleveref alone does not tell shared counters apart with this TeX Live).
% Write heuristics as remarks whose title says "(heuristic)".
% =============================================================================================
\theoremstyle{plain}
\newtheorem{theorem}{Theorem}[section]
\crefname{theorem}{Theorem}{Theorems}
% \sharedthm[style]{env}{Name}{Names}: numbered like theorem, cross-referenced as Name.
\newcommand{\sharedthm}[4][plain]{%
  \theoremstyle{#1}\newaliascnt{#2}{theorem}\newtheorem{#2}[#2]{#3}\aliascntresetthe{#2}%
  \crefname{#2}{#3}{#4}}
\sharedthm{proposition}{Proposition}{Propositions}
\sharedthm{lemma}{Lemma}{Lemmas}
\sharedthm{corollary}{Corollary}{Corollaries}
\sharedthm{fact}{Fact}{Facts}                       % standard result quoted with a reference
\sharedthm[definition]{definition}{Definition}{Definitions}
\sharedthm[definition]{assumption}{Assumption}{Assumptions}
\sharedthm[definition]{example}{Example}{Examples}
\sharedthm[remark]{remark}{Remark}{Remarks}
\theoremstyle{plain}

\crefname{equation}{equation}{equations}
\Crefname{equation}{Equation}{Equations}
\crefname{appendix}{Appendix}{Appendices}       % \crefalias{section}{appendix} after \appendix

\numberwithin{equation}{section}

% =============================================================================================
% Editorial macros
% =============================================================================================
% \todo{...}: visible red placeholder for anything not yet available (numbers from running
% experiments, missing references, decisions). Works in text, math and captions.
\newcommand{\todoinner}[1]{\textcolor{red}{\textbf{[TODO:}~#1\textbf{]}}}
\DeclareRobustCommand{\todo}[1]{\ifmmode\text{\todoinner{#1}}\else\todoinner{#1}\fi}

% Generated tables and figures. They may not exist yet: these wrappers print a \todo instead of
% stopping the compilation. Templates: paper/OUTLINE.md, Section 11.
%   \inputtable{e2_llr}    -> tables/e2_llr.tex (a bare tabular written by the experiments);
%                             use it inside threeparttable when the table fits the text width
%   \fittable{e3_power}    -> the same, shrunk to \linewidth if it is wider (never enlarged).
%                             NOT inside threeparttable (it cannot measure a boxed tabular):
%                             put the notes in \tablenote{...} instead
%   \tablenote{text}       -> table notes in a \linewidth minipage, \footnotesize
%   \inputnote{e2_size_note} -> tables/e2_size_note.tex (generated notes), inside
%                             \tablenote{} or a tablenotes \item[]
%   \inputfigure[width]{e2_power} -> figures/e2_power.pdf (default width \linewidth)
\newcommand{\inputtable}[1]{%
  \IfFileExists{tables/#1.tex}{\input{tables/#1.tex}}{\todo{table \texttt{\detokenize{#1}.tex} not yet generated}}}
\newcommand{\fittable}[1]{%
  \IfFileExists{tables/#1.tex}%
    {\adjustbox{max width=\linewidth}{\input{tables/#1.tex}}}%
    {\todo{table \texttt{\detokenize{#1}.tex} not yet generated}}}
\newcommand{\tablenote}[1]{%
  \par\smallskip\begin{minipage}{\linewidth}\footnotesize\raggedright #1\end{minipage}}
\newcommand{\inputnote}[1]{%
  \IfFileExists{tables/#1.tex}{\input{tables/#1.tex}}{\todo{note \texttt{\detokenize{#1}.tex} not yet generated}}}
\newcommand{\inputfigure}[2][\linewidth]{%
  \IfFileExists{figures/#2.pdf}{\includegraphics[width=#1]{#2}}%
    {\fbox{\parbox{0.9\linewidth}{\centering\todo{figure \texttt{\detokenize{#2}.pdf} not yet generated}}}}}

% =============================================================================================
% Notation macros (meaning of every macro: paper/notation.md). Use them; do not retype.
% =============================================================================================
% --- Probability, analysis -------------------------------------------------------------------
\newcommand{\E}{\mathbb{E}}
\newcommand{\Prob}{\mathbb{P}}
\DeclareMathOperator{\Var}{Var}
\DeclareMathOperator{\Cov}{Cov}
\DeclareMathOperator{\Corr}{Corr}
\DeclareMathOperator{\Law}{Law}
\newcommand{\W}{W_{1}}                          % Wasserstein-1 distance: \W(P,Q)
\DeclareMathOperator{\Lip}{Lip}                 % Lipschitz constant / class: \Lip(f), \Lip_c
\DeclareMathOperator{\KL}{KL}
\DeclareMathOperator{\TV}{TV}
\DeclareMathOperator{\MMD}{MMD}
\newcommand{\Jeff}{J}                           % Jeffreys divergence KL(P||Q)+KL(Q||P)
\DeclareMathOperator*{\argmin}{arg\,min}
\DeclareMathOperator*{\argmax}{arg\,max}
\DeclareMathOperator{\diag}{diag}
\DeclareMathOperator{\Tr}{tr}                   % trace
\DeclareMathOperator{\med}{med}                 % median
\DeclareMathOperator{\cond}{cond}               % condition number
\DeclareMathOperator{\supp}{supp}
\newcommand{\one}{\mathbf{1}}                   % indicator \one\{A\} and vector of ones \one
\newcommand{\R}{\mathbb{R}}
\newcommand{\Nat}{\mathbb{N}}                   % (N is the number of assets)
\newcommand{\Pspace}[1]{\mathcal{P}_{#1}}       % measures with finite #1-th moment
\newcommand{\iid}{i.i.d.}
\newcommand{\eqd}{\overset{d}{=}}
\newcommand{\dto}{\rightsquigarrow}             % convergence in distribution
\newcommand{\pto}{\xrightarrow{\ \Prob\ }}      % convergence in probability
\newcommand{\indep}{\perp\!\!\!\perp}
\DeclarePairedDelimiter{\norm}{\lVert}{\rVert}
\DeclarePairedDelimiter{\abs}{\lvert}{\rvert}
\newcommand{\Gsig}{\mathcal{G}}                 % sigma-field of the training data and randomness

% --- Data, samples, periods ------------------------------------------------------------------
\newcommand{\TRAIN}{\textsc{train}}             % 2018-01-02 to 2022-12-30 (calibration, real past)
\newcommand{\TEST}{\textsc{test}}               % 2023-01-03 to 2025-06-16 (evaluation, real new)
\newcommand{\Ttr}{T_{\mathrm{tr}}}              % TRAIN length in trading days (1258)
\newcommand{\Tte}{T_{\mathrm{te}}}              % TEST length in trading days (615)
\newcommand{\Rpast}{R_{\mathrm{past}}}          % real windows of TRAIN
\newcommand{\Rnew}{R_{\mathrm{new}}}            % real windows of TEST
\newcommand{\Ppast}{P_{\mathrm{past}}}          % window law of R_past
\newcommand{\Pnew}{P_{\mathrm{new}}}            % window law of R_new
\newcommand{\QG}{Q_{G}}                         % window law of the generator G

% --- Train / evaluation split ----------------------------------------------------------------
\newcommand{\trn}{\mathrm{tr}}                  % superscript/subscript "training half"
\newcommand{\evl}{\mathrm{ev}}                  % superscript/subscript "evaluation half"
\newcommand{\Itr}{I_{\trn}}
\newcommand{\Iev}{I_{\evl}}
\newcommand{\Jtr}{J_{\trn}}
\newcommand{\Jev}{J_{\evl}}
\newcommand{\Pev}{P^{\evl}}                     % empirical mean over real eval windows: \Pev f
\newcommand{\Qev}{Q^{\evl}}
\newcommand{\Ptr}{P^{\trn}}
\newcommand{\Qtr}{Q^{\trn}}

% --- Critic, statistics, p-values ------------------------------------------------------------
\newcommand{\fhat}{\widehat{f}}                 % trained critic (network o standardisation)
\newcommand{\ghat}{\widehat{g}}                 % network part, \fhat = \ghat \circ A
\newcommand{\lhat}{\widehat{\ell}}              % learned logit (LR-critic, C2ST)
\newcommand{\lstar}{\ell^{*}}                   % Bayes-optimal logit
\newcommand{\lgp}{\lambda_{\mathrm{GP}}}        % gradient-penalty weight (10)
\newcommand{\Dcrit}{D}                          % critic discrepancy D(f) = Pf - Qf (target of T)
\newcommand{\Tlogit}{T_{\mathrm{logit}}}
\newcommand{\Tacc}{T_{\mathrm{acc}}}
\newcommand{\TMMD}{T_{\mathrm{MMD}}}
\newcommand{\Tsig}{T_{\mathrm{sig}}}
\newcommand{\TKS}{T_{\mathrm{KS}}}
\newcommand{\TLR}{T_{\mathrm{LR}}}              % Neyman-Pearson oracle statistic
\newcommand{\Danc}{\Delta}                      % anchored relative statistic (primary)
\newcommand{\Dtwo}{\Delta^{\mathrm{2c}}}        % two-critic relative statistic (secondary)
\newcommand{\pperm}{p_{\mathrm{perm}}}
\newcommand{\phac}{p_{\mathrm{hac}}}
\newcommand{\pwild}{p_{\mathrm{wild}}}
\newcommand{\pmed}{p_{\mathrm{med}}}            % min(1, 2 x median p over synthetic seeds)
\newcommand{\zhac}{z_{\mathrm{hac}}}
\newcommand{\LRV}{\Omega}                       % long-run variance of a sequence of critic values
\newcommand{\bNP}{\beta_{\mathrm{NP}}}          % power of the Neyman-Pearson oracle
\newcommand{\Sig}{\mathrm{Sig}}                 % truncated path signature

% --- Decision rule ---------------------------------------------------------------------------
\newcommand{\GOOD}{\textsf{GOOD}}
\newcommand{\BAD}{\textsf{BAD}}
\newcommand{\INCONCL}{\textsf{INCONCLUSIVE}}
\newcommand{\INVALID}{\textsf{INVALID}}
\newcommand{\Cdesign}{\mathcal{C}_{\mathrm{ref}}}  % reference design of the power check
\newcommand{\abd}{\alpha_{\mathrm{bd}}}         % boundary size of the relative test
\newcommand{\bmin}{\beta_{\min}}                % required power (0.8)
\newcommand{\betaL}{\beta_{L}}                  % Clopper-Pearson lower bound of the power

% --- Tests and models (names in text) --------------------------------------------------------
\newcommand{\WC}{W-critic}                      % write \WC{} test
\newcommand{\LRC}{LR-critic}
\newcommand{\SigMMD}{Sig-MMD}
\newcommand{\GARCHt}{GARCH(1,1)-$t$}

% --- Portfolios ------------------------------------------------------------------------------
\newcommand{\wgmv}{w_{\mathrm{gmv}}}
\newcommand{\wmv}{w_{\mathrm{mv}}}
\newcommand{\CE}{\mathrm{CE}}
\newcommand{\SR}{\mathrm{SR}}
\newcommand{\SigLW}{\widehat{\Sigma}_{\mathrm{LW}}}

% =============================================================================================
\title{Is Synthetic Financial Data Good Enough?\\[0.3em]
  \large A Validated Critic Two-Sample Test with Likelihood Benchmarks and Markowitz Utility}
\author{\todo{Author names and affiliations}}
\date{Draft of \today}

\begin{document}

\maketitle

\begin{abstract}
% =============================================================================================
% ABSTRACT (input inside \begin{abstract}...\end{abstract} in main.tex; no \section here).
% Target: 180-220 words, one paragraph, no citations, no numbers that are not in a generated
% table. Plan: paper/OUTLINE.md, "Abstract".
%
% Planned content, in this order:
%  1. Problem: synthetic return data is used for back-testing, data sharing and portfolio
%     construction, but its quality is judged by stylized facts or uncalibrated scores.
%  2. Method: the critic of a Wasserstein GAN, trained alone on a training split of real versus
%     synthetic windows, gives a held-out two-sample statistic; exact permutation level for
%     independent windows; time-block split + Driscoll-Kraay HAC inference for dependent panels;
%     an anchored relative test (is the generator worse than the real past?) that is exact at its
%     boundary; a pre-specified power check and a three-valued verdict (GOOD / BAD /
%     INCONCLUSIVE).
%  3. Theory: exact level, consistency under idealised training, Neyman-Pearson upper bound and
%     the local-power factor, anti-conservativeness of the two-critic relative test, bias of
%     ticker splits on common dates.
%  4. Simulation: valid size everywhere with independent windows; window tests detect volatility
%     clustering; the pooled KS test is more powerful against marginal deviations but invalid
%     under within-window dependence and blind to dependence-only deviations. State plainly that
%     the critic is not uniformly most powerful.
%  5. Application: S&P 500 (2018-2025), five univariate generators, verdicts, agreement with
%     likelihood benchmarks; Markowitz portfolios trained on synthetic data and whether the
%     verdicts predict utility (they need not: Section 8). Headline numbers: \todo until the
%     E3/E4 tables are final.
% =============================================================================================
\todo{Abstract to be written last, from the final tables (see OUTLINE.md).}

\end{abstract}

\noindent\textbf{Keywords:} two-sample test; Wasserstein critic; synthetic data; permutation
test; HAC inference; likelihood ratio; generative models; portfolio choice.\\
\noindent\textbf{JEL codes:} C12, C14, C15, C45, C52, C58, G11.
\todo{confirm keywords and JEL codes}

\section{Introduction}\label{sec:intro}
% =============================================================================================
% Owner: front writer. Plan and claims: paper/OUTLINE.md, Section 3. Notation: notation.md. No
% tables or figures in this section. Every number is taken from a table that a result section
% inputs: tab:e2-size, tab:e2-power, tab:e2dep-power (fig:e2dep-power), tab:e2-llr
% (manual_e2_llr_main), tab:e2-relative-boundary (manual_e2_relative_boundary_main),
% tab:e2-dependent-size, tab:e1-size-sim (manual_e1_size_sim_main), tab:e1-nonstat
% (e1_nonstat), tab:e3-stylized, tab:e3-power (manual_e3_power_main), tab:e3-pooled-tests,
% tab:e3-verdicts, tab:e3-agreement, tab:e4-correlation, tab:e4-utility, tab:e4-utility-agg,
% tab:e4-tests, tab:e4-link. The 45% figure is the theory.md pilot (Section 11.5), quoted as a
% pilot only. The count "all 14 configurations" (critic ranking of the likelihood models in the
% E5 grid) was checked on results/e5/e5_sensitivity.csv (median absolute statistic).
% Post-hoc architecture ablation: tab:e5-arch (e5_arch_sim.tex; 0.98 / 0.10 / 1.00 for t10 vs t4,
% n=500); "pilot on four E2 designs" and "D2 and seed combination added on the day of the runs"
% from the docs/DESIGN.md change log; "six tickers" from results/e3_quick/meta (development runs).
% "Barely reproducible" (E4 critic statistic): seed-0 vs seed-1 Spearman over the 30 universes
% within each generator, -0.03 to 0.30 (tab:manual-e4-reliability, from results/e4/e4_tests.csv).
% Pooled vs replica-wise variance ratios (MV Gaussian GMV 1.038 / 1.004): tab:e4-utility-agg.
% Fixed realised-volatility statistic (0.69-0.94 of the critic's z_hac, same ranking of GBM,
% Merton, GARCH-t):
% tab:e3-rv-baseline. Indep-cell oversize vs standard error: tab:e1-sim-se (sec:sim-dependent).
% =============================================================================================

Synthetic return data, simulated from a fitted model or from a trained generative network, now
serve several purposes in finance. They are used to back-test trading strategies on many paths
instead of the single historical one, which a strategy selected in-sample tends to overfit
\citep{bailey2017probability}; to share data that cannot be disclosed for confidentiality or
privacy reasons \citep{assefa2020generating}; to train hedging strategies
\citep{buehler2019deep,kondratyev2019market}; and to augment short samples, for instance when
portfolio inputs are estimated \citep{buehler2020data,cetingoz2026synthetic}. The generators
range from classical parametric models (geometric Brownian motion, jump-diffusions, GARCH,
stochastic volatility) to generative networks
\citep{wiese2020quant,takahashi2019modeling,buehler2020data,ni2021sig}. Each of these uses
assumes an answer to the question in our title: is the synthetic data good enough?

The question is usually answered informally. Most papers compare stylized facts of real and
generated paths, such as heavy tails, volatility clustering and leverage
\citep{cont2001empirical}, visually or through summary statistics, and evaluation practice varies
widely \citep{stenger2024evaluation}. Discriminator-based scores, such as the held-out error of a
classifier trained to tell real from synthetic series, are used without a null distribution
\citep{yoon2019time,francopedroso2018ets,storchan2021learning}. Without a null distribution and a
known power, a low score cannot be told apart from a weak classifier. Formal two-sample tests have
been used
to evaluate market generators \citep{buehler2020data,chevyrev2022signature}, but without an
analysis of their size and power on financial data. Different evaluation criteria for generative
models can also disagree sharply \citep{theis2016note}, so a criterion has to answer a precise
statistical question and, ideally, say something about the intended use.

We treat the question as a hypothesis test. A yes/no answer that a statistician can defend needs
five things.
\begin{enumerate}
  \item[(i)] \emph{Level control on dependent data.} Real return windows are serially dependent
  through volatility clustering and cross-sectionally dependent through common market shocks, so
  permutation p-values, which assume exchangeability, are not valid for them. With windows of many
  stocks pooled and the training and evaluation halves split by ticker, a pilot simulation of a
  stylised one-factor panel rejected a true null in about 45\% of replications at nominal 5\%. In
  the full placebo of E1, a factor panel fitted to our data and run through the complete pipeline,
  the same protocol rejected in all 200 replications (\cref{tab:e1-size-sim}).
  \item[(ii)] \emph{Power against the deviations that matter in finance}, quantified before the
  test is applied to real data. Without known power, a non-rejection says nothing about quality.
  \item[(iii)] \emph{A null that respects non-stationarity.} A generator $G$ is calibrated on the
  real past $\Rpast$ and judged against later real data $\Rnew$, and it inherits features of the
  past that need not persist. In our data every generator fitted to 2018--2022 is more volatile
  than the 2023--2025 returns (panel median daily standard deviation between 1.89\% and 2.13\%,
  against 1.62\%; \cref{tab:e3-stylized}). The absolute null $\Law(G)=\Law(\Rnew)$ asks the
  generator to anticipate such changes. Users who back-test or manage risk need exactly that, and
  we report this test, but its rejection does not separate a poor generator from a change of
  regime after the calibration period. As in forecast comparisons against a benchmark model
  \citep{diebold1995comparing,giacomini2006tests}, we therefore also compare $G$ with the naive
  benchmark of replaying history, that is, historical simulation from $\Rpast$: is $G$ further from
  $\Rnew$ than $\Rpast$ is? A rejection then means worse than that benchmark, not merely different
  from the future.
  \item[(iv)] \emph{A link to likelihood-based evaluation.} Econometricians trust it, and where a
  likelihood exists it provides a benchmark verdict and a power envelope.
  \item[(v)] \emph{Evidence on use}: whether the verdict predicts how useful the data are for a
  given decision.
\end{enumerate}

This paper builds such a test around the critic of a Wasserstein GAN, trained alone, with no
generator. A critic is a neural-network test function chosen to maximise a mean difference
between two samples under a Lipschitz penalty \citep{arjovsky2017wasserstein,gulrajani2017improved}.
Let $X_1,\dots,X_n$ be windows of $L$ consecutive daily returns of real data, with law $P$, and let
$Y_1,\dots,Y_m$ be windows simulated from $G$, with law $Q$. We split both samples into a training
half and an evaluation half and train the critic $\fhat$ on the training halves. The statistic for
the absolute null $P=Q$ (\cref{def:abs-null}) is the held-out mean difference
\begin{equation*}
  T=\Pev\fhat-\Qev\fhat=\frac{1}{n_e}\sum_{i\in\Iev}\fhat(X_i)-\frac{1}{m_e}\sum_{j\in\Jev}\fhat(Y_j).
\end{equation*}
Given the training half, $T$ is unbiased for $P\fhat-Q\fhat$. This quantity is bounded by the
Lipschitz constant of the network times the Wasserstein-1 distance between the standardised
window laws (\cref{prop:unbiased}). We read $T$ as a learned moment that can be evaluated window
by window (\cref{sec:emp-witness}), not as an estimate of $\W$.

For an econometrician the construction is a learned, sample-split relative of the consistent
specification tests of \citet{bierens1990consistent} and \citet{stinchcombe1998consistent}: they
use a whole class of test functions, whereas we select one on a training split and test with it
on the other. It is also close in spirit to the simulation-based specification tests of
\citet{corradi2005bootstrap}, with one difference that matters. Their null is that the model class
is correctly specified. Ours is that the fitted generator, with parameters estimated on a
calibration period and then held fixed, reproduces the window law of a later evaluation period
(\cref{sec:lit-econometrics}). None of the building blocks is new on its own, and
\cref{sec:lit-summary} lists what we take from earlier work.

\paragraph{Contributions.}
\begin{enumerate}
  \item \emph{A critic test with known level.} For independent windows, the permutation test with
  the critic held fixed has exact finite-sample level whatever the training did
  (\cref{prop:exact-level}). For dependent panels we split by contiguous time blocks rather than
  by ticker, so that training and evaluation windows cover disjoint dates, and we studentise $T$
  with a Driscoll--Kraay variance \citep{driscoll1998consistent}. The resulting p-value is
  asymptotically valid under mixing conditions for a fixed critic whose evaluation data are
  independent of its training data (\cref{prop:hac}). A ticker split on common dates would
  instead couple the two halves through the realised market path, and the statistic is then
  biased under the null (\cref{prop:ticker-bias}). Our time-block split interleaves training and
  evaluation blocks without a gap, so it does not deliver that independence
  (\cref{lem:decoupling}), and we measure the size of the complete real-data pipeline in a
  simulated placebo panel calibrated to the data: at nominal 5\% it rejects a true null at 0.085,
  against 1.000 for the ticker split with permutation p-values, but it over-rejects when the
  synthetic tickers are independent, as they are for our univariate generators
  (\cref{sec:emp-placebo,sec:sim-dependent}).

  \item \emph{An anchored relative test and a three-valued verdict.} Relative similarity tests
  compare two candidate models with one reference sample \citep{bounliphone2016test}. In ours one
  candidate is the real past, so the null states that $G$ is at least as close to the future as
  the past is. A single critic is trained to separate $\Rnew$ from an equally weighted pool of
  synthetic and past training windows. The statistic $\Danc$ is its mean on held-out past windows
  minus its mean on held-out synthetic windows (\cref{eq:anchored-stat}). At the boundary
  $\Law(G)=\Law(\Rpast)$ the two held-out samples are exchangeable given the critic, so for
  independent windows the permutation test is exact for every critic (\cref{prop:anchored}). The
  obvious alternative, the difference of two separately trained critic statistics, has no such
  guarantee: at that boundary its size is at least the nominal level, and for fixed critics and
  growing evaluation sets it tends to $1/2$ (\cref{prop:boundary-size}). A generator is \BAD{} if
  the anchored test rejects. It is \GOOD{} if the test does not reject and a power check shows
  that in simulated reference designs the same pipeline detects a synthetic volatility 20\% too
  high with probability at least $0.8$. It is \INCONCL{} otherwise
  (\cref{def:power-check,def:verdict,prop:verdict}). On real data, where windows are dependent and
  the past is itself non-stationary, the verdict combines Driscoll--Kraay p-values over several
  synthetic seeds; the exactness results above do not cover this rule, so the power check also
  estimates its size at the boundary, in a stationary design and in one that keeps the real
  volatility path of the past (\cref{sec:method-decision}). The check was fixed before the full
  real-data runs; the second design and the seed combination were added on the day of those runs,
  after development runs on six tickers (\cref{sec:emp-power}).

  \item \emph{Calibration against ground truth and likelihood-ratio oracles.} On independent
  simulated windows we measure the size and power of the \WC{} test and of five competitors: MMD
  \citep{gretton2012kernel}, a signature MMD \citep{chevyrev2022signature}, the classifier
  two-sample test \citep[C2ST;][]{lopezpaz2017revisiting}, a learned log-likelihood-ratio
  statistic \citep[\LRC{};][]{cheng2022classification} and the Kolmogorov--Smirnov (KS) test on
  pooled returns. The benchmark is the Neyman--Pearson oracle, which no two-sample test can beat
  (\cref{prop:oracle-bound}). The theory predicts a large gap. With a 50/50 split and equal sample
  sizes, the local noncentrality of any test based on the held-out mean difference of a fixed
critic is at most half the oracle's
  (\cref{prop:local-power}; proved for a fixed critic, sketched for a trained one). The comparison
  states where simpler tests win, and it includes alternatives that change only the dependence
  within a window (\cref{sec:sim}).

  \item \emph{Likelihood benchmarks on S\&P~500 data.} Five univariate generators (GBM, Merton,
  \GARCHt{}, Heston and the stationary bootstrap) are fitted to 100 stocks on 2018--2022 and
  evaluated on 2023-01-03 to 2025-06-16. Three of them have a tractable likelihood (GBM, Merton,
  \GARCHt{}). For these we compare the verdicts with out-of-sample predictive log-likelihood,
  probability-integral-transform tests and Amisano--Giacomini comparisons
  \citep{diebold1998evaluating,berkowitz2001testing,amisano2007comparing}. Heston has no
  closed-form likelihood \citep{aitsahalia2007maximum} and the bootstrap no parametric one; these
  are the cases that call for a likelihood-free test (\cref{sec:emp}). The critic ranks GBM,
  Merton and \GARCHt{} exactly as the predictive log-score does; GBM, Merton and Heston are \BAD{}
  and \GARCHt{} is \INCONCL{} (\cref{sec:emp-verdicts}).

  \item \emph{A link to Markowitz utility.} We show that critic distance and portfolio utility can
  disagree in both directions (\cref{sec:mk-theory}). The utility of the global minimum-variance
  (GMV) portfolio depends on the covariance matrix only up to scale, so a generator whose only
  error is a scalar factor loses nothing for GMV, yet a consistent critic test rejects it
  (\cref{ex:scale-blind}). Conversely, an error in a low-variance direction, which GMV weights
  heavily, can be nearly invisible in $\W$ (\cref{ex:spread}). Empirically, we build GMV and
  mean-variance portfolios from the output of five multivariate generators, evaluate them on the
  real evaluation period against the same portfolios estimated on the real calibration period,
  and relate the utility losses to multivariate critic and MMD statistics. The statistics do not
  predict the losses; the distance between the synthetic and the past average correlation does,
  with a negative sign, consistent with the fall of the real correlations, a sign that could
  reverse in a period of rising correlations (\cref{sec:markowitz}).
\end{enumerate}

\paragraph{Main findings.}
We call \emph{window tests} the tests that compare the laws of whole windows (\WC{}, \LRC{}, C2ST,
MMD, \SigMMD{}), as opposed to KS on the $nL$ pooled returns. With independent windows every
window test holds its level: in each family and at both sample sizes, the Wilson 95\% interval of
its rejection rate at nominal 5\% contains 0.05 (\cref{tab:e2-size}). KS, which treats the pooled
returns as independent, is not valid when returns within a window are dependent: returns within a
GARCH window share their conditional variance, and KS rejects a true null at 0.110 [0.074, 0.161]
with $n=500$ windows and at 0.080 [0.050, 0.126] with $n=2000$ (Wilson 95\% intervals in
brackets); in the other three families its intervals contain 0.05.
Against deviations of the one-dimensional marginal (volatility level, tails, jumps), KS is
nevertheless far more powerful than every window test. At $n=2000$ it rejects Gaussian against
Student-$t_4$ returns at matched variance, and a variance-matched GBM against a Merton
jump-diffusion, in every replication; the best window test, the \WC{} test, rejects at 0.67 and
0.20. Against volatility clustering at matched variance the window tests have power: with $n=2000$
and GARCH persistence 0.95 against 0.98, they reject at 0.79 (\WC{}), 0.93 (MMD), 0.86
(\SigMMD{}), 0.55 (C2ST) and 0.30 (\LRC{}) (\cref{tab:e2-power,fig:e2-power}). KS rejects even more
often there, but it is oversized and its power comes from the thinner tails that lower persistence
gives the marginal. When the marginal is held exactly fixed, KS has no power and the window tests
detect the deviation: the \WC{} test rejects at 0.46, 0.94 and 1.00 as the share of windows made
of independent returns rises from one half to one, below MMD and \SigMMD{} at one half (0.70)
and comparable to them at the larger shares, and above C2ST and the \LRC{}
(\cref{fig:e2dep-power}). No window test dominates. A kernel test is the most
powerful window test, or tied for it, at every alternative of volatility level and of
clustering, while the \WC{} test is the most powerful against the lightest-tailed alternatives to
the Student-$t_4$ law ($t_{30}$ and Gaussian). At the smallest deviation of each family with
$n=500$, every window test stays far below the oracle, which rejects in at least 0.98 of the
replications at every alternative of the four families; against a quarter
or a half of GBM windows in a Merton sample no window test rejects more often than 0.08. The
learned scores are tests, not likelihood estimators: their Spearman correlation with the exact
log-likelihood ratio lies between 0.024 and 0.556 (\cref{tab:e2-llr}). At the relative boundary
$\Law(G)=\Law(\Rpast)$, with $n=500$, the anchored test rejects at 0.065 [0.038, 0.108] and 0.055
[0.031, 0.096] in two designs, and the two-critic test at 0.040 and 0.065, so the excess that
\cref{prop:boundary-size} allows is not visible at this sample size. A generator that resamples
the past makes the anchored test reject at 0.965 with a random split, against 0.010 when the split
respects the source window (\cref{tab:e2-relative-boundary}). With consecutive windows of one GARCH
path, the random-split permutation test is oversized (0.145 for the \WC{} test, 0.290 for MMD),
and the Driscoll--Kraay p-value after a time-block split and the wild bootstrap bring the rates to
0.070 and 0.075 (\cref{tab:e2-dependent-size}).

In a simulated factor panel with the dimensions of our data, the time-block Driscoll--Kraay
protocol rejects a true null at 0.085 [0.054, 0.132] for the \WC{} test (0.050 against a fixed-$b$
reference) and the wild-bootstrap MMD test at 0.045. The ticker split rejects at 0.990 even with
Driscoll--Kraay standard errors, the permutation p-value after a time-block split at 0.410 and KS
at 0.970 (\cref{tab:e1-size-sim}). When the synthetic tickers are drawn independently, as the
univariate generators draw them, the learned tests over-reject (0.270 for the absolute \WC{} test,
0.155 for the anchored test). The statistic itself is no more biased than when the synthetic
panel has a factor, but it is negatively correlated with its Driscoll--Kraay standard error, and
the rejections concentrate in the replications with a small standard error
(\cref{tab:e1-sim-se}); \cref{cor:block-bias} gives a model-based explanation, which we have not
tested. Tested against each other with the same protocol, the real \TRAIN{} and \TEST{} windows are not rejected as
different ($\pmed=0.452$ for the \WC{} test; \cref{tab:e1-nonstat}), but with few evaluation dates
(\cref{rem:small-T}) this non-rejection says little about whether the two periods are similar. The
relative null is motivated by the volatility level that the generators inherit, not by a
rejection of the past. The power check shows the cost of non-stationarity. In a stationary
reference design the verdict rule detects a 20\% volatility error in all 100 replications and
never rejects at the boundary; with the real volatility path of 2018--2022 its power is 0.45 [0.36, 0.55] and its boundary size 0.08 [0.04,
0.15] (\cref{tab:e3-power}), so on this panel and split no generator can be \GOOD. The absolute
\WC{} test rejects all five generators (\cref{tab:e3-pooled-tests}). The anchored test rejects GBM
($\pmed<0.001$), Merton (0.005) and Heston (0.017), which, as calibrated on 2018--2022, are \BAD{},
Heston borderline. It does not reject \GARCHt{} (0.432), which is distinguishable from the future
but not detectably worse than history, and \INCONCL{}; the stationary bootstrap, which copies the
past, receives no verdict (\cref{tab:e3-verdicts}). The witness shows that the critic mainly
detects a volatility level: every generator is too volatile for 2023--2025
(\cref{sec:emp-witness}). Each verdict therefore concerns a model together with its calibration
on a period that contains the 2020 and 2022 turbulence, not the model class alone. The critic ranks
\GARCHt{}, Merton and GBM exactly as the predictive log-score does (Kendall $\tau=1$), whereas its
per-ticker rejections agree little with Berkowitz tests (Cohen $\kappa=0.11$), which check one-day
calibration rather than simulated windows (\cref{tab:e3-agreement}). In the sensitivity analysis
this ranking holds in all 14 configurations, but the \BAD{} verdicts of Merton and Heston do not,
and with an alternative split every parametric generator is \INCONCL{} (\cref{sec:robust}).
Fitting is itself informative. Simulated with the quadratic-exponential scheme of \citet{andersen_2008_simple},
Heston needs a bounded volatility of variance to avoid days with zero variance, and it then falls
well short of the kurtosis of the data (\cref{tab:e3-heston-moments}).

In the Markowitz study, the average pairwise correlation of real returns fell from 0.419 in
\TRAIN{} to 0.268 in \TEST{} (\cref{tab:e4-correlation}). Four of the five generators give
minimum-variance portfolios whose \TEST{} variance is within 6\% of that of the portfolio trained
on real data (replica-wise variance ratios). With pooled moments, CCC-GARCH-$t$ and MV Student-$t$
beat it (variance ratios 0.869 and 0.896); their average correlations are the furthest below the
past one. Factor GARCH-$t$ loses most (1.466 with pooled moments, 1.436 replica-wise) although its
average correlation is close to the past one (\cref{tab:e4-utility,tab:e4-utility-agg}). In this
evaluation period test results do not predict these losses. MV Gaussian is rejected in all 30
universes after size adjustment (first seed), yet its minimum-variance portfolios have \TEST{}
variance ratios close to one (1.038 replica-wise, 1.004 with pooled moments). Within a generator
the critic statistic is barely reproducible across synthetic seeds (Spearman correlation between
the two seeds from $-0.03$ to 0.30; \cref{tab:manual-e4-reliability}), so its near-zero
correlation with the minimum-variance loss (0.01) is weak evidence. The losses are point
estimates over one \TEST{} period, without inference over dates. The one consistent correlate is the distance between the synthetic and the past average correlation, with
a negative sign ($-0.36$ within generators; \cref{tab:e4-tests,tab:e4-link}), which reflects the
fall in correlations and need not hold in a period in which they rise.

\paragraph{What we do not claim.}
The \WC{} test is one well-calibrated and interpretable tool within a validated protocol, not a
uniformly most powerful test. KS is more powerful against marginal deviations, and a kernel test
is at least as powerful as the \WC{} test at every alternative of volatility level and of
clustering in our simulations. The power of the critic depends on its optimiser settings and its
architecture (\cref{sec:robust}). In a post-hoc ablation, whose architecture was chosen after the
real-data results of the pre-specified perceptron were known and after a pilot on four of the
simulated designs, a time-pooled convolutional critic keeps the level and comes close to KS
against tails and jumps: with $n=500$ it rejects $t_{10}$ against $t_4$ windows at 0.98, where the
perceptron rejects at 0.10 and KS at 1.00 (\cref{tab:e5-arch}). The comparison with KS is thus a
statement about the pre-specified network, whose verdicts are the ones we report. We do not
compare with power-optimised deep-kernel tests \citep{liu2020learning}, and on real data, where
the critic behaves almost like a realised-volatility statistic (\cref{sec:emp-witness}), a fixed
log realised-volatility statistic on the same seed-0 windows ranks the three likelihood models
in the same order and reaches between 0.69 and 0.94 of the critic's $\zhac$
(\cref{tab:e3-rv-baseline}); we did not run it inside the verdict rule. The protocol contributes exact or asymptotically valid inference on
dependent panels, the relative null, the link with likelihood and an extension to multivariate
windows, where pooled KS sees only the marginal; the size of its dependent-data version is
measured, not proved, and it is not nominal in every design. A \GOOD{} verdict means that a
procedure with demonstrated power against a deviation of pre-specified size did not detect one. It
never means that the generator equals the truth, and its error control rests on a premise that
cannot be verified (\cref{prop:verdict}). The relative null does not penalise a generator error
that happens to point towards the change of regime, and the power check measures power only
against a volatility 20\% too high, an error that on our data points away from the calmer
evaluation period. A rejection is conditional on the fitted parameters, so
it does not separate misspecification from estimation error or from parameter drift between the
two periods; the relative test is designed to absorb the drift. Under non-stationarity a verdict
concerns one pair of periods. Univariate verdicts say nothing about the cross-sectional
dependence that drives portfolio risk (\cref{ex:corr-blind}), and whether synthetic data are good
for a task has to be checked on that task. Finally, the panel consists of the S\&P~500
constituents at the end of the sample. This is a selection on outcomes of the evaluation period
that goes beyond survivorship bias \citep{brown1992survivorship,shumway1997delisting}, and it may
make the real evaluation panel calmer than the one a forecaster faced at the end of 2022.

\paragraph{Outline.}
\Cref{sec:lit} reviews related work and \cref{sec:method} defines the procedure. \Cref{sec:theory}
states the theory: exact level, consistency, the likelihood link, the relative test, inference
under dependence and what the verdict controls. \Cref{sec:sim} reports the simulation evidence,
\cref{sec:data} describes the data and the generators, \cref{sec:emp} gives the S\&P~500 results
and \cref{sec:markowitz} studies Markowitz utility. \Cref{sec:robust} checks robustness and
\cref{sec:concl} concludes. Proofs are in \cref{app:proofs}, implementation details in
\cref{app:impl} and further results in \cref{app:sim,app:data,app:emp,app:mk,app:robust}.

\section{Related literature}\label{sec:lit}
% =============================================================================================
% Source: docs/literature.md (entries, "Closest competitors", "Positioning for a CREST
% audience", "Citation map"). Only keys of references.bib (Sections 9 and 10 of the .bib hold
% the keys proposed in OUTLINE.md Section 13 and the six keys added in the referee revision:
% im_2018_quantitatively, gulrajani_2019_towards, sutherland_2017_generative,
% jitkrittum_2016_interpretable, meehan_2020_three, lucic_2018_gans; each checked against the
% publisher or arXiv record at integration, 2026-11-16). Positioning only: no
% results, no numbers. Never "first". Subsection labels: OUTLINE.md Section 12, plus
% sec:lit-summary (added here).
% =============================================================================================

This section positions the paper in five literatures; \cref{sec:lit-summary} summarises what is
and is not ours.

\subsection{Generative models for financial time series and their evaluation}
\label{sec:lit-synthetic}

Synthetic return data is used for back-testing, data sharing and hedging
\citep{assefa2020generating,kondratyev2019market,buehler2019deep,bailey2017probability}.
Generators include GANs \citep{wiese2020quant,takahashi2019modeling}, restricted Boltzmann
machines \citep{kondratyev2019market}, signature-based variational autoencoders
\citep{buehler2020data}, generators trained with signature losses
\citep{ni2021sig,liao2024sig,issa2023nonadversarial}, generic time-series GANs
\citep{esteban2017real,yoon2019time} and generators with use-specific losses such as portfolio
tail risk \citep{cont2026tail}. Parametric models remain natural benchmarks (GARCH(1,1), for
instance, is hard to beat out of sample; \citealp{hansen2005forecast}); together with a stationary
bootstrap they are the generators of this paper (\cref{sec:data-generators}).

Evaluation practice is heterogeneous \citep{stenger2024evaluation}. Most papers compare stylized
facts \citep{cont2001empirical} visually or through summary statistics
\citep{wiese2020quant,takahashi2019modeling,vyetrenko2020get}. Others report fidelity scores
\citep{alaa2022faithful} or discriminator scores: the held-out error of a classifier trained after
the fact \citep{yoon2019time}, a real-versus-simulated classification competition
\citep{francopedroso2018ets}, or a discriminator output used to calibrate a market simulator
\citep{storchan2021learning}. In machine learning, critics for several GAN objectives, the
Wasserstein one included, have been trained after the fact on training samples and the resulting
divergences evaluated on held-out samples as metrics for image generators
\citep{im_2018_quantitatively}, and \citet{gulrajani_2019_towards} define such neural-network
divergences on held-out data so that a generator that memorises its training set does not score
well. Such scores have no null distribution, and a low discriminative score
cannot be told apart from a weak classifier unless the power of the procedure is known; hence the
pre-specified power check in our decision rule (\cref{sec:method-decision}).

Formal testing of market generators is not new. \Citet{buehler2020data} argue for two-sample tests
as evaluation metrics and apply marginal Kolmogorov--Smirnov tests and the signature-MMD test of
\citet{chevyrev2022signature}, whose distribution-free threshold is valid for every null law and
hence conservative; they report no size or power analysis. \Citet{bai2026calibrating} come closer
to this paper: to calibrate non-identifiable simulators, among them a limit-order-book market
simulator, they compare real and simulated outputs through a feature map, possibly the output of a
trained WGAN critic, with Bonferroni-corrected Kolmogorov--Smirnov or mean-difference statistics
and asymptotic guarantees for independent outputs. A WGAN critic has thus already been used, with
frequentist guarantees, to validate a financial simulator against real data. Our goal differs (a
verdict on one fitted generator, not a set of eligible parameters), and so does the statistic:
their critic is a fixed pre-trained feature map (the most conservative of their features), ours is
trained on one half of the two samples and its mean difference on the other half is the statistic.
Dependent windows, a relative null and the likelihood-ratio oracle are outside their scope.

\subsection{Two-sample tests with learned test functions}\label{sec:lit-tests}

An integral probability metric (IPM) is $\sup_{f\in\mathcal F}\abs{Pf-Qf}$ over a class
$\mathcal F$ of test functions \citep{muller1997integral}; over the 1-Lipschitz functions it is
$\W$, by Kantorovich--Rubinstein duality \citep{villani2009optimal}. The critic of a Wasserstein GAN maximises the empirical
objective over a class of networks \citep{arjovsky2017wasserstein}, with a gradient penalty in
place of the Lipschitz constraint in WGAN-GP \citep{gulrajani2017improved}. Our critic, trained
alone, is a neural-network test function chosen to maximise a mean difference under a Lipschitz
penalty. An IPM over a restricted class is estimable at the parametric rate
\citep{sriperumbudur2012empirical}, whereas the empirical measure converges in $\W$ at rate
$n^{-1/d}$ for $d\ge3$ \citep{fournier2015rate}. A finite critic measures a weaker ``neural net distance''
\citep{arora2017generalization}, and with a soft penalty a good critic need not estimate $\W$ well
\citep{mallasto2019well}. Our statistic is therefore an IPM over trained critics, whose
conditional mean is bounded by the critic's Lipschitz constant times $\W$ between the standardised
window laws (\cref{prop:unbiased,rem:gp-soft}), not an estimate of $\W$. A trained critic is
also a random function of its seed and hyper-parameters, and large-scale comparisons of GANs
depend on random restarts and on the hyper-parameter search budget \citep{lucic_2018_gans}. We
therefore combine p-values over seeds by a rule that preserves the validity of the individual
p-values under arbitrary dependence between them (\cref{rem:combine-seeds}) and report how the
verdicts change with the critic's hyper-parameters (\cref{sec:robust}).

The MMD test compares kernel mean embeddings \citep{gretton2012kernel}.
\Citet{sutherland_2017_generative} choose the kernel on a training split by maximising an estimate
of a power proxy, test on the held-out split, and use the optimised MMD and its witness function
to criticise generative models; \citet{liu2020learning} extend this to deep kernels (MMD-D), and
\citet{jitkrittum_2016_interpretable} choose test locations by maximising a lower bound on power,
so that the selected features show where the two laws differ. These tests, like the evaluation
metrics of \cref{sec:lit-synthetic}, treat the samples as independent. The classifier
two-sample test (C2ST) of \citet{lopezpaz2017revisiting} tests whether held-out accuracy differs
from one half. Accuracy discards magnitude: in the Gaussian mean-shift problem accuracy tests are
only $1/\sqrt{\pi}$ as efficient as Hotelling-type tests \citep{kim2021classification}. The
held-out difference of mean logits \citep{cheng2022classification} keeps it; since the Bayes logit
is the log density ratio
\citep{goodfellow2014generative,nguyen2010estimating,cranmer2015approximating}, it is a learned
log-likelihood-ratio statistic, our \LRC{} competitor (\cref{sec:method-competitors}).

The \WC{} test instantiates a known template: learn a real-valued witness on one split and test the
equality of its means on the other, with a permutation null \citep{kubler2022witness}; see also
the location tests on learned representations of \citet{kirchler2020two}, the logit instance of
\citet{cheng2022classification} and the witness-based model criticism of
\citet{lloyd2015statistical}. The split-then-permute structure of a Wasserstein test is not ours
either: \citet{wang2022two} (linear version: \citealp{wang2021two}) train a kernel projector on one
split to maximise the Wasserstein distance between projected samples and permute on the held-out
split with the projector fixed. \Citet{hu2026neural} use a deep, approximately 1-Lipschitz
witness, as we do, but with learned projections, max-type aggregation and asymptotically pivotal
calibration for \iid{} data; see also \citet{ramdas2017wasserstein}. Our test is the template
with a WGAN-GP critic: the statistic is the held-out mean difference of one network (no transport
problem is solved on the evaluation split), calibrated by permutation for independent windows and
by a Driscoll--Kraay variance for dependent ones. We claim this instance, its validation and the
protocol around it, not the template. We do not show that it is more powerful than MMD-D, which is
not among our baselines (\cref{sec:concl-limitations}), and it does not beat the baselines we run
against every deviation: against changes of the one-dimensional marginal of returns, the pooled
Kolmogorov--Smirnov test is far more powerful than every window test in our simulations, although
it is invalid when returns within a window are dependent and blind to deviations that leave the
marginal unchanged (\cref{sec:sim-power}).

\Citet{chevyrev2022signature} show that (robust) signature moments characterise laws of
stochastic processes and build an MMD test on them
\citep[see also][]{lyons1998differential,kiraly2019kernels}. Signature discrepancies also serve as
training losses for market generators \citep{ni2021sig,liao2024sig,issa2023nonadversarial}, one
reason to evaluate with a test that was not the training objective. Our \SigMMD{} baseline is a
truncated (depth-4) signature-feature MMD in their spirit, without tensor normalisation; it does
not inherit their characteristic property.

A relative test compares two candidate models with one reference sample.
\Citet{bounliphone2016test} use the difference of two MMD estimates and their joint asymptotic
normality, which requires independent \iid{} samples and two candidates that both differ from the
reference, the MMD U-statistic being degenerate at equality; our MMD relative test follows them
and inherits this condition. See also \citet{jitkrittum2018informative} and
\citet{gerber2024likelihood} on deciding which of two simulators produced the data. Closest in
aim, \citet{zhou_2025_anchor} ask which of two distributions is closer to an anchor distribution:
they maximise the difference between the two anchored discrepancies over deep kernels, which also
selects the direction of the hypothesis, and then calibrate the selected test, for independent
samples. Our relative test differs in four respects. One of the two candidates is the real past,
so the question is whether a generator does better than historical simulation; the windows are
serially and cross-sectionally dependent; the discrepancy is the held-out mean difference of one
critic, whose witness can be inspected; and the test enters a decision rule with a pre-specified
power check.
% zhou_2025_anchor: description checked against the arXiv:2510.10477 abstract (anchor-based
% maximum discrepancy over a deep-kernel space; Phase I infers the hypothesis, Phase II tests it).
In our relative null, then, $\Rpast$ is one of the candidates: is $G$ as good a proxy for the future
$\Rnew$ as history is? With learned critics, the difference of two separately trained
discrepancies compares IPMs over two data-dependent classes, and at the boundary $\QG=\Ppast$
its size is at least the nominal level and can exceed it (\cref{prop:boundary-size}). The
anchored test trains one critic to separate $\Rnew$ from a balanced pool of synthetic and past
windows and compares its means on held-out synthetic and past windows; its permutation p-value is
exact at the boundary for every critic (\cref{prop:anchored}). The argument is the elementary
fixed-witness one; its value is exactness at the boundary, where the other two versions are not
exact. Exactness needs the synthetic and past windows to be exchangeable at the boundary
(\cref{ass:anc}). Serially dependent windows are not, and the Driscoll--Kraay version of the test
used on such panels has no such guarantee: its size is measured
(\cref{sec:sim-dependent,sec:emp-placebo}). A generator that copies its training data violates the
assumption as well. Data copying is a known failure of generator evaluation:
\citet{meehan_2020_three} test for it by comparing the distances of generated and of held-out real
samples to the training sample. Our stationary bootstrap copies blocks of $\Rpast$ by
construction, so we do not test for copying; we withhold its verdict (\INVALID{};
\cref{sec:method-relative,tab:e2-relative-boundary}).

\subsection{Specification testing and density forecast evaluation in econometrics}
\label{sec:lit-econometrics}

For an econometrician the natural frame is the specification testing of simulation-based models.
Models without a closed-form likelihood are estimated by simulated moments
\citep{mcfadden1989method,duffie1993simulated}, indirect inference
\citep{gourieroux1993indirect,smith1993estimating}, the efficient method of moments
\citep{gallant1996which} or simulated maximum likelihood \citep{fermanian2004nonparametric}; see
\citet{gourieroux1997simulation}, and \citet{bernton2019approximate} and \citet{cherief2022finite}
for Wasserstein ABC and minimum-MMD estimation under misspecification. Discriminators are accepted econometric objects: in the
adversarial estimator of \citet{kaji2023adversarial} a logistic discriminator gives optimally
weighted simulated moments and the oracle discriminator maximum likelihood \citep[see
also][]{kaji2023metropolis}. Our paper is a testing counterpart: we test a calibrated generator
instead of estimating one, use an IPM instead of cross-entropy, and work with serially dependent
windows, which their theory excludes. Their bridge from discriminators to maximum likelihood is the
estimation counterpart of our comparison of the \LRC{} statistic with the Neyman--Pearson oracle
(\cref{prop:bayes-logit,prop:oracle-bound}).

Specification tests of financial models compare the marginal density of a diffusion with a kernel
estimate \citep{aitsahalia1996testing} or test transition densities through probability integral
transforms \citep{hong2005nonparametric}, which requires the transition density. Closest to us,
\citet{corradi2005bootstrap} and \citet{bhardwaj2008simulation} compare the empirical (marginal,
respectively conditional) distribution of the data with that of data simulated from the fitted
model, covering stochastic volatility without a closed-form likelihood, with block-bootstrap
critical values that account for serial dependence and parameter estimation error. Our test
differs in three respects. (a) \emph{Null}: they test the model class at the pseudo-true
parameter, so estimation error must be corrected; we test the fitted generator (parameters
estimated on \TRAIN{} and held fixed, null $\QG=\Pnew$ on \TEST{}), so no correction is needed,
but a rejection does not separate misspecification from estimation error or parameter drift, which
the relative test is designed to absorb. (b) \emph{Statistic}: a fixed Kolmogorov-type functional
versus a learned function. (c) \emph{Calibration}: a block bootstrap versus permutation or a
Driscoll--Kraay variance, conditional on the fitted parameters and the critic.

In this language the critic is a test function chosen from the data. Consistent specification
tests use a whole class of test functions, such as exponential weights
\citep{bierens1990consistent} or neural hidden units indexed by a nuisance parameter present only
under the alternative \citep{stinchcombe1998consistent,lee1993testing}, with a supremum or integral
over the class and asymptotic calibration. The \WC{} test is a learned, sample-split relative: it
selects one function on a training split and calibrates on the other. The analogy concerns the test
functions only (their null is a conditional moment restriction), and the \WC{} test inherits the
power of the learned function, not the consistency of Bierens-type tests; it is consistent only
under conditions on training (\cref{prop:consistency,rem:c-realism}).

When a likelihood exists, density forecast evaluation supplies trusted benchmarks: probability
integral transforms are \iid{} uniform under a correct predictive density
\citep{diebold1998evaluating}, \citet{berkowitz2001testing} tests their normal transforms against a
Gaussian AR(1), the log score is strictly proper \citep{gneiting2007strictly}, and density forecasts
are compared by out-of-sample log-likelihood ratios with a HAC variance
\citep{amisano2007comparing}, in the spirit of \citet{vuong1989likelihood}. We use them for GBM,
Merton and GARCH-$t$ (\cref{sec:emp-likelihood}); Heston has no closed-form likelihood
\citep{aitsahalia2007maximum} and the stationary bootstrap no parametric likelihood. Likelihoods rank models by
Kullback--Leibler divergence and critics by an IPM, so disagreement is informative
(\cref{rem:kl-vs-w1}). The relative test is a loss-differential comparison
\citep{diebold1995comparing} that, like \citet{giacomini2006tests}, evaluates a forecasting method
(``use the real past''); its closest econometric relative is the simulation-based comparison of
misspecified diffusion models by \citet{corradi2011predictive}.

\subsection{Inference under dependence}\label{sec:lit-dependence}

A permutation test is exact when the pooled observations are exchangeable under the null
\citep{lehmann2022testing}, also with random permutations and the $+1$ correction
\citep{hemerik2018exact,phipson2010permutation}; with a critic fixed by the training half this
gives exact level for independent windows (\cref{prop:exact-level}). Studentised permutation tests
are asymptotically valid for the weaker null of equal means, for independent samples
\citep{janssen1997studentized,chung2013exact}; under time-series dependence the closest results
concern autocorrelation tests \citep{romano2022permutation}. Real return windows are not
exchangeable: consecutive windows share volatility regimes, and windows of different stocks on the
same dates share market shocks. We studentise the held-out mean difference with a Bartlett-kernel
HAC variance \citep{newey1987simple} and the bandwidth of \citet{andrews1991heteroskedasticity},
after summing centred critic values by date \citep{driscoll1998consistent}. The normal
approximation rests on a central limit theorem for mixing sequences \citep{herrndorf1984functional},
and GARCH and a class of stochastic-volatility models are $\beta$-mixing
\citep{carrasco2002mixing}. With few time periods HAC tests over-reject; fixed-$b$ critical values
\citep{kiefer_2005_new} are one correction, reported for comparison (\cref{rem:small-T}). The MMD
statistic is a degenerate U-statistic under the null rather than a mean, and
\citet{chwialkowski2014wild} calibrate it under dependence by a wild bootstrap with dependent
multipliers, which we use for MMD and \SigMMD{} (\cref{rem:wild}); block bootstraps
\citep{kunsch1989jackknife,politis1994stationary,lahiri2003resampling} are the alternative.

Two features are less standard. First, the critic is trained on data that depend on the evaluation
data. A ticker split puts both halves on the same dates and biases the statistic
(\cref{prop:ticker-bias}); a split by contiguous time blocks makes training and evaluation dates
disjoint, which removes the common-date bias but does not make the two halves independent, and
a persistent factor can still bias the statistic across adjacent blocks (\cref{cor:block-bias}). Under
mixing, \cref{lem:decoupling} bounds the remaining dependence across block boundaries only when a
gap separates the two halves; the implemented split has none, and under persistent volatility
adjacent blocks remain strongly dependent, so the effect of this dependence on size is measured,
not bounded (\cref{sec:sim-dependent,sec:emp-placebo}). Second, per-ticker tests are
dependent through the market factor: Benjamini--Hochberg control relies on positive regression
dependence, and arbitrary dependence needs the Benjamini--Yekutieli correction
\citep{benjamini1995controlling,benjamini2001control}.

\subsection{Portfolio choice and use-based evaluation}\label{sec:lit-portfolio}

Mean-variance portfolios \citep{markowitz1952portfolio} are very sensitive to estimation error
\citep{michaud1989markowitz,chopra1993effect,kan2007optimal}; shrinkage \citep{ledoit2004well},
long-only constraints \citep{jagannathan2003risk} and the equally weighted portfolio
\citep{demiguel2009optimal} are standard responses. A generator fitted on \TRAIN{} is a randomised
function of \TRAIN, so its output carries no more information about the future than \TRAIN{}
itself. \Citet{cetingoz2026synthetic} make this precise for Markowitz and minimum-variance
portfolios: the initial sample size limits what synthetic data can add, generating far more data
than was available can bias the estimated statistic, and generic generators fit high-variance
principal components that matter little for minimum-variance portfolios; they argue for
use-targeted evaluation. Synthetic-trained portfolios can thus beat a specific real-data rule, such
as the sample-covariance GMV portfolio, only when a parsimonious generator acts as a structured,
shrinkage-like estimator (\cref{prop:no-free-lunch}).

Judging synthetic data by a downstream decision is the train-on-synthetic, test-on-real (TSTR)
principle \citep{esteban2017real,yoon2019time}; in econometrics, \citet{athey2024using} judge
generated data by the ranking of estimators it produces, and evaluation criteria for generative
models are known to disagree \citep{theis2016note}. Our Markowitz study uses TSTR as the utility
measure but asks whether a distributional test verdict predicts the TSTR utility gap across
universes and generators, which \citet{cetingoz2026synthetic} do not address. Theory suggests at
most a weak link (\cref{sec:mk-theory}): GMV utility depends on second moments only, so critic
distance and utility can disagree in both directions, and under non-stationarity a generator that
reproduces the past better is not necessarily more useful (\cref{tab:e4-correlation}).

\subsection{Summary of the positioning}\label{sec:lit-summary}

None of the following is ours: a critic or neural-network divergence trained after the fact to
evaluate a generator \citep{im_2018_quantitatively,gulrajani_2019_towards}; a kernel learned on one
split to criticise a generative model \citep{sutherland_2017_generative}; a witness learned on one
split and tested on the other
\citep{kubler2022witness,kirchler2020two,cheng2022classification}; a split-then-permute
Wasserstein test \citep{wang2022two}; a WGAN critic used to validate a financial simulator
\citep{bai2026calibrating}; a formal two-sample test of a market generator
\citep{buehler2020data}; a specification test comparing real data with data simulated from the
fitted model, with dependence handled by a block bootstrap \citep{corradi2005bootstrap}. What the paper
adds is their combination and its validation, as listed in \cref{sec:intro}: an inference protocol
for dependent panels with measured size, an anchored relative null benchmarked on the real past,
oracle-based power comparisons that show where simpler tests win, a likelihood link and a
Markowitz study. The \WC{} test is one well-calibrated, interpretable tool within this protocol,
not a uniformly more powerful test.

\section{Methodology}\label{sec:method}
% =============================================================================================
% Owner: methodology writer. Describes what src/synthval does (twosample/critic.py, nets.py,
% c2st.py, mmd.py, signature.py, ks.py, relative.py; pipeline.py; experiments/e3_sp500.py for
% the verdict). Formal statements live in Section 4 and are only referenced here.
% Equations defined here: eq:wgan-gp, eq:critic-stat, eq:mmd-u, eq:lr-oracle, eq:perm-pvalue,
% eq:dk-var, eq:hac-z, eq:anchored-obj, eq:anchored-stat, eq:two-critic-stat, eq:pmed,
% eq:verdict. Tables: tab:method-tests, tab:method-defaults. Figure: fig:alg-verdict.
% Numbers (referee revision): E1 sizes 0.085/0.080 (joint), 0.270/0.155 (indep) from
% tables/e1_size_sim.tex; D1/D2 sizes and power from results/e3/e3_power_summary.csv and
% e3_realvol_summary.csv (single-seed boundary intervals are Clopper--Pearson, as in
% tables/e3_power.tex); learning-rate discordant pairs 19:0 from results/e5/e5_simpower_raw.csv
% (family a, 1000 steps, B=1000, alpha=0.05); protocol dates from docs/DESIGN.md change log and
% the written/quick fields of results/{e3,e3_quick}/meta_*.json.
% =============================================================================================

This section defines the tests exactly as they are implemented in the code that produces every
number in the paper. All of them share one template: a test function is learned on a
\emph{training half} of the two samples and evaluated, frozen, on the held-out
\emph{evaluation half}. What changes from one setting to another is how the evaluation
statistic is calibrated. For independent windows we use permutation p-values, which have exact
finite-sample level. For the serially and cross-sectionally dependent windows of real data we
use a time-block split with Driscoll--Kraay standard errors. This studentisation is
asymptotically valid for a fixed critic whose evaluation data are independent of its training
data, a condition that the implemented split meets only approximately, so we measure the size of
the whole pipeline in a placebo calibrated to the data instead of assuming it. On
top of the absolute test we build an anchored relative test and a three-valued verdict with a
pre-specified power check. Formal properties are in \cref{sec:theory}; further implementation
details are in \cref{app:impl}. \Cref{tab:method-tests,tab:method-defaults} summarise the tests
and their default settings.

% ---------------------------------------------------------------------------------------------
\subsection{Windows, samples and hypotheses}\label{sec:method-setup}

\paragraph{Returns and scaling.}
The data are daily log-returns $r_{it}=\log S_{it}-\log S_{i,t-1}$ of tickers $i=1,\dots,N$.
\TRAIN{} (returns from 2018-01-03 to 2022-12-30, $\Ttr=1258$ trading days; the price of
2018-01-02 is the base of the first return) is the calibration period and \TEST{} (2023-01-03 to
2025-06-16, $\Tte=615$ days) the evaluation period. For each ticker let
$s_i$ be the standard deviation (divisor $\Ttr-1$) of its \TRAIN{} returns. Real \emph{and}
synthetic returns of ticker $i$ are divided by the same constant $s_i$ before anything else.
Given \TRAIN{}, this is a fixed bijection applied to both samples, so it leaves every per-ticker
null hypothesis unchanged; its only role is to put tickers on a common scale so that their
windows can be pooled (\cref{rem:pooling}).

\paragraph{Windows.}
A window is a block of $L$ consecutive normalised returns. Windows do not overlap (step $L$),
and an incomplete tail is dropped. The default is $L=20$, about one trading month; the
per-ticker tests of E3 and the multivariate tests of E4 use $L=5$. A univariate window lies in
$\R^{L}$. A multivariate window of $N$ assets lies in $\R^{N\times L}$ (E4: $N=30$). Each window
carries the \emph{date rank} $\tau$ of its first day on the common trading calendar: \TRAIN{}
days have ranks $0,\dots,\Ttr-1$ and \TEST{} days $\Ttr,\dots,\Ttr+\Tte-1$. Synthetic windows
simulated over the \TEST{} horizon receive the \TEST{} ranks, and windows of different tickers
that start on the same day share their rank. At $L=20$ a ticker yields 62 \TRAIN{} and 30
\TEST{} windows, so the 100-ticker panel of E1, E3 and E5 gives 6200 and 3000 windows.

\paragraph{Samples.}
The roles of the samples never change. $X_1,\dots,X_n$, with law $P$, are the real windows of
the comparison sample; on real data $X=\Rnew$, the \TEST{} windows (law $\Pnew$ in relative
statements). $Y_1,\dots,Y_m$, with law $Q$, are synthetic windows from a generator $G$ (law $\QG$
in relative statements). In E3, $G$ is fitted per ticker on \TRAIN{} and simulates one
\TEST{}-length path per ticker, started from the fitted end-of-\TRAIN{} state when the model has
one, so each ticker contributes as many synthetic as real windows. This matters: the pooled null
follows from the per-ticker nulls only when the per-ticker proportions are the same in both
samples (\cref{rem:pooling}). $Z_1,\dots,Z_k$, with law $\Ppast$, are the real past windows
$\Rpast$ (\TRAIN{}); only the relative tests use them.
In the simulations of E2, $X\sim P$ and $Y\sim Q$ have known laws and every window is simulated
from its own path after a burn-in, so the windows are independent and identically distributed
(\iid{}). On real data, and on windows cut from long simulated paths, they are not. Volatility
clustering makes the windows of a ticker serially dependent, and common market shocks make
windows of different tickers on the same dates dependent.

\paragraph{Hypotheses.}
The absolute null is $H_0: P=Q$ (\cref{def:abs-null}). It concerns the law of \emph{one}
window. Dependence \emph{between} windows is not part of $H_0$; it is a nuisance that the
inference protocol of \cref{sec:method-inference} must absorb. A window law fixes all the
dependence of the return process up to lag $L-1$ (\cref{prop:window-ident}). A window test can
therefore detect volatility clustering within a month, but not features that act only at longer
horizons.

Markets are not stationary. \TRAIN{} contains the 2020 crash and the 2022 drawdown. The average
correlation of the E4 universes fell from \TRAIN{} to \TEST{} (\cref{tab:e4-correlation}), and
so did the volatility level of the panel: the panel median of the daily standard deviation is
1.62\% in \TEST{}, against 2.11\% for the GBM volatility fitted on \TRAIN{}
(\cref{tab:e3-stylized,tab:e3-fit-diagnostics}). E1 compares the \TRAIN{} and \TEST{} windows
directly; the window tests with valid p-values do not reject, but with about 15 \TEST{} evaluation
dates they have little power for this comparison (\cref{tab:e1-nonstat,sec:emp-placebo}). A generator calibrated on \TRAIN{} inherits features of that period,
such as its volatility level, and cannot be expected to reproduce the \TEST{} law exactly, so
absolute rejections are to be expected on real data. The relative tests of
\cref{sec:method-relative} ask a weaker and more useful question: is $G$ detectably
further from $\Rnew$ than the real past $\Rpast$ is?

% ---------------------------------------------------------------------------------------------
\subsection{The Wasserstein-critic statistic}\label{sec:method-critic}

\paragraph{Split and standardisation.}
Each sample is split into a training and an evaluation half, with index sets $\Itr,\Iev$
(real) and $\Jtr,\Jev$ (synthetic) and sizes $n_t,n_e$ and $m_t,m_e$; the default is 50/50.
For independent windows the halves are drawn at random within each sample (half of the windows,
rounded, for training). For dependent windows we use the time-block split of
\cref{sec:method-inference}. The split never depends on the values of the windows. Inputs are
standardised by $A(x)=\Lambda^{-1}(x-m_A)$, where $m_A$ and the positive diagonal $\Lambda$ hold
the mean and standard deviation of the pooled training half $\{X_i\}_{i\in\Itr}\cup\{Y_j\}_{j\in\Jtr}$.
They are computed per coordinate for univariate windows and per asset channel (pooled over the
$L$ positions) for multivariate windows. The same $A$ is applied to every window, training or
evaluation, real or synthetic.

\paragraph{Training.}
The critic is $\fhat=\ghat\circ A$, where $\ghat$ is a neural network. It is trained \emph{alone},
with no generator, with the gradient-penalised Wasserstein GAN objective of
\citet{gulrajani2017improved} (see also \citealp{arjovsky2017wasserstein}):
\begin{equation}\label{eq:wgan-gp}
  \ghat\approx\argmax_{g}\;\Big\{\Ptr(g\circ A)-\Qtr(g\circ A)
  -\lgp\,\E\big(\norm{\nabla g(\tilde x)}-1\big)^2\Big\},\qquad
  \tilde x=u\,A(X)+(1-u)\,A(Y),
\end{equation}
where $\Ptr$ and $\Qtr$ are empirical means over the training halves, $X$ and $Y$ are training
windows drawn independently, $u\sim U(0,1)$ is independent of them, and $\lgp=10$. The gradient
is taken with respect to the standardised input. Each optimiser step draws $b=128$ indices with
replacement from each training half and $b$ independent values of $u$, pairs the $r$-th real
and synthetic draws to form $\tilde x_r$, and takes one Adam step \citep{kingma_2015_adam} on
the mini-batch loss
$-b^{-1}\sum_{r}\{g(A X_{i_r})-g(A Y_{j_r})\}+\lgp b^{-1}\sum_r(\norm{\nabla g(\tilde x_r)}-1)^2$.
We follow \citet{gulrajani2017improved} for the penalty weight and the Adam moment parameters
$(0.5,0.9)$, but use a learning rate of $10^{-3}$. It was chosen before any E2 experiment on the
basis of an informal check, which is not part of the reported experiments or of the repository,
in which their rate of $10^{-4}$ appeared to under-train the critic within 1000 steps and to lose
power. E5 contradicts this check in one of its two designs: on simulated windows the rate
$10^{-4}$ is more powerful against a 10\% volatility increase (0.78 against 0.40; in all 19
replications in which the two rates disagree, only $10^{-4}$ rejects) and not detectably
different against a change of persistence (\cref{sec:robust}). The default was fixed before the experiments and not
re-tuned, so the critic power reported in this paper is that of a pre-specified critic, not of an
optimised one. There is no early stopping: the weights after step 1000 are frozen.

\paragraph{Statistic.}
With $\fhat$ frozen,
\begin{equation}\label{eq:critic-stat}
  T=\Pev\fhat-\Qev\fhat=\frac1{n_e}\sum_{i\in\Iev}\fhat(X_i)-\frac1{m_e}\sum_{j\in\Jev}\fhat(Y_j).
\end{equation}
The critic is trained to score real windows high, so large values of $T$ are evidence against
$H_0$, and every critic-based test is one-sided. Given the training $\sigma$-field $\Gsig$
(training halves, split, training randomness), $T$ is unbiased for the critic discrepancy
$\Dcrit(\fhat)=P\fhat-Q\fhat$, and $\Dcrit(\fhat)\le\Lip(\ghat)\,\W(P_A,Q_A)$, where $P_A$ and
$Q_A$ are the laws of the standardised windows (\cref{prop:unbiased}). If $\ghat$ were
1-Lipschitz, $T$ would therefore estimate a lower bound on the Wasserstein-1 distance between
the standardised window laws (Kantorovich--Rubinstein duality, \cref{fact:kr-duality}). The
gradient penalty enforces the Lipschitz constraint only softly and only at the interpolates, and
WGAN critics are known to be poor estimators of $\W$ itself \citep{mallasto2019well}. We
therefore read $T$ as an \emph{approximate} lower bound and use it only as a test statistic
(\cref{rem:gp-soft}). The level guarantees of \cref{sec:method-inference} do not depend on any
Lipschitz property, on the convergence of training or on the architecture: a poor critic costs
power, never level.

Why hold data out? The plug-in distance $\W(P_n,Q_m)$ and the in-sample training objective are
biased upwards, and the bias vanishes only at rate $n^{-1/d}$ in dimension $d$
\citep{fournier2015rate}. The held-out $T$ is conditionally unbiased and has an exact null
distribution under exchangeability (\cref{rem:holdout}). The price is that only half of the
data trains the critic and only half evaluates it; \cref{prop:local-power} quantifies the loss.
The held-out-witness template is that of \citet{kirchler2020two} and \citet{kubler2022witness},
and \citet{wang2022two} split and permute for Wasserstein testing. Our contribution is not the
template but its instance, its validation and the protocols of
\cref{sec:method-inference,sec:method-relative,sec:method-decision}.

\paragraph{Architectures.}
For univariate windows $\ghat$ is a multilayer perceptron $\R^L\to\R^{128}\to\R^{128}\to\R^{128}\to\R$
with LeakyReLU activations (slope 0.2) and no normalisation layers. For multivariate windows
$(N,L)$ (E4) it is a one-dimensional convolutional network with the assets as input channels.
It has two blocks, each a convolution with 64 channels, kernel size $\min(5,L)$ and ``same''
padding, followed by a LayerNorm over the $(64,L)$ feature map and a LeakyReLU. The feature map
is then flattened and passed through a dense layer of width 64, a LeakyReLU and a linear
output. Flattening, rather than pooling over time, keeps the position of each day inside the
window. BatchNorm is never used: the gradient penalty is a per-window quantity, and BatchNorm
would couple the windows of a mini-batch. The multilayer perceptron is the network of every
test on univariate windows in this paper (E1, E2, E2-dep, E3 and E5), including the anchored test
behind the verdicts. A time-pooled convolutional critic, which averages convolutional features
over the days of a window, was examined only after the real-data results were known, and was
selected by a pilot on four E2 designs that is not part of the reported experiments; this
post-hoc ablation is reported in \cref{sec:robust} and does not replace the pre-specified
perceptron.

\begin{table}[tbp]
  \centering
  \caption{Default training settings of the learned tests (the code defaults used in every
    experiment unless stated). The anchored and two-critic relative tests use the \WC{} settings.}
  \label{tab:method-defaults}
  \small
  \begin{tabular}{@{}>{\raggedright\arraybackslash}p{0.21\linewidth}
                  >{\raggedright\arraybackslash}p{0.37\linewidth}
                  >{\raggedright\arraybackslash}p{0.36\linewidth}@{}}
    \toprule
    & \WC{} & \LRC{} and C2ST \\
    \midrule
    Network, $L$-windows & MLP $L$--128--128--128--1, LeakyReLU(0.2) & same \\
    Network, $(N,L)$-windows & CNN: 2 $\times$ [Conv1d(64, kernel $\min(5,L)$), LayerNorm, LeakyReLU], dense 64 & same \\
    Input scaling & per coordinate (per asset channel), pooled training half & same \\
    Objective & WGAN-GP \eqref{eq:wgan-gp}, $\lgp=10$ & class-balanced cross-entropy \\
    Optimiser & Adam, learning rate $10^{-3}$, $(\beta_1,\beta_2)=(0.5,0.9)$ & Adam, learning rate $10^{-3}$, $(0.9,0.999)$, weight decay $10^{-4}$ \\
    Mini-batch & 128 per sample, drawn with replacement & 128 per class \\
    Steps & 1000, no early stopping & at most 1000; early stopping (20\% validation, checked every 50 steps, patience 5, best checkpoint kept) \\
    Training / evaluation & 50/50 & 50/50 \\
    Permutations $B$ & 1000 & 1000 (accuracy: randomised ties) \\
    \bottomrule
  \end{tabular}
\end{table}

% ---------------------------------------------------------------------------------------------
\subsection{Competing tests and the likelihood-ratio oracle}\label{sec:method-competitors}

\paragraph{Likelihood-ratio critic and C2ST.}
Both use the split, standardisation and architecture of the \WC{} test, but train a classifier.
Its logit $\lhat$ (the network output) minimises the class-balanced binary cross-entropy
$\tfrac12\{\Ptr\log(1+e^{-\lhat})+\Qtr\log(1+e^{\lhat})\}$, with real windows labelled 1 and
mini-batches of 128 windows per class. The optimiser is Adam with learning rate $10^{-3}$,
default moments $(0.9,0.999)$ and weight decay $10^{-4}$, for at most 1000 steps. Early stopping
holds out a random 20\% of each training half, evaluates the balanced validation cross-entropy
every 50 steps, stops after five evaluations without improvement and restores the best
checkpoint. At $n=500$ the selected checkpoint is typically the first one (step 50), because
the classifier overfits early. A pilot of alternative settings found none with higher
power at valid size (\cref{app:impl-nets}), so we keep the defaults and record the selected step instead of tuning it.

With balanced classes, the Bayes-optimal logit is $\lstar=\log(p/q)$ (\cref{prop:bayes-logit};
\citealp{goodfellow2014generative}). The \LRC{} statistic is $\Tlogit=\Pev\lhat-\Qev\lhat$. If
$\lhat=\lstar+\text{const}$, its conditional mean is the Jeffreys divergence $\Jeff(P,Q)$
(\cref{cor:jeffreys}), so $\Tlogit$ is a learned log-likelihood-ratio statistic
\citep{cheng2022classification}. Unlike $T$ it has no lower-bound property: rescaling $\lhat$
inflates it. The C2ST statistic \citep{lopezpaz2017revisiting,kim2021classification} is the
balanced accuracy of the rule $\{\lhat>0\}$,
$\Tacc=\tfrac12+\tfrac12(\Pev h-\Qev h)$ with $h=\one\{\lhat>0\}$; its conditional mean is at
most $\tfrac12\{1+\TV(P,Q)\}$. The three learned statistics are thus affine functions of a
difference of evaluation means of one function that is fixed given $\Gsig$: the critic, the
logit or the indicator $h$. This is the only structure the inference of
\cref{sec:method-inference} uses.

\paragraph{Kernel tests.}
The MMD test \citep{gretton2012kernel} uses the unbiased estimate
\begin{equation}\label{eq:mmd-u}
  \TMMD=\frac{\sum_{i\neq i'}k(X_i,X_{i'})}{n(n-1)}+\frac{\sum_{j\neq j'}k(Y_j,Y_{j'})}{m(m-1)}
  -\frac{2}{nm}\sum_{i,j}k(X_i,Y_j),\qquad k(x,y)=e^{-\norm{x-y}^2/(2\varsigma^2)},
\end{equation}
on the flattened normalised windows, with no further standardisation. The bandwidth
$\varsigma$ is the median heuristic: the median pairwise Euclidean distance in a random subset
of at most 2000 points of the pooled sample. It depends only on the pooled sample, so it is
invariant to the labels. Samples larger than 3000 windows are subsampled without replacement;
E4 fixes the bandwidth per universe (\cref{sec:mk-design}). There is no training split: the
kernel tests use all windows.

The \SigMMD{} test \citep{chevyrev2022signature,kiraly2019kernels} maps a window
$(r_1,\dots,r_L)$ to the piecewise-linear path through the points
$(k/L,\sum_{l\le k}r_l)$, $k=0,\dots,L$: time and cumulative return. Its signature truncated at
depth 4 is computed exactly by Chen's identity \citep{lyons1998differential,chevyrev2016primer},
which gives $2+4+8+16=30$ features. Each feature is standardised with the mean and standard
deviation of the pooled sample, and $\Tsig$ is the statistic \eqref{eq:mmd-u} on these features,
with its own median-heuristic bandwidth.

\paragraph{Kolmogorov--Smirnov.}
$\TKS$ is the two-sample KS statistic between the $nL$ pooled real returns and the $mL$ pooled
synthetic returns, with a two-sided asymptotic p-value. It sees only the one-dimensional
marginal law of a daily return, so it has no power against deviations that leave this marginal
unchanged (\cref{fig:e2dep-power}). Its p-value assumes \iid{} returns. That assumption fails
when the returns of a window share a conditional variance, which makes KS oversized on GARCH
windows (\cref{tab:e2-size}), and it fails again across dependent windows. We keep KS because
practitioners use it, and because against deviations of the marginal (volatility level, tails,
jumps) it is, by a wide margin, the most powerful of the two-sample tests we compare
(\cref{sec:sim-power}).
Any honest use of a window test has to be weighed against this baseline.

\paragraph{Neyman--Pearson oracle (E2 only).}
When $P$ and $Q$ have known window densities $p,q$, the oracle tests $Y\sim P$ against
$Y\sim Q$ on the synthetic sample with
\begin{equation}\label{eq:lr-oracle}
  \TLR=\sum_{j=1}^m\log\frac{q(Y_j)}{p(Y_j)}.
\end{equation}
Its null distribution is computed from 2000 independent samples of $m$ windows drawn under $P$,
and the Monte Carlo p-value with randomised tie-breaking has exact level
\citep{dufour_2006_monte}. For i.i.d.\ returns (volatility level, tails, jumps) the window
densities are exact, so this is the most powerful test of $P$ against $Q$, and its power bounds
that of every level-$\alpha$ two-sample test (\cref{prop:oracle-bound}). For GARCH windows the
density has no closed form. We integrate over the conditional variance at the start of the
window with $M=512$ Monte Carlo draws from its law after the burn-in. The draws are shared by
the observed and the null samples, so the test keeps exact level, and $\TLR$ approaches the
exact log-likelihood ratio as $M$ grows.

% ---------------------------------------------------------------------------------------------
\subsection{Inference for independent and for dependent windows}\label{sec:method-inference}

For the learned tests, conditionally on $\Gsig$, every statistic is a difference
$\bar a-\bar b$ of evaluation means of a fixed function $g$, with $a_k=g(X_k)$, $k\in\Iev$, and
$b_k=g(Y_k)$, $k\in\Jev$. Two protocols turn this difference into a p-value.

\paragraph{Protocol for independent windows (E2).}
Pool the evaluation values, $(c_1,\dots,c_{n_e+m_e})=(a,b)$. For $b=1,\dots,B$, draw a uniform
random permutation of the pooled vector, independently of everything else, and let $T_b$ be the
statistic recomputed with the permuted labels. The p-value is
\begin{equation}\label{eq:perm-pvalue}
  \pperm=\frac{1+\#\{b\le B:\ T_b\ge T\}}{1+B},
\end{equation}
with $B=1000$; ties with $T$ count against it (up to a relative tolerance of $10^{-12}$). The
C2ST accuracy is discrete, so for it ties are broken at random:
$\pperm=(\#\{T_b>T\}+U\,(1+\#\{T_b=T\}))/(1+B)$ with $U\sim U(0,1)$, which is exactly uniform
under exchangeability. If the windows are \iid{} and the two samples are independent, the pooled
evaluation windows are exchangeable given $\Gsig$ under $H_0$. The test then has level at most
$\alpha$ in finite samples, for every critic, every $B$ and every split that does not depend on
the data (\cref{prop:exact-level}; \citealp{hemerik2018exact}). The kernel tests have no training
split. They permute the labels of the pooled sample with the kernel matrix and the bandwidth held
fixed ($B=500$), which is exact under the same conditions.

\paragraph{Why real data need another protocol.}
Two things go wrong when windows are dependent (\cref{sec:theory-dependence}).
\begin{enumerate}
  \item \emph{Variance.} The real evaluation windows of a ticker are autocorrelated through
    volatility clustering, and windows of different tickers on the same dates are correlated
    through the market. The permutation distribution ignores both, so the test is oversized
    (\cref{rem:naive-perm}).
  \item \emph{Bias.} If training and evaluation windows are close in time, the critic learns the
    realised volatility path of the training dates (``high volatility in spring 2023 means
    real''). Real evaluation windows from the same or neighbouring dates share that path;
    synthetic windows do not. Then $\E[T\mid\Gsig]>0$ under $H_0$. A random split of
    consecutive windows creates this coupling. A split by ticker in a panel creates it in full,
    because every training date is also an evaluation date (\cref{prop:ticker-bias}). No
    variance estimator can remove a bias.
\end{enumerate}
The naive protocol, a ticker split with permutation p-values, is run in E1 as a comparison.
In the simulated placebo panel it rejects a true null at nominal 5\% in all 200 replications
(\cref{tab:e1-size-sim}).

\paragraph{Protocol for dependent windows (E1, E3, E4, E5 and E2 part g).}
It has three components.

\emph{Time-block split.} The date range $[\tau_{\min},\tau_{\max}]$ of each sample is cut into
$K=10$ intervals of equal length: a window with date rank $\tau$ gets the block label
$\min\{K-1,\lfloor K(\tau-\tau_{\min})/(\tau_{\max}-\tau_{\min})\rfloor\}$. Labels are shared
across samples, so the blocks of $X$ and of $Y$ that carry the same label always fall in the
same half.
The $K$ labels are shuffled, and the training half receives them in that order up to the point
where the cumulative number of windows (over all samples) is closest to half the total. In E3,
$X$ and $Y$ both cover the 30 \TEST{} window dates. Every block then holds three dates (300
windows per sample), five blocks train and five evaluate, and each half has 15 dates and 1500
windows per sample. Training and evaluation windows never share a date, which removes the
common-date bias of \cref{prop:ticker-bias}. Two features of the implementation limit what the
theory covers. The blocks interleave in time, and there is no gap between a training block and
an adjacent evaluation block. The dependence across these boundaries is therefore not
controlled by \cref{lem:decoupling}, which needs a gap. We measure its effect instead
(\cref{tab:e2-dependent-size,tab:e1-size-sim}).

\emph{Driscoll--Kraay variance.} Let $\tau_k$ be the date rank of evaluation window $k$ and
$u_k=(\tau_k-\min_{k'}\tau_{k'})/L$ its position on the grid of window start dates, of length
$U=\max_ku_k+1$. Positions without an evaluation window (the training blocks) count as zeros.
With the date sums $\mathcal S_u=\sum_{k:\,u_k=u}(a_k-\bar a)$,
\begin{equation}\label{eq:dk-var}
  \widehat V_{\mathrm{DK}}(\bar a)=\frac{1}{n_e(n_e-1)}\Big[\sum_{u}\mathcal S_u^2
  +2\sum_{h=1}^{H}\Big(1-\frac{h}{H+1}\Big)\sum_{u}\mathcal S_u\mathcal S_{u+h}\Big].
\end{equation}
Summing within a date absorbs arbitrary cross-sectional dependence \citep{driscoll1998consistent},
and the Bartlett weights \citep{newey1987simple} handle serial dependence. The factor
$n_e/(n_e-1)$, folded into the normalisation, makes \eqref{eq:dk-var} equal to the usual \iid{}
variance estimator when $H=0$ and all dates are distinct. Lags are counted in windows of $L$
trading days, never in days. The truncation lag $H$ follows the AR(1) plug-in rule of
\citet{andrews1991heteroskedasticity} for the Bartlett kernel. The coefficient $\hat\rho$ is the
lag-one autocorrelation of $(\mathcal S_u)$, computed only over adjacent grid positions that both
carry evaluation windows (so that the holes left by the training blocks are not mistaken for low
persistence) and clipped to $[-0.97,0.97]$. Then
$\hat\eta=4\hat\rho^2/\{(1-\hat\rho)^2(1+\hat\rho)^2\}$ and
$H=\lceil1.1447(\hat\eta U)^{1/3}\rceil-1$, truncated to $[0,U-1]$. If fewer than ten adjacent
pairs are observed, the rule is applied to the sequence of observed dates. The rule of
\citet{newey1994automatic} is an E5 robustness option. The studentised statistic and its
p-value are
\begin{equation}\label{eq:hac-z}
  \zhac=\frac{\bar a-\bar b}{\sqrt{\widehat V_{\mathrm{DK}}(\bar a)+\widehat V_{\mathrm{DK}}(\bar b)}},
  \qquad \phac=1-\Phi(\zhac).
\end{equation}
The two variances add because the synthetic sample is independent of the real one. For a fixed
critic and a growing number of evaluation dates, $\zhac$ is asymptotically $N(0,1)$ under
$H_0$ when the window sequences are mixing (\cref{prop:hac}). With only 15 evaluation dates the
Driscoll--Kraay estimate is noisy and biased downwards, so some oversize remains. E1 measures the
size of exactly this pipeline and reports fixed-$b$ p-values \citep{kiefer_2005_new} as
comparison rows (\cref{rem:small-T}). At nominal 5\% the \WC{} test rejects a true null at 0.085
(the anchored relative test at 0.080) when the synthetic panel shares the factor structure of the
real one, but at 0.270 (anchored: 0.155) when the synthetic tickers are drawn independently, as
the univariate generators of E3 draw them (\cref{tab:e1-size-sim}). In that case the studentised
statistic is also shifted, not only over-dispersed (\cref{rem:small-T}), so a better variance
estimate alone would not remove the excess; \cref{sec:emp} takes this into account. Under this protocol the code also computes the permutation
p-value under the time-block split. E1 reports its size, but it is never used for inference.

\emph{Wild bootstrap for the kernel tests.} The kernel tests have no split and no witness
function, so they cannot be studentised in this way. They use the wild bootstrap for degenerate
V-statistics of \citet{chwialkowski2014wild}. Let $\mathsf K$ be the pooled kernel matrix and
$v=(\one_n/n,-\one_m/m)$. The p-value refers to the V-statistic $v^\top\mathsf Kv$, the squared
distance between the two empirical kernel mean embeddings. In replicate $b$, a stationary
Gaussian AR(1) process $(\zeta_u)$ with $\Cov(\zeta_u,\zeta_{u'})=\exp(-\abs{u-u'}/\ell_w)$ is
drawn on the pooled grid of window dates. Window $k$ receives $\zeta_{u_k}$, so windows that
start on the same date share one multiplier, in both samples. The multipliers are centred
within each sample, $\tilde\zeta$, and the replicate is $\tilde v^\top\mathsf K\tilde v$ with
$\tilde v=(\tilde\zeta_X/n,-\tilde\zeta_Y/m)$. The p-value $\pwild$ is the share of the
$B=500$ replicates at or above $v^\top\mathsf Kv$, computed as in \eqref{eq:perm-pvalue}. The
correlation length is $\ell_w=(\hat\eta U)^{1/3}$, where $\hat\eta$ is now the several-series
version of the Andrews constant \citep{andrews1991heteroskedasticity}, computed from the date
sums of the scores on the five leading kernel principal components ($\ell_w=0$, i.e.\
independent multipliers, when no dependence is detected). This plug-in is our heuristic:
\citet{chwialkowski2014wild} leave $\ell_w$ as a tuning parameter, and we give no proof of
validity in our setting (\cref{rem:wild}). The tables report $\TMMD$ and $\Tsig$ as the
unbiased statistic \eqref{eq:mmd-u}, while $\pwild$ refers to the V-statistic.

\emph{KS} is computed and reported, but flagged as invalid under this protocol.

\begin{table}[tbp]
  \centering
  \caption{The tests as implemented. ``Dependent windows'' is the protocol of real data and of
    windows cut from long simulated paths; ``DK'' is the Driscoll--Kraay studentisation
    \eqref{eq:hac-z} after the time-block split.}
  \label{tab:method-tests}
  \footnotesize
  \setlength{\tabcolsep}{4pt}
  \begin{tabular}{@{}>{\raggedright\arraybackslash}p{0.13\linewidth}
                  >{\raggedright\arraybackslash}p{0.27\linewidth}
                  >{\raggedright\arraybackslash}p{0.15\linewidth}
                  >{\raggedright\arraybackslash}p{0.18\linewidth}
                  >{\raggedright\arraybackslash}p{0.19\linewidth}@{}}
    \toprule
    Test & Statistic (target) & Learned on training half & p-value, independent windows &
    p-value, dependent windows \\
    \midrule
    \WC{} & $T=\Pev\fhat-\Qev\fhat$ (approximate lower bound on $\W(P_A,Q_A)$) & WGAN-GP critic &
      permutation, $B=1000$ & DK $\phac$ \\
    \LRC{} & $\Tlogit=\Pev\lhat-\Qev\lhat$ ($\Jeff(P,Q)$ if $\lhat=\lstar$) & classifier &
      permutation, $B=1000$ & DK $\phac$ \\
    C2ST & balanced accuracy $\Tacc$ (at most $\tfrac12\{1+\TV(P,Q)\}$) & classifier &
      permutation, randomised ties & DK $\phac$ \\
    MMD & $\TMMD$, Gaussian kernel on windows & none & permutation, $B=500$ &
      wild bootstrap, $B=500$ \\
    \SigMMD{} & $\Tsig$, Gaussian kernel on depth-4 signatures & none & permutation, $B=500$ &
      wild bootstrap, $B=500$ \\
    KS & $\TKS$ on pooled returns (one-dimensional marginal) & none & asymptotic (needs \iid{}
      returns) & reported, marked invalid \\
    \addlinespace
    Anchored relative & $\Danc=\Ppast^{\evl}\fhat-\QG^{\evl}\fhat$ & one critic, $\Rnew$ vs pool &
      permutation of $G$ / $\Rpast$ labels & DK (verdict p-value) \\
    Two-critic relative & $\Dtwo$, \eqref{eq:two-critic-stat} & two critics &
      paired bootstrap, $B=500$ & DK (secondary) \\
    \addlinespace
    LR oracle (E2) & $\TLR$, \eqref{eq:lr-oracle} & none (true densities) &
      Monte Carlo, 2000 null samples & -- \\
    \bottomrule
  \end{tabular}
\end{table}

% ---------------------------------------------------------------------------------------------
\subsection{Relative tests}\label{sec:method-relative}

\paragraph{The anchored relative test.}
The inputs are $\Rnew$ ($X$, $n$ windows), the synthetic sample $G$ ($Y$, $m$ windows) and
$\Rpast$ ($Z$, $k$ windows).
\begin{enumerate}
  \item \emph{Split.} $Y$ and $Z$ are split into training and evaluation halves, of sizes
    $m_t,m_e$ and $k_t,k_e$: by time blocks under the dependent protocol, at random otherwise.
    The block labels are shared, so a block of $G$ (one tenth of \TEST{}) and the block of
    $\Rpast$ with the same label (the corresponding tenth of \TRAIN{}) fall in the same half; in
    E3 this puts about half of each sample in each half. $\Rnew$ is not split: all of it is
    used for training, and none of it for evaluation.
  \item \emph{Training.} One critic $\fhat=\ghat\circ A$, with the architecture and optimiser of
    the \WC{} test and $A$ fitted on all training windows, is trained to separate $\Rnew$
    (positive class) from the pool $Y^{\trn}\cup Z^{\trn}$ (negative class). Inside the pool the
    two parts carry equal total weight, whatever their sizes:
    \begin{equation}\label{eq:anchored-obj}
      \ghat\approx\argmax_g\Big\{\widehat P_{\mathrm{new}}(g\circ A)
      -\tfrac12\QG^{\trn}(g\circ A)-\tfrac12\Ppast^{\trn}(g\circ A)
      -\lgp\,\E\big(\norm{\nabla g(\tilde x)}-1\big)^2\Big\},
    \end{equation}
    where $\widehat P_{\mathrm{new}}$ is the empirical measure of all $\Rnew$ windows and
    $\QG^{\trn},\Ppast^{\trn}$ are those of the two training halves. In practice each mini-batch
    draws pool windows with replacement, with probability $1/(2m_t)$ for each window of
    $Y^{\trn}$ and $1/(2k_t)$ for each window of $Z^{\trn}$, and the penalty is evaluated at
    interpolates between an $\Rnew$ window and a pool window.
  \item \emph{Statistic.} With $\fhat$ frozen,
    \begin{equation}\label{eq:anchored-stat}
      \Danc=\Ppast^{\evl}\fhat-\QG^{\evl}\fhat
      =\frac1{k_e}\sum_{Z_l\in Z^{\evl}}\fhat(Z_l)-\frac1{m_e}\sum_{Y_j\in Y^{\evl}}\fhat(Y_j).
    \end{equation}
\end{enumerate}
The critic scores ``$\Rnew$-likeness''. $\Danc>0$ means that real past windows look more like
the new data than synthetic windows do: $G$ is worse than simply replaying history. The test is
one-sided against $\Danc>0$, and its conditional target is $\delta=\Ppast\fhat-\QG\fhat$, with
null $\delta\le0$ (\cref{def:rel-null}). The name ``anchored'' refers to the role of $\Rnew$.
For one shared critic, the difference of the two absolute discrepancies is
$(\Pnew\fhat-\QG\fhat)-(\Pnew\fhat-\Ppast\fhat)=\Ppast\fhat-\QG\fhat$: the $\Rnew$ terms cancel
exactly. $\Rnew$ anchors what the critic looks for, but is never evaluated, which is why all of
it can be used for training.

\emph{p-values.} For independent windows, the labels of the pooled $G$ and $\Rpast$ evaluation
windows are permuted with $\fhat$ fixed ($B=1000$), and the p-value is \eqref{eq:perm-pvalue}
with $(a,b)=(\fhat(Z^{\evl}),\fhat(Y^{\evl}))$. At the boundary $\Law(G)=\Law(\Rpast)$ this test
is exact for every critic and every law of $\Rnew$ (\cref{prop:anchored}). Given $\Gsig$, the
$G$ and $\Rpast$ evaluation windows are then \iid{} from one law, whatever the training did. The
pool weights affect power only. In the interior of the null, with unequal evaluation sizes, the
plain permutation is only approximately calibrated. The code therefore also reports the
studentised permutation p-value (Welch statistic recomputed under each permutation), which is
asymptotically valid for $\delta\le0$ and still exact at the boundary
\citep{janssen1997studentized,chung2013exact}. For dependent windows the p-value is
\eqref{eq:hac-z} with $(a,b)=(\fhat(Z^{\evl}),\fhat(Y^{\evl}))$, each Driscoll--Kraay variance
computed on its own date ranks. This HAC p-value enters the verdict, and the reported interval
is $\Danc\pm1.96\,\widehat{\mathrm{se}}_{\mathrm{DK}}$.

\emph{Caveats.} (i) $G$ is calibrated on $\Rpast$, so independence between $G$ and the $\Rpast$
windows is an idealisation; for parametric generators each past window moves the fitted
parameters by $O(1/k)$, and we treat the effect as negligible without proof
(\cref{rem:rel-guarantee}). (ii) The stationary bootstrap generator copies blocks of $\Rpast$.
Some synthetic evaluation windows then duplicate windows that the critic saw as pool negatives,
exchangeability fails, and the test is invalid; its verdict is \INVALID{}. E2 part e runs this
case with a naive and with a grouped split (\cref{tab:e2-relative-boundary}). (iii) On real data
$\Rpast$ is not stationary (2020, 2022). The boundary is then the pooled \TRAIN{} window law,
and the time-block split compares the critic's view of $G$ with a random half of the \TRAIN{}
history. A single-seed p-value is not calibrated in this situation. The verdict therefore uses
the seed-combined p-value of \cref{sec:method-decision}, whose power and boundary size are
estimated under a design with the real \TRAIN{} volatility path (design D2 below). (iv) The last \TRAIN{} block
borders the first \TEST{} dates. Under volatility clustering the critic can learn a regime
shared by $\Rnew$ and the end of $\Rpast$, a small effect in the direction of rejection.

\paragraph{The two-critic relative test (secondary).}
All three samples are split, with shared block labels under the dependent protocol. Two critics
are trained separately with the same settings: $\fhat_G$ on $(X^{\trn},Y^{\trn})$ and $\fhat_P$
on $(X^{\trn},Z^{\trn})$. The statistic is the difference of the two held-out discrepancies,
\begin{equation}\label{eq:two-critic-stat}
  \Dtwo=\big(\Pev\fhat_G-\QG^{\evl}\fhat_G\big)-\big(\Pev\fhat_P-\Ppast^{\evl}\fhat_P\big),
\end{equation}
a sum of three evaluation means given $\Gsig$. The p-value comes from a paired bootstrap with
$B=500$ replicates. In each replicate the $\Rnew$ evaluation windows are resampled once and used
for both critics, and the $G$ and $\Rpast$ evaluation windows are resampled independently. The
p-value is centred, $(1+\#\{\Dtwo_{*b}-\Dtwo\ge\Dtwo\})/(1+B)$ \citep{hall1991two}, and the
interval is the percentile interval; under the dependent protocol the three Driscoll--Kraay
variances are added instead. This test is valid conditionally on the two trained critics
(\cref{prop:paired-boot}) but not exact at the boundary $\Law(G)=\Law(\Rpast)$, where its
unconditional size is at least $\alpha$ (\cref{prop:boundary-size}). There, the conditional target $\Dcrit_G-\Dcrit_P$ is the difference
of two exchangeable replicates of a quantity that depends on training. It is symmetric about
zero, with a spread of the order of the evaluation noise, and for fixed critics and growing
evaluation sets the rejection probability tends to $1/2$. In E2 part e, with $n=500$, the excess
was too small to be seen (\cref{sec:sim-relative}). We report the test for comparison only; it
never enters the verdict.

\paragraph{Kernel version.}
For MMD and \SigMMD{} there is no learned witness. The relative test then follows
\citet{bounliphone2016test}: $\TMMD(\Rnew,G)-\TMMD(\Rnew,\Rpast)$ with one kernel (median
heuristic on the three pooled samples). Its normal p-value uses the variance of the first-order
Hoeffding projections, with Driscoll--Kraay variances of the projections under dependence. The
approximation degenerates when all three laws coincide. This version is available in the code
but is not used for any verdict.

% ---------------------------------------------------------------------------------------------
\subsection{The decision rule}\label{sec:method-decision}

\paragraph{Seeds.}
A single anchored p-value depends on the synthetic path and on the randomness of the test (block
assignment, initialisation, mini-batches). Reporting the most favourable of several draws would
be invalid. E3 therefore runs $S=5$ synthetic seeds per generator. Seed $s$ draws a fresh
synthetic path for every ticker and fresh test randomness, and the seeds are shared across
generators, so at a given seed every generator faces the same block assignment and critic
initialisation (common random numbers). With $p_s$ the HAC p-value of the anchored test at seed
$s$, the seeds are combined by
\begin{equation}\label{eq:pmed}
  \pmed=\min\big\{1,\;2\med(p_1,\dots,p_S)\big\},
\end{equation}
which is a valid p-value under arbitrary dependence between the $p_s$
\citep{meinshausen_2009_pvalues}; see \cref{rem:combine-seeds}. With $S=5$ the median is the
third smallest $p_s$.

\paragraph{Power check.}
A non-rejection is evidence of quality only if the test could have detected a deviation of
relevant size. The power check (\cref{def:power-check}) simulates the whole rule, five seeds
included, with the anchored \WC{} test, the dependent protocol and the real sample sizes,
scales, date ranks and split. The pre-specified deviation $\vartheta^*$ multiplies synthetic
volatility by 1.2. There are two reference designs.

\emph{D1 (stationary).} The reference model is the per-ticker \GARCHt{} fitted on
\TRAIN{} (the E3 fits). In each replication, $\Rpast^*$ is a \TRAIN{}-length path started
from stationarity (250 burn-in days), and $\Rnew^*$ is an independent \TEST{}-length path
started from the fitted end-of-\TRAIN{} state. For each seed the alternative $G^*$ is 1.2
times an independent \TEST{}-length draw from the end-of-\TRAIN{} state. The boundary
generator $G^0$ is an independent \TEST{}-length draw started from stationarity, so that
$\Law(G^0)=\Law(\Rpast^*)$ exactly. $G^*$ and $G^0$ are tested with the same test randomness.
The D1 alternative thus differs from the boundary in its start state as well as in its
scale, so its power is not that of the factor 1.2 alone.

\emph{D2 (real volatility path).} Let $v_{ij}$ be the root mean square of the
normalised real \TRAIN{} window $j$ of ticker $i$. Window $(i,j)$ of $\Rpast^*$ is
$v_{ij}\,\xi$ with $\xi\sim N(0,I_L)$, drawn once per replication; this keeps the 2020
and 2022 regimes and their co-movement across tickers. $\Rnew$ is the real \TEST{} sample.
For each seed, the $j'$-th window of ticker $i$ in $G^0$ is $v_{iU}\,\xi'$ with $U$ uniform
over the \TRAIN{} windows of ticker $i$, so $G^0$ has exactly the pooled \TRAIN{} window law
of each ticker; it carries the \TEST{} date ranks. The alternative is $G^*=1.2\,G^0$. The
boundary of D2 is a placebo of the anchored test under the real non-stationarity of
$\Rpast$, with two simplifications. Within a window the returns of $\Rpast^*$ and $G^0$ are Gaussian with a
constant scale, and the noise is independent across tickers and windows; $G^0$ also draws
its scales independently across tickers and windows, so it has no within-window tails or
clustering and no regime common to several tickers. The D2 boundary size is therefore an
estimate for this approximation of the pooled past, not for every generator at the
boundary, in particular not for fat-tailed generators with volatility clustering such as
\GARCHt{}.

Each design runs $R=100$ replications. The estimated power $\hat\beta_{\mathrm{D}}$ is the share
of replications with $\pmed\le\alpha$ under the alternative; the boundary size
$\hat\alpha_{\mathrm{bd},\mathrm{D}}$ is the same share under $G^0$. The power is reported with
its Wilson interval and its one-sided Clopper--Pearson lower bound $\betaL$
\citep{clopper_1934_use}, the boundary size with its Clopper--Pearson interval. The power of the
rule is the smaller of the two, $\hat\beta=\min(\hat\beta_{\mathrm{D1}},\hat\beta_{\mathrm{D2}})$. The designs, the deviation and
the threshold were fixed, and both designs were run, before the full real-data tests; D2 and the
seed rule entered the protocol on the day of those runs, when the real-data tests had been run
only in smoke runs on a six-ticker subset. The check does not involve the generators being
evaluated, but D2 uses the real \TEST{} sample, so its power and boundary size are conditional on
the realised \TEST{} sample and \TRAIN{} volatility path (\cref{tab:e3-power}).

\paragraph{Verdict.}
With $\alpha=0.05$ and $\bmin=0.8$,
\begin{equation}\label{eq:verdict}
  V=\begin{cases}
    \BAD & \text{if } \pmed\le\alpha,\\
    \GOOD & \text{if } \pmed>\alpha \text{ and } \hat\beta\ge\bmin,\\
    \INCONCL & \text{otherwise,}
  \end{cases}
\end{equation}
and $V=\INVALID$ for the stationary bootstrap, whose relative test is invalid (caveat (ii)
above). \Cref{def:verdict} restates the rule formally, and \cref{prop:verdict} states what it
controls. The variant that requires the Clopper--Pearson bounds of both designs to reach $\bmin$
is reported alongside. Each verdict comes with $\pmed$ of the absolute \WC{} test over the same
seeds, the median and range of $\Danc$ over seeds, and the boundary sizes of both designs
(\cref{tab:e3-verdicts}). The combination ``absolute test rejects, relative test does not''
reads ``distinguishable from the future, but not worse than history''. It is a finding, not a
failure of the method.

\paragraph{What the verdict means.}
\BAD{} is a rejection of the anchored null for the trained critics. With independent windows,
the permutation version of the anchored p-value would bound its error rate at the boundary by
$\alpha$ in finite samples (\cref{prop:anchored}(a)), and $\pmed$ would inherit the bound
(\cref{rem:combine-seeds}). The rule as implemented combines Driscoll--Kraay p-values instead,
which are only asymptotically valid (\cref{prop:anchored}(c)), under conditions that the
interleaved split meets only approximately (\cref{sec:method-inference}). The seed combination
is valid only if the single-seed p-values are. In D2 they are not: the single-seed boundary size
is 0.25 [0.17, 0.35] (0.09 [0.04, 0.16] in D1), and the boundary size of the rule is
0.08 [0.04, 0.15] (0.00 [0.00, 0.04] in D1; Clopper--Pearson 95\% intervals,
\cref{tab:e3-power}). On real data the error rate of \BAD{} at the boundary is therefore this
estimated boundary size, which we report rather than assume, and a \BAD{} verdict has a real-data
error interpretation only when the Clopper--Pearson interval of the estimated boundary size covers
$\alpha$ or lies below it (\cref{prop:verdict}(a)), as at the main split. The size-adjusted power
of \cref{def:power-check} separates the discrimination of the rule from this miscalibration
(\cref{tab:manual-theory-calibration}), and \cref{sec:emp-verdicts} also compares each observed
$\pmed$ with the D2 boundary replications.

\GOOD{} means ``no detectable deviation of the pre-specified size and type in designs that
resemble the data''. It never means ``equal to the truth''. Its guarantee rests on a premise
that cannot be verified: the actual deviation must be at least as detectable as a 20\%
volatility error in D1 and D2 (\cref{prop:verdict}). The premise is also one-directional.
$\vartheta^*$ is an upward volatility error; in D2, where $\Rnew$ is the real \TEST{} sample and
\TEST{} was calmer than \TRAIN{}, it moves the generator away from the evaluation period, which
is the direction in which the anchored test is built to reject. The check covers neither
downward volatility errors nor errors in other features, such as tails or clustering, and its
power against a factor 0.8 was not estimated. A downward
error moves $G$ towards \TEST{} and need not violate the anchored null at all, since that null
asks only whether $G$ is further from $\Rnew$ than the past is (\cref{rem:anchored-null}). Under
a non-stationary past, a non-rejection can therefore reflect a generator bias that happens to
point in the direction of the regime change, and a \GOOD{} verdict would credit that bias as
quality. In E3 the rule power in D2 is 0.45, so no generator can be \GOOD{} and this restriction
changes no reported verdict, but it limits what a \GOOD{} verdict means in any application with
a regime change between the two periods. The per-ticker tests of E3 ($L=5$,
calibrated p-values, Benjamini--Hochberg; \cref{sec:emp-perticker,rem:bh}) are descriptive and
do not enter the verdict. \Cref{fig:alg-verdict} collects the steps.

\begin{figure}[tbp]
  \centering
  \fbox{\begin{minipage}{0.95\linewidth}
  \small
  \textbf{Input.} Real \TRAIN{} and \TEST{} panels (100 tickers); a generator class
  $\mathcal M$; $L=20$, $K=10$, $S=5$, $\alpha=0.05$, $\bmin=0.8$, $R=100$.
  \begin{enumerate}
    \setlength{\itemsep}{1pt}
    \item \emph{Power check} (once, for all generators). For $\mathrm D\in\{\mathrm{D1},\mathrm{D2}\}$
      and $r=1,\dots,R$: draw $\Rpast^*$ (and $\Rnew^*$ in D1); for $s=1,\dots,S$ draw $G^*_s$ and
      $G^0_s$ and run step 4 on each; set $p^{\mathrm{alt}}_r$ and $p^{\mathrm{bd}}_r$ by
      \eqref{eq:pmed}. Then $\hat\beta_{\mathrm D}=R^{-1}\sum_r\one\{p^{\mathrm{alt}}_r\le\alpha\}$,
      $\hat\alpha_{\mathrm{bd},\mathrm D}=R^{-1}\sum_r\one\{p^{\mathrm{bd}}_r\le\alpha\}$ and
      $\hat\beta=\min_{\mathrm D}\hat\beta_{\mathrm D}$.
    \item \emph{Fit.} Fit $\mathcal M$ on the \TRAIN{} returns of each ticker.
    \item \emph{Simulate} (for $s=1,\dots,S$). One \TEST{}-length path per ticker with seed $s$,
      from the fitted end-of-\TRAIN{} state where the model has one; divide by $s_i$; cut
      windows; assign \TEST{} date ranks.
    \item \emph{Anchored test} (for each $s$). Time-block split of $G_s$ and $\Rpast$ (shared
      labels); train one \WC{} on all of $\Rnew$ against the pool with weights $\tfrac12,\tfrac12$
      \eqref{eq:anchored-obj}; compute $\Danc_s$ \eqref{eq:anchored-stat} and
      $p_s=1-\Phi(\zhac)$ from \eqref{eq:dk-var}--\eqref{eq:hac-z}.
    \item \emph{Combine.} $\pmed=\min\{1,2\med_sp_s\}$.
    \item \emph{Verdict.} \INVALID{} if $G$ resamples $\Rpast$ (stationary bootstrap); else
      \BAD{} if $\pmed\le\alpha$, \GOOD{} if $\hat\beta\ge\bmin$, \INCONCL{} otherwise.
    \item \emph{Report.} $V$; the Clopper--Pearson variant; $\hat\beta_{\mathrm D}$ and
      $\hat\alpha_{\mathrm{bd},\mathrm D}$; the median and range of $\Danc_s$; $\pmed$ of the
      absolute \WC{} test (time-block split, $\phac$) over the same seeds.
  \end{enumerate}
  \end{minipage}}
  \caption{The decision rule for one generator, as run in E3. Step 1 does not involve the
    generators under evaluation.}
  \label{fig:alg-verdict}
\end{figure}

\section{Theory}\label{sec:theory}
% =============================================================================================
% Owner: theory writer. Source: docs/theory.md (Sections 1-9), checked against src/synthval.
% Proofs: sections/appendix_proofs.tex (input by appendix.tex under \section{Proofs}).
% Labels follow paper/OUTLINE.md, Section 6. New labels (not in the outline registry):
%   sec:theory-markowitz, ass:anc, ass:rel, prop:anchored-target, rem:anchored-null,
%   rem:nonstat-past, rem:ks-marginal, tab:theory-status, app:proofs-* (appendix subsections),
%   cor:block-bias, rem:calibrated-bad, tab:manual-theory-calibration
%   (paper/scripts/theory_calibration.py), tab:manual-e1-sim-stat (paper/scripts/e1_sim_stat.py).
% =============================================================================================

This section states what the protocol of \cref{sec:method} guarantees and what it does not. A
Proposition, Lemma or Corollary is proved, in \cref{app:proofs} unless the proof takes a few
lines; ``sketch'' marks a proof that gives every key step but omits routine measure-theoretic
detail. A Fact is a standard result quoted with a reference. A Remark whose title contains
``heuristic'' is an argument, not a theorem. \Cref{tab:theory-status} at the end of the section
lists the status and the assumptions of each result. Throughout, $P$ and $Q$ are the laws of one
real and one synthetic window on $\R^d$ ($d=L$, or $d=NL$ for multivariate windows), $\norm{\cdot}$
is the Euclidean (Frobenius) norm, $Pf=\int f\,dP$, and $\Pev f$, $\Qev f$, $\Ptr f$, $\Qtr f$
are means of $f$ over the evaluation and training halves of the two samples.

% ---------------------------------------------------------------------------------------------
\subsection{What a window test can detect}\label{sec:theory-windows}

Let $(r_t)_{t\in\mathbb Z}$ be a return process with values in $\R$ or $\R^N$. The real sample
consists of non-overlapping windows $X_i=(r_{s_0+(i-1)L+1},\dots,r_{s_0+iL})$ of the observed
series, and the synthetic sample of windows $Y_j$ of paths simulated from the generator $G$.

\begin{definition}[Absolute null]\label{def:abs-null}
With $P=\Law(X_i)$ and $Q=\Law(Y_j)$, the absolute null is $H_0:P=Q$, against $H_1:P\neq Q$.
\end{definition}

$H_0$ concerns the marginal law of one window. Dependence between windows, such as successive
windows of one path or windows of different tickers on the same dates, is not part of $H_0$; it is
a nuisance that changes the null distribution of every statistic.
\Cref{sec:theory-critic,sec:theory-level,sec:theory-consistency,sec:theory-likelihood} assume
independent windows and \cref{sec:theory-dependence} removes this assumption.

\begin{proposition}[What a window law identifies]\label{prop:window-ident}
Let $(r_t)$ be strictly stationary and let $P_L$ be the law of $(r_1,\dots,r_L)$.
\begin{enumerate}
  \item[(a)] $P_L$ determines the joint law of $(r_{t_1},\dots,r_{t_k})$ whenever
  $t_k-t_1\le L-1$: every autocovariance and every moment $\E g(r_t,r_{t+h})$ with
  $0\le h\le L-1$ (the autocorrelations of $r_t$, $\abs{r_t}$ and $r_t^2$, leverage correlations),
  and the law of aggregated returns $\sum_{t=1}^h r_t$ for $h\le L$.
  \item[(b)] The family $(P_L)_{L\ge1}$ determines the law of the process.
  \item[(c)] For every $L\ge1$ there are strictly stationary processes with the same $P_L$ and
  different $P_{L+1}$.
\end{enumerate}
\end{proposition}
\noindent\textit{Proof.} In \cref{app:proofs-windows}.

A test on windows of length $L$ therefore tests the necessary condition $P_L=Q_L$: equality of
every dependence feature at lags $0,\dots,L-1$ and of aggregated returns up to horizon $L$. It has
no power against deviations that appear only at lags $L$ or more, however large the sample.
Longer windows see longer-range dependence but give fewer windows ($n=\lfloor T/L\rfloor$) in a
higher dimension; E5 varies $L\in\{5,10,20,60\}$ (\cref{sec:robust}).

\begin{proposition}[Pairs determine a stationary Markov chain]\label{prop:markov-pairs}
Let $(r_t)$ and $(r'_t)$ be strictly stationary, time-homogeneous, first-order Markov chains with
values in a Polish space. If $(r_0,r_1)\eqd(r'_0,r'_1)$, then $(r_t)\eqd(r'_t)$ as processes. For
chains of order $k$, $P_{k+1}=P'_{k+1}$ implies equality of the processes.
\end{proposition}
\noindent\textit{Proof.} In \cref{app:proofs-windows}.

\Cref{prop:markov-pairs} compares two Markov chains: if only the generator is Markov, $P_2=Q_2$
says that the real process has the generator's pair law, nothing more. The models of this paper
are Markov only in an augmented state. GBM, Merton and the i.i.d.\ Student-$t$ laws of E2 have
i.i.d.\ returns, and a window density is a product of daily densities. In the \GARCHt{} model,
$\sigma_{t+1}^2=\omega/(1-b)+a\sum_{j\ge0}b^j\varepsilon_{t-j}^2$ is a function of past returns
(the model is observation-driven), so the likelihood is explicit given $\sigma_1^2$; yet the
returns are not Markov of any finite order when $a,b>0$ (sketch in \cref{app:proofs-windows}).
The conditional variance has persistence $a+b$: shocks to $\sigma_t^2$ decay like $(a+b)^h$, and
so does the autocorrelation of the squared shocks $\varepsilon_t^2$ when $\E\varepsilon_t^4<\infty$
\citep{francq_2019_garch}. With
standardised Student-$t_\nu$ innovations that requires $(a+b)^2+2a^2(\nu-1)/(\nu-4)<1$, which fails
for the E2 parameters ($a=0.08$, $b=0.90$, $\nu=5$), so for them only the statement about
$\sigma_t^2$ applies. At persistence $0.98$ the volatility half-life is
$\log0.5/\log0.98\approx34$ days, beyond $L=20$.
In the Heston model the variance is a latent continuous state with its own shock and no closed-form
filter, so the likelihood is intractable \citep{aitsahalia2007maximum}; we simulate it with the
quadratic-exponential scheme of \citet{andersen_2008_simple} (\cref{sec:data-generators}). The
stationary bootstrap has no likelihood. Likelihood benchmarks thus exist for GBM, Merton and
\GARCHt{} only, and the critic is the only formal tool of the paper for the other two generators.

\begin{remark}[Pooling tickers]\label{rem:pooling}
Real and synthetic returns of ticker $i$ are divided by the same constant $s_i$ computed on
\TRAIN. Given \TRAIN{} this is a fixed bijection, so the per-ticker null $P_i=Q_i$ is unchanged.
If ticker $i$ contributes a fraction $\pi_i$ of the real windows and $\pi'_i$ of the synthetic
ones, the pooled laws are $\bar P=\sum_i\pi_iP_i$ and $\bar Q=\sum_i\pi'_iQ_i$, and the per-ticker
nulls imply the pooled null $\bar P=\bar Q$ only if $\pi_i=\pi'_i$. Each ticker therefore
contributes the same number of real and synthetic windows. The pooled null is weaker than the
intersection of the per-ticker nulls, because per-ticker errors can cancel in the mixture; the
per-ticker tests of \cref{rem:bh} address this.
\end{remark}

% ---------------------------------------------------------------------------------------------
\subsection{The held-out critic statistic}\label{sec:theory-critic}

Each sample is split into a training and an evaluation half, with index sets $\Itr,\Iev$ (real,
sizes $n_t,n_e$) and $\Jtr,\Jev$ (synthetic, sizes $m_t,m_e$). Let $\Gsig$ be the $\sigma$-field
generated by the training halves, the index sets and the training randomness $\xi$
(initialisation, mini-batches). The critic is $\fhat=\ghat\circ A$, where
$A(x)=\Lambda^{-1}(x-m_A)$ is the coordinate-wise standardisation fitted on the pooled training
half (per asset channel for multivariate windows) and $\ghat$ is the network trained by
\eqref{eq:wgan-gp}. Conditioning on $\Gsig$ fixes $\fhat$, and the statistic is
$T=\Pev\fhat-\Qev\fhat$ \eqref{eq:critic-stat}. The networks are compositions of globally
Lipschitz maps (affine and convolutional layers, LeakyReLU, LayerNorm with a positive
$\epsilon$), so $\Lip(\ghat)<\infty$ and $\fhat\in L^p(P)$ whenever $P\in\Pspace{p}(\R^d)$, although
the constant is not controlled.

\begin{fact}[Kantorovich--Rubinstein duality]\label{fact:kr-duality}
\citep[Ch.~5--6]{villani2009optimal}. For $P,Q\in\Pspace{1}(\R^d)$,
$\W(P,Q):=\inf_{\pi\in\Pi(P,Q)}\int\norm{x-y}\,\pi(dx,dy)=\sup_{f\in\Lip_1}(Pf-Qf)$, where
$\Pi(P,Q)$ is the set of couplings and $\Lip_c=\{f:\Lip(f)\le c\}$. $\W$ is a metric on
$\Pspace{1}$, and it metrises weak convergence together with convergence of first moments.
\end{fact}

$\W$ is the integral probability metric of $\Lip_1$ \citep{muller1997integral}; on $\Pspace{1}$,
$H_0$ holds if and only if $\W(P,Q)=0$.

\begin{remark}[The critic sees standardised windows]\label{rem:representation}
Given $\Gsig$, $A$ is a fixed affine bijection, so $A_\#P=A_\#Q$ if and only if $P=Q$. The
gradient penalty differentiates $\ghat$ with respect to the standardised input, so the Lipschitz
constant it pushes towards one is $\Lip(\ghat)$, the Lipschitz constant of $\fhat$ for the metric
$\norm{x-y}_A=\norm{\Lambda^{-1}(x-y)}$. Write $P_A=A_\#P$ and $Q_A=A_\#Q$. Then $\W(P_A,Q_A)$ is
the Wasserstein-1 distance of $P$ and $Q$ under $\norm{\cdot}_A$, and
$\norm{\Lambda}_{\mathrm{op}}^{-1}\W(P,Q)\le\W(P_A,Q_A)\le\norm{\Lambda^{-1}}_{\mathrm{op}}\W(P,Q)$.
In E2 the windows are raw daily returns with volatility $0.02$, so $\norm{\Lambda^{-1}}_{\mathrm{op}}$
is about $50$; in E3 and E4 the returns are already divided by $s_i$ and $\Lambda$ is of order one.
\end{remark}

\begin{proposition}[Conditional unbiasedness and a bound by $\W$ of the standardised laws]
\label{prop:unbiased}
Assume that, conditionally on $\Gsig$, every evaluation window $X_i$ has law $P$, every
evaluation window $Y_j$ has law $Q$, and $P,Q\in\Pspace{1}(\R^d)$.
\begin{enumerate}
  \item[(a)] $\E[T\mid\Gsig]=P\fhat-Q\fhat=P_A\ghat-Q_A\ghat=:\Dcrit(\fhat)$.
  \item[(b)] $\Dcrit(\fhat)\le\Lip_{S_A}(\ghat)\,\W(P_A,Q_A)\le\Lip(\ghat)\,\W(P_A,Q_A)$, where
  $\Lip_{S_A}$ is the Lipschitz constant on $S_A=\supp P_A\cup\supp Q_A$.
  \item[(c)] If the evaluation windows are also conditionally independent and
  $P,Q\in\Pspace{2}(\R^d)$, then
  $\Var(T\mid\Gsig)=\Var_P\fhat/n_e+\Var_Q\fhat/m_e\le\Lip(\ghat)^2\big(\Tr\Sigma_{P_A}/n_e+\Tr\Sigma_{Q_A}/m_e\big)$,
  where $\Sigma_{P_A}=\Lambda^{-1}\Sigma_P\Lambda^{-1}$ is the covariance matrix of $P_A$.
\end{enumerate}
\end{proposition}
\noindent\textit{Proof.} In \cref{app:proofs-critic}.

\begin{remark}[What the $\W$ reading of $T$ does and does not say]\label{rem:gp-soft}
If $\Lip(\ghat)\le1$, part (b) makes the conditional mean of $T$ a lower bound on $\W(P_A,Q_A)$;
this is the sense in which $T$ is ``a $\W$ lower bound''. The reading has five limits.
(i) \emph{Scale:} it concerns standardised windows, not return units (\cref{rem:representation}).
(ii) \emph{Sampling error:} $T$ itself can exceed $\W(P_A,Q_A)$ by a term of order
$n_e^{-1/2}+m_e^{-1/2}$; the conservative summary is the lower confidence bound
$T-z_{1-\alpha}(s_X^2/n_e+s_Y^2/m_e)^{1/2}$.
(iii) \emph{Soft constraint:} the penalty enforces the Lipschitz constraint only in expectation
over interpolates of training points, and it is two-sided, pushing gradient norms towards one
rather than below one \citep{petzka_2018_regularization}. Nothing guarantees $\Lip(\ghat)\le1$,
and WGAN critics are known to be poor estimators of $\W$ itself
\citep{mallasto2019well,stanczuk_2021_wasserstein}.
(iv) \emph{Looseness:} a small $T$ says that the trained critic found no separating direction, not
that $\W$ is small; \cref{prop:consistency} says what training would have to deliver.
(v) \emph{Relevance:} $\W$ of standardised windows weights directions by the inverse coordinate
scale and ignores what a user of the data values; \cref{sec:theory-markowitz} shows that it can
rank generators against their portfolio utility.
We therefore use $T$ as a test statistic. Its level (\cref{sec:theory-level}) does not depend on
any Lipschitz property, and its $\W$ reading is approximate and descriptive.
\end{remark}

\begin{remark}[Why hold data out]\label{rem:holdout}
The plug-in distance $\W(P_n,Q_m)$ is biased upwards, since $\E\,\W(P_n,P)$ decays at the rate
$n^{-1/d}$ for $d\ge3$ \citep{fournier2015rate}, with a constant that depends on $d$ and on $P$. With
one window the bias is $\E\,\W(\delta_{X_1},P)=\E\norm{X_1-X_1'}$, $X_1'$ an independent copy, which
is of the order of the scale of the windows. The rate factor is still $0.68$ at $d=20$ and
$n=2000$, so at this rate going from one window to 2000 removes only about a third of that bias.
The in-sample training objective is biased upwards in the same way. The
held-out $T$ is conditionally unbiased for $\Dcrit(\fhat)$ and has an exact null distribution
(\cref{prop:exact-level}), at the cost of sample splitting, which \cref{prop:local-power}
quantifies. This is the template of sample-split two-sample tests
\citep{lopezpaz2017revisiting,kim2021classification,liu2020learning,kubler2022witness}; for
Wasserstein-based testing see \citet{ramdas2017wasserstein}. The critic is a learned, sample-split
relative of the test functions of \citet{bierens1990consistent} and
\citet{stinchcombe1998consistent}.
\end{remark}

% ---------------------------------------------------------------------------------------------
\subsection{Exact level under exchangeability}\label{sec:theory-level}

Pool the evaluation windows as $V=(V_1,\dots,V_{n_e+m_e})=(X_{\Iev},Y_{\Jev})$, real first, and let
$t(v)=n_e^{-1}\sum_{k\le n_e}\fhat(v_k)-m_e^{-1}\sum_{k>n_e}\fhat(v_k)$. Then $T=t(V)$, the
permuted statistics are $T_b=t(\pi_b\cdot V)$ with $(\pi\cdot v)_k=v_{\pi^{-1}(k)}$, and the
p-value \eqref{eq:perm-pvalue} is $\pperm=(1+\#\{b\ge1:T_b\ge T\})/(1+B)$.

\begin{assumption}[Exchangeability]\label{ass:exch}
\emph{(A1)} (i) $X_1,\dots,X_n$ are i.i.d.\ $P$, $Y_1,\dots,Y_m$ are i.i.d.\ $Q$, and the two
samples are independent. (ii) The index sets are chosen independently of the window values;
random, deterministic, group-wise and time-block splits qualify. (iii) $\fhat$ is a measurable
function of the training windows and of randomness $\xi$ independent of the data. (iv) The
permutations $\pi_1,\dots,\pi_B$ are i.i.d.\ uniform on the permutations of $\{1,\dots,n_e+m_e\}$,
independent of everything else.\\
\emph{(A1$'$)} Conditionally on $\Gsig$, the law of $V$ is invariant under every permutation of
its coordinates. Under $H_0$, (A1) implies (A1$'$).
\end{assumption}

\begin{proposition}[Exact level]\label{prop:exact-level}
Under (A1$'$) and (A1)(iv), in particular under (A1) and $H_0$, for every $\alpha\in(0,1)$, every
$B\ge1$, every $n_e,m_e\ge1$ and every critic, however it was trained,
\[
  \Prob(\pperm\le\alpha\mid\Gsig)\le\frac{\lfloor\alpha(B+1)\rfloor}{B+1}\le\alpha\quad\text{a.s.},
  \qquad\text{hence}\qquad \Prob(\pperm\le\alpha)\le\alpha .
\]
\end{proposition}
\noindent\textit{Proof.} In \cref{app:proofs-level} (\cref{lem:exchangeable,lem:rank}).

This is the randomisation principle \citep{lehmann2022testing} in the random-permutation form of
\citet{phipson2010permutation} and \citet{hemerik2018exact}. The guarantee needs no Lipschitz
constraint, no convergence of training and no continuity: a poor or overfitted critic costs power,
never level. For a discrete statistic such as the C2ST accuracy, randomised tie-breaking makes the
p-value exactly uniform under (A1$'$). The guarantee does need conditional exchangeability of the
evaluation windows given the training data, which fails when real windows depend on each other or
on the training windows (\cref{sec:theory-dependence}).

\begin{remark}[Weaker nulls and studentisation]\label{rem:studentised}
(a) Conditionally on $\Gsig$, $z=T/(s_X^2/n_e+s_Y^2/m_e)^{1/2}$ is asymptotically $N(0,1)$
whenever $P\fhat=Q\fhat$ and second moments are finite (Lindeberg); under $H_0$ its normal p-value
is asymptotically equivalent to $\pperm$ \citep{hoeffding_1952_large}.
(b) The permutation test is exact for $P=Q$. For the weaker hypothesis $H_0':P\fhat=Q\fhat$ the
unstudentised permutation test can be miscalibrated when $n_e\neq m_e$ and the two variances
differ \citep{romano_1990_behavior}; with $n_e=m_e$ the permutation variance matches the true
variance asymptotically. Permuting the studentised statistic is asymptotically valid for $H_0'$
and stays exact under $H_0$ \citep{janssen1997studentized,chung2013exact}. These results concern
unequal i.i.d.\ laws, not serial dependence. They matter for the relative tests, whose nulls are
statements about means of a trained critic.
\end{remark}

\begin{remark}[Seeds and the p-value lottery]\label{rem:combine-seeds}
The p-value depends on the random split, on $\xi$ and, on real data, on the synthetic seed.
Reporting the smallest p-value over $S$ seeds is invalid. If every $p_s$ is valid, that is
$\Prob(p_s\le u)\le u$ for all $u\in(0,1)$, then $\pmed=\min\{1,2\med_sp_s\}$ \eqref{eq:pmed} is
valid under arbitrary dependence between the $p_s$ \citep{meinshausen_2009_pvalues}: the event
$\{\pmed\le\alpha\}$ implies that at least $S/2$ of the $p_s$ are at most $\alpha/2$, and Markov's
inequality bounds its probability by $(2/S)\sum_s\Prob(p_s\le\alpha/2)\le\alpha$. Twice the mean
is valid as well \citep{vovk_2020_combining}. The combination inherits any miscalibration of the
individual $p_s$.
\end{remark}

% ---------------------------------------------------------------------------------------------
\subsection{Consistency}\label{sec:theory-consistency}

The standardisation $A$ is estimated on the training half. Under the conditions of
\cref{ass:consistency} below, if moreover the real share $n_t/(n_t+m_t)$ of the training windows
converges and the limiting pooled law has positive coordinate variances, $A$ converges in
probability (almost surely for nested training halves) to a population standardisation $A^*$, and
the estimation error adds $o_\Prob(1)$ terms to every step (\cref{app:proofs-consistency}).
We therefore state the results for the standardised laws $A^*_\#P$ and $A^*_\#Q$ and for the
network $\ghat$, with Lipschitz constants in standardised coordinates, and, in this subsection
only, write $P,Q$ for $A^*_\#P,A^*_\#Q$ and $\fhat$ for $\ghat$.

\begin{assumption}[Consistency regime]\label{ass:consistency}
\emph{(C1)} (A1) holds, $P,Q\in\Pspace{2}(\R^d)$ and $P\neq Q$.
\emph{(C2)} $n_t,m_t,n_e,m_e\to\infty$ and $B\ge\lceil2/\alpha\rceil$ ($\pperm\le\alpha$ requires
only $B+1\ge1/\alpha$; the stronger condition is used in the proof).
\emph{(C3)} (Hard Lipschitz constraint, an idealisation of the gradient penalty.) There are
$c<\infty$ and classes $\mathcal F_n\subset\Lip_c$ with $\Prob(\fhat\in\mathcal F_n)\to1$.
\emph{(C4)} (Richness.) There are $n_0$ and $f_0\in\bigcap_{n\ge n_0}\mathcal F_n$ with
$Pf_0-Qf_0=:\eta_0>0$.
\emph{(C5)} (Approximate optimisation.) With probability tending to one,
$(\Ptr-\Qtr)\fhat\ge\sup_{f\in\mathcal F_n}(\Ptr-\Qtr)f-\varepsilon_{\mathrm{opt}}$ with
$\varepsilon_{\mathrm{opt}}\le\eta_0/2$.
\end{assumption}

\begin{proposition}[Consistency of the critic test]\label{prop:consistency}
Under \cref{ass:consistency}, $\Prob(\pperm\le\alpha)\to1$.
\end{proposition}
\noindent\textit{Proof.} In \cref{app:proofs-consistency} (uses \cref{lem:uniform-lip}).

\begin{corollary}[$T$ estimates $\W$ of the standardised laws under idealised training]
\label{cor:w1-limit}
Suppose (C1)--(C3) hold with $c=1$, $\sup_{f\in\mathcal F_n}(Pf-Qf)\to\W(P,Q)$, and (C5) holds
with $\varepsilon_{\mathrm{opt}}\to0$ in probability. Then $T\pto\W(A^*_\#P,A^*_\#Q)$, the
Wasserstein-1 distance under the metric $\norm{\Lambda_*^{-1}(x-y)}$, not in return units.
\end{corollary}
\noindent\textit{Proof.} In \cref{app:proofs-consistency}.

\begin{remark}[How realistic are (C3)--(C5)?]\label{rem:c-realism}
(C4) is a population property. It holds if, for some $c$, the $c$-Lipschitz members of
$\bigcup_n\mathcal F_n$ are dense in the bounded 1-Lipschitz functions for uniform convergence on
compacts, because bounded Lipschitz functions separate probability measures
\citep{dudley_2002_real}; norm-constrained GroupSort networks have this density property
\citep{anil_2019_sorting}. For the LeakyReLU and LayerNorm networks used here, trained with a soft
penalty, we do not prove it. (C3) idealises the penalty (\cref{rem:gp-soft}), and (C5) is an
assumption about stochastic-gradient training that we cannot verify. \Cref{prop:consistency} says
what has to hold for consistency, not that WGAN-GP training delivers it; the power study of
\cref{sec:sim-power} measures what training delivers.
\end{remark}

% ---------------------------------------------------------------------------------------------
\subsection{Likelihood connection and the Neyman--Pearson bound}\label{sec:theory-likelihood}

\begin{proposition}[Bayes-optimal logit; \citealp{goodfellow2014generative}]\label{prop:bayes-logit}
Let $\nu=P+Q$, $p=dP/d\nu$, $q=dQ/d\nu$, and let $(V,\iota)$ satisfy
$\iota\sim\mathrm{Bernoulli}(\pi)$, $V\mid\iota=1\sim P$ and $V\mid\iota=0\sim Q$. Among measurable
$D:\R^d\to[0,1]$, the cross-entropy risk $-\E[\iota\log D(V)+(1-\iota)\log(1-D(V))]$ is minimised
by $D^*=\pi p/(\pi p+(1-\pi)q)$, $\nu$-a.e., whose logit is
$\lstar(v)=\log\{p(v)/q(v)\}+\log\{\pi/(1-\pi)\}$.
\end{proposition}
\begin{proof}
The risk equals $\int-[\pi p\log D+(1-\pi)q\log(1-D)]\,d\nu$. Where $p,q>0$ the integrand is
strictly convex in $D\in[0,1]$ and minimised at $D^*$; where $q=0<p$ the minimiser is $D=1$, and
conversely.
\end{proof}

\begin{corollary}[Targets of the logit and accuracy statistics]\label{cor:jeffreys}
(a) If $\lhat=\lstar+\text{const}$ and $\KL(P\Vert Q)$, $\KL(Q\Vert P)$ are finite, then the
\LRC{} statistic $\Tlogit=\Pev\lhat-\Qev\lhat$ satisfies
$\E[\Tlogit\mid\Gsig]=\KL(P\Vert Q)+\KL(Q\Vert P)=\Jeff(P,Q)$, the Jeffreys divergence; the class
prior cancels. (b) For a decision region $\hat A=\{\lhat>0\}$ the balanced accuracy satisfies
$\E[\Tacc\mid\Gsig]=\tfrac12+\tfrac12\{P(\hat A)-Q(\hat A)\}\le\tfrac12\{1+\TV(P,Q)\}$, with
equality for $\hat A=\{p>q\}$.
\end{corollary}
\begin{proof}
Apply \cref{prop:unbiased}(a), which needs no Lipschitz property, to $f=\lstar$ and to
$f=\one_{\hat A}$.
\end{proof}

Unlike the \WC{}, nothing makes $\E[\Tlogit\mid\Gsig]$ a lower bound on $\Jeff$ for an imperfect
logit: rescaling $\lhat$ inflates it. $\Tlogit$ is a plug-in estimator, consistent only if
$\lhat\to\lstar$ in $L^1(P)\cap L^1(Q)$; the level of the test does not depend on this. The
\LRC{} is thus a learned log-likelihood-ratio statistic \citep{cheng2022classification}, and
\cref{sec:sim-llr} measures how close the learned logit is to $\log(p/q)$ (\cref{tab:e2-llr}).

\begin{proposition}[No two-sample test beats the Neyman--Pearson oracle]\label{prop:oracle-bound}
Let $P,Q$ have densities $p,q$ with respect to a common $\sigma$-finite measure. Write
$\bNP(F\to F';m,\alpha)$ for the power of the most powerful level-$\alpha$ test of
$F^{\otimes m}$ against $F'^{\otimes m}$ \citep{lehmann2022testing}. Let
$\phi(X_{1:n},Y_{1:m},U)\in[0,1]$ be any test whose auxiliary randomness $U$ (splits, training,
permutations) is independent of the data.
\begin{enumerate}
  \item[(a)] If $\E_{P^{\otimes n}\otimes P^{\otimes m}}\phi\le\alpha$, then
  $\E_{P^{\otimes n}\otimes Q^{\otimes m}}\phi\le\bNP(P\to Q;m,\alpha)$.
  \item[(b)] If $\E_{Q^{\otimes n}\otimes Q^{\otimes m}}\phi\le\alpha$, then
  $\E_{P^{\otimes n}\otimes Q^{\otimes m}}\phi\le\bNP(Q\to P;n,\alpha)$.
\end{enumerate}
\end{proposition}
\noindent\textit{Proof.} In \cref{app:proofs-likelihood}.

A two-sample test of exact level $\alpha$ thus has power at most
$\min\{\bNP(P\to Q;m,\alpha),\bNP(Q\to P;n,\alpha)\}$: its level must hold at every null point
$P=Q$, while the oracle knows both densities. The E2 oracle \eqref{eq:lr-oracle} applies the
Neyman--Pearson statistic $\TLR=\sum_j\log\{q(Y_j)/p(Y_j)\}$ to the synthetic sample with a Monte
Carlo critical value. As a Monte Carlo test it has exact level, and its power tends to $\bNP$ as
the number of null samples grows \citep{dufour_2006_monte}. For the i.i.d.\ families (a), (b) and
(d) of E2 the window density is a product of daily densities and the oracle is exact. For the
\GARCHt{} family (c) the window density
$p(x_{1:L})=\int\prod_tp(x_t\mid\sigma_t^2)\,\Pi(d\sigma_1^2)$, with $\Pi$ the law of the
conditional variance at the window start, has no closed form. E2 replaces it by an average over
$M=512$ draws from $\Pi$, shared by the observed and the null samples. The Monte Carlo test keeps
its exact level, but its statistic only approximates the log-likelihood ratio (the logarithm of an
unbiased density estimate is biased downwards by a term of order $1/M$). Its power need not equal
$\bNP$, so the family-(c) curve is an approximate envelope, and a learned test above it would not
contradict \cref{prop:oracle-bound}.

\begin{proposition}[Local power of fixed-critic tests]\label{prop:local-power}
Let $\{Q_\theta\}_{\theta\in[0,\theta_0)}$, with $Q_0=P$, be differentiable in quadratic mean at
$0$ with score $h\in L^2(P)$ and $I=Ph^2>0$. Let $m\to\infty$ with $n/m\to\kappa\in(0,\infty)$,
$n_e=\lfloor\rho n\rfloor$ and $m_e=\lfloor\rho m\rfloor$ for some $\rho\in(0,1]$, and
$\theta_m=\tau/\sqrt m$ with $\tau>0$.
\begin{enumerate}
  \item[(a)] $\bNP(P\to Q_{\theta_m};m,\alpha)\to1-\Phi(z_{1-\alpha}-\tau\sqrt I)$.
  \item[(b)] Let $f\in L^2(P)$ be fixed, with $\sigma_f^2=\Var_Pf>0$ and orientation
  $r_f=-\Corr_P(f,h)\ge0$ (the critic scores real windows higher). The test that rejects when
  $z>z_{1-\alpha}$, and the permutation test with the same $f$ and $B\to\infty$, have power
  under
  $P^{\otimes n}\otimes Q_{\theta_m}^{\otimes m}$ tending to
  $1-\Phi\big(z_{1-\alpha}-\tau\sqrt I\,r_f\sqrt{\rho\kappa/(1+\kappa)}\big)$.
  \item[(c)] At $\kappa=1$ and $\rho=1/2$ (the defaults), the noncentrality of every test based on
  the held-out mean difference of a fixed critic (studentised or permutation) is at most half the
  oracle's. Where the oracle has power $0.80$ at $\alpha=0.05$, every such test has asymptotic
  power at most $0.34$.
\end{enumerate}
\end{proposition}
\noindent\textit{Proof.} In \cref{app:proofs-likelihood}; proved for a fixed critic, sketched for
a learned one.

The factor $\sqrt{\rho\kappa/(1+\kappa)}$ separates the two costs of not knowing the densities:
the two-sample cost $\sqrt{\kappa/(1+\kappa)}$, which tends to one only when $P$ is effectively
known, and the splitting cost $\sqrt\rho$. The factor $r_f\le1$ is the price of a sub-optimal
direction. The bound concerns the mean-difference statistic $T$ that the protocol uses; other
statistics of the held-out critic values, such as rank or quantile comparisons, are not covered.
With a learned critic the argument applies conditionally on the training half with $f=\fhat$, so
the bound in (c) holds for mean-difference tests with learned critics too; a rigorous version
needs a central limit theorem uniform over the critic class, which we do not provide.

\begin{example}[Volatility level]\label{ex:vol-level}
Let $P=N(0,\sigma^2I_L)$ and $Q_\theta=N(0,(1+\theta)^2\sigma^2I_L)$. The score is proportional to
$\norm{x}^2-L\sigma^2$. The critic $\pm\norm{x}$, oriented to score real windows higher, is an
optimal Kantorovich potential for every scale factor, and its local efficiency is
$r_f^2=\Corr(\chi_L,\chi_L^2)^2=0.876$, $0.957$, $0.988$ for $L=1,5,20$. The population
standardisation is a multiple of the identity here, so standardising changes nothing. Against
volatility-level deviations the population \WC{} is nearly as efficient as the likelihood-ratio
direction (details in \cref{app:proofs-likelihood}).
\end{example}

\begin{remark}[\LRC{} versus \WC{} (heuristic)]\label{rem:lr-vs-w1}
By \cref{prop:bayes-logit} the population logit critic is $\log(p/q_\theta)\approx-\theta h$, the
locally optimal direction ($r_f=1$). The population \WC{} is a Kantorovich potential, which is in
general not proportional to $h$ ($r_f<1$). If both were learned well, the \LRC{} would be locally
at least as powerful. The Lipschitz penalty may instead reduce the estimation variance of the
\WC{} in samples of a few thousand windows. Which effect dominates is an empirical question; in
E2 the \WC{} is the more powerful of the two in most cells (\cref{sec:sim-power}).
\end{remark}

\begin{remark}[What the pooled KS test can and cannot see]\label{rem:ks-marginal}
(i) Suppose the returns are i.i.d.\ within windows under both laws, $P=P_1^{\otimes L}$ and
$Q=Q_1^{\otimes L}$, as in E2 families (a) and (b). Then the pair formed by the empirical
distribution of the $nL$ pooled real returns and that of the $mL$ pooled synthetic returns is
sufficient for $(P_1,Q_1)$. For any test $\phi$, the conditional expectation of $\phi$ given this
pair has the same rejection probability at every $(P_1,Q_1)$, so nothing is lost by looking only
at the two pooled marginals. This is what KS does, with no splitting and no learning; the oracle
statistic $\TLR=\sum_{j,t}\log\{q_1(y_{jt})/p_1(y_{jt})\}$ is also a function of the pooled
synthetic marginal. A window critic must learn this additive structure in dimension $L$ from half
of the data. This is the regime in which E2 finds KS far ahead of every window test
(\cref{tab:e2-power}).
(ii) When returns within a window are dependent, the null distribution of KS, derived for $nL$
independent draws, is wrong: in E2 family (c) KS is oversized under $H_0$ (\cref{tab:e2-size}).
(iii) When $P$ and $Q$ have the same one-dimensional marginal and differ only in their dependence
(the E2-dep families), the scaled KS statistic converges in law to the supremum of a Gaussian
process under the alternative as under the null, so its rejection probability converges to a
limit below one: KS is not consistent against dependence-only deviations, whereas the window tests
are consistent under \cref{ass:consistency} if the deviation is visible at lags below $L$
(\cref{prop:window-ident}; \cref{fig:e2dep-power}). Parts (ii) and (iii) follow from the
empirical-process central limit theorem for window-level sums of indicators; we do not write out
the details.
\end{remark}

\begin{remark}[Likelihood ranks by $\KL$, the critic by $\W$]\label{rem:kl-vs-w1}
The expected out-of-sample log-likelihood of a model $G$ is $\E_P\log g=-\mathrm{Ent}(P)-\KL(P\Vert G)$
whenever finite, and for conditional densities of stationary ergodic processes it is the
corresponding relative-entropy rate. Rankings by predictive log-score, and the Vuong and
Amisano--Giacomini tests \citep{vuong1989likelihood,amisano2007comparing}, therefore rank by
$\KL(P\Vert G)$; the critic ranks by (lower bounds on) $\W$ of window laws. Both vanish only at
$G=P$, but they can rank misspecified models differently. Let $P$ be Student-$t_4$ scaled to unit
variance, $G_1=N(0,1)$ (right variance, wrong tails) and $G_2$ the same $t_4$ law scaled by $1.4$
(right tails, wrong scale). Numerical quadrature gives $\KL(P\Vert G_1)\approx0.084$,
$\W(P,G_1)\approx0.139$, $\KL(P\Vert G_2)\approx0.060$ and $\W(P,G_2)=0.4\,\E_P\abs{X}\approx0.283$:
the likelihood prefers $G_2$, $\W$ prefers $G_1$. Disagreement between the two evaluations in
\cref{sec:emp-likelihood} therefore shows that they measure different divergences, not that
either is broken.
\end{remark}

% ---------------------------------------------------------------------------------------------
\subsection{The relative tests}\label{sec:theory-relative}

There are three samples: $\Rnew$, windows $X_1,\dots,X_n$ with law $\Pnew$; the synthetic sample,
$Y_1,\dots,Y_m$ with law $\QG$; and $\Rpast$, windows $Z_1,\dots,Z_k$ with law $\Ppast$. Markets
are non-stationary, and a generator calibrated on the past inherits features of the past, such
as its volatility level, that need not persist; on our data the absolute \WC{} test rejects every
generator (\cref{tab:e3-pooled-tests}), whereas the window tests do not reject the past windows
themselves against the new ones (E1(c), \cref{tab:e1-nonstat}), a comparison for which they have
little power with about 15 \TEST{} evaluation dates (\cref{sec:emp-placebo}). The relative question is whether $G$ is at
least as close to $\Rnew$ as $\Rpast$ is. For a sample $W$ and a part $\mathrm{p}\in\{\trn,\evl\}$
of it, $W^{\mathrm p}f$ denotes the mean of $f$ over that part.

\emph{Anchored test} (primary, \cref{sec:method-relative}). $Y$ and $Z$ are split into training
and evaluation parts of sizes $m_t,m_e$ and $k_t,k_e$; all of $X$ is used for training. One critic
$\fhat$ is trained to separate $X$ (positive class) from the pool $\{Y^{\trn},Z^{\trn}\}$, in which
the two samples carry weight $1/2$ each. The statistic and its conditional target are
\[
  \Danc=Z^{\evl}\fhat-Y^{\evl}\fhat\quad\eqref{eq:anchored-stat},\qquad
  \delta=\Ppast\fhat-\QG\fhat ,
\]
where now $\Gsig$ is generated by $X$, $Y^{\trn}$, $Z^{\trn}$, the index sets and $\xi$. Large
$\Danc$ means that the critic which scores ``$\Rnew$-likeness'' finds the past more $\Rnew$-like
than $G$.

\emph{Two-critic test} (secondary). Critics $\fhat_G$ on $(X^{\trn},Y^{\trn})$ and $\fhat_P$ on
$(X^{\trn},Z^{\trn})$ are trained separately with independent randomness and evaluated on the
shared $X^{\evl}$: $\Dtwo=(X^{\evl}-Y^{\evl})\fhat_G-(X^{\evl}-Z^{\evl})\fhat_P$
\eqref{eq:two-critic-stat}, with targets $D_G=(\Pnew-\QG)\fhat_G$, $D_P=(\Pnew-\Ppast)\fhat_P$ and
$\delta^{\mathrm{2c}}=D_G-D_P$.

\begin{definition}[Relative nulls]\label{def:rel-null}
(a) The anchored null is $H_0^{\mathrm{anc}}:\delta\le0$, that is $\QG\fhat\ge\Ppast\fhat$ for the
trained critic, against $\delta>0$. (b) The two-critic null is
$H_0^{\mathrm{2c}}:\delta^{\mathrm{2c}}\le0$. (c) The idealised $\W$ null is
$H_0^{W}:\W(\QG,\Pnew)\le\W(\Ppast,\Pnew)$. The boundary point $\QG=\Ppast$ belongs to $H_0^W$ and
to the analogous null for any divergence.
\end{definition}

Both conditional nulls depend on the data through the trained critics. \Cref{prop:anchored}
gives exact or asymptotic level for $H_0^{\mathrm{anc}}$; \cref{prop:anchored-target} says what
$\delta$ measures; \cref{prop:paired-boot,prop:boundary-size} treat the two-critic test.

\begin{assumption}[Anchored exchangeability]\label{ass:anc}
(i) Conditionally on $X$, the $k+m$ windows $Z_1,\dots,Z_k,Y_1,\dots,Y_m$ are i.i.d. This holds
when the three samples are independent, each is i.i.d., and $\QG=\Ppast$, whatever $\Pnew$.
(ii) The split of $Y$ and $Z$ depends only on time stamps, group labels and randomness independent
of the window values. (iii) $\fhat$ is a measurable function of $X$, $Y^{\trn}$, $Z^{\trn}$ and of
randomness $\xi$ independent of the data. (iv) The permutations of the $k_e+m_e$ evaluation labels
are i.i.d.\ uniform and independent of everything else.
\end{assumption}

\begin{proposition}[Level of the anchored test]\label{prop:anchored}
\begin{enumerate}
  \item[(a)] \emph{Exact level at the boundary.} Under \cref{ass:anc}, the permutation p-value of
  $\Danc$ (labels of the pooled $Z^{\evl}$ and $Y^{\evl}$ permuted, $\fhat$ fixed) satisfies
  $\Prob(\pperm\le\alpha\mid\Gsig)\le\lfloor\alpha(B+1)\rfloor/(B+1)\le\alpha$ almost surely, for
  every critic, every $B\ge1$ and $k_e,m_e\ge1$, and whatever $\Pnew$.
  \item[(b)] \emph{Weaker null, independent windows.} Suppose that, conditionally on $\Gsig$,
  $Z^{\evl}$ is i.i.d.\ $\Ppast$, $Y^{\evl}$ is i.i.d.\ $\QG$, the two are independent, and
  $\fhat$ has finite positive variances under both laws. As $k_e,m_e\to\infty$ with
  $k_e/(k_e+m_e)\to\kappa\in(0,1)$ and $\fhat$ fixed, the studentised permutation test satisfies
  $\limsup\Prob(\text{reject}\mid\Gsig)\le\alpha$ on $\{\delta\le0\}$. When $k_e=m_e$, so does the
  unstudentised permutation test.
  \item[(c)] \emph{Dependent windows.} Suppose that, conditionally on $\Gsig$, the evaluation
  windows of $Z$ and of $Y$, ordered by date rank, are independent of each other and satisfy
  \cref{ass:mixing} with marginals $\Ppast$ and $\QG$ and $f=\fhat$. Then the Driscoll--Kraay
  statistic $\zhac$ of $\Danc$ \eqref{eq:hac-z} is asymptotically $N(0,1)$ on $\{\delta=0\}$ and
  tends to $-\infty$ in probability on $\{\delta<0\}$, so the test that rejects when
  $\zhac>z_{1-\alpha}$ has asymptotic level $\alpha$ for $H_0^{\mathrm{anc}}$.
\end{enumerate}
\end{proposition}
\noindent\textit{Proof.} In \cref{app:proofs-relative}.

\begin{proposition}[What the anchored target measures]\label{prop:anchored-target}
Lipschitz constants and $\W$ are taken in the standardised coordinates $A$ of the anchored critic,
and $M=\tfrac12(\QG+\Ppast)$ is the law of the pool.
\begin{enumerate}
  \item[(a)] If $\Ppast,\QG\in\Pspace{1}$ and $\Lip(\ghat)\le c$, then
  $\abs{\delta}\le c\,\W(A_\#\Ppast,A_\#\QG)$. In particular $\delta=0$ whenever $\QG=\Ppast$,
  whatever $\Pnew$ and the critic.
  \item[(b)] \emph{(\LRC{}.)} Let $\Pnew,\QG,\Ppast$ be mutually absolutely continuous with
  $\KL(\QG\Vert\Pnew)$ and $\KL(\Ppast\Vert\Pnew)$ finite. If $\fhat=\lstar$ is the Bayes logit of
  the population training problem ($\Pnew$ against $M$ with equal class weights), then
  \begin{equation}\label{eq:anchored-kl}
    \delta=\big[\KL(\QG\Vert\Pnew)-\KL(\Ppast\Vert\Pnew)\big]-\big[\KL(\QG\Vert M)-\KL(\Ppast\Vert M)\big],
  \end{equation}
  and the second bracket lies in $[-\log2,\log2]$. If moreover $d\QG/d\Ppast=1+\varepsilon h$ with
  $h$ bounded and $\Ppast h=0$, the second bracket equals $-\varepsilon^3\Ppast(h^3)/24+O(\varepsilon^4)$
  as $\varepsilon\to0$.
  \item[(c)] \emph{(\WC{}.)} If $\fhat=f^*$ is a 1-Lipschitz maximiser of $(\Pnew-M)f$, then
  $\delta=(\Pnew-\QG)f^*-(\Pnew-\Ppast)f^*$, and the two terms sum to $2\,\W(A_\#\Pnew,A_\#M)$.
\end{enumerate}
\end{proposition}
\noindent\textit{Proof.} In \cref{app:proofs-relative}.

\begin{remark}[What the anchored null means away from the boundary]\label{rem:anchored-null}
By \cref{prop:anchored-target}(a) the anchored test is a two-sample test of $G$ against the past
whose witness has been trained to point towards $\Rnew$: it can reject only if $G$ differs from the
past, and only in a direction that separates $\Rnew$ from the pool. If the market were stationary
($\Pnew=\Ppast$) it would be a critic test of $\QG=\Ppast$. By part (b), near the boundary the
idealised \LRC{} version compares $\KL(\QG\Vert\Pnew)$ with $\KL(\Ppast\Vert\Pnew)$ up to a
third-order term, so to first order $H_0^{\mathrm{anc}}$ says that $G$ is at least as close to the
new law as the past is, in the Kullback--Leibler divergence from the candidate to the new law; this
is the reverse of the direction $\KL(\Pnew\Vert G)$ by which predictive log-scores rank
(\cref{rem:kl-vs-w1}). Away from the boundary the correction in \eqref{eq:anchored-kl} is bounded
but not small. By part (c), the \WC{} version compares the two absolute discrepancies with a common
witness. Each is at most the corresponding $\W$ distance, but a common witness need not be optimal
for either pair, so $\delta>0$ is possible under $H_0^W$ and conversely: neither relative test is a
test of $H_0^W$. A rejection (\BAD) means that $G$ is detectably less $\Rnew$-like than the past
for a critic trained to recognise $\Rnew$ against an even mixture of both. A non-rejection is not
evidence that $G$ is good, which is why the verdict needs a power check (\cref{sec:theory-verdict}).
Two idealisations remain. First, $G$ is calibrated on $\Rpast$, so its windows depend on the past
windows through the fitted parameters and \cref{ass:anc}(i) holds only approximately; for a
parametric generator each past window moves the estimate by $O(1/k)$, which we expect to be
negligible but do not prove. Second, the stationary bootstrap copies past windows that the critic
saw as pool negatives; evaluation windows of $G$ then duplicate training points, exchangeability
fails in the direction of rejection, and the generator receives the verdict \INVALID{} (E2 part e
documents this leakage, \cref{tab:e2-relative-boundary}).
\end{remark}

\begin{remark}[Non-stationary past (heuristic)]\label{rem:nonstat-past}
On real data $\Rpast$ is the \TRAIN{} panel, whose windows are not identically distributed (the
2020 crash, the 2022 drawdown). \Cref{ass:anc}(i) then fails, and the natural boundary is
$\QG$ equal to the pooled \TRAIN{} window law. The time-block split sends a random half of the
\TRAIN{} blocks to the evaluation part, so $Z^{\evl}\fhat$ estimates the critic mean over a random
half of the \TRAIN{} regimes while $\QG\fhat$ averages over all of them. The target $\delta$ is then
a random variable, not necessarily centred at zero, whose spread reflects the heterogeneity between
blocks, which the Driscoll--Kraay variance of the evaluation part does not capture. By the argument of
\cref{prop:boundary-size}(a), the rejection rate at this boundary is about
$\E[1-\Phi(z_{1-\alpha}-\delta/\sigma)]$, which exceeds $\alpha$ when $\delta$ is roughly symmetric
and not small relative to the evaluation standard deviation $\sigma$. In addition, the last
\TRAIN{} block borders the first \TEST{} dates; with volatility clustering, when that block falls
in the evaluation part it raises $Z^{\evl}\fhat$, in the direction of rejection. We therefore claim
no calibration for a single-seed anchored p-value on real data. The verdict uses $\pmed$ over the
synthetic seeds, and the size of this rule at the pooled-\TRAIN{} boundary and its power are
estimated in design D2, which keeps the real \TRAIN{} volatility path, as well as in the
stationary design D1 (\cref{def:power-check}, \cref{tab:e3-power}). The estimate depends on the
periods: at the main split the boundary size of the rule in D2 is 0.08 [0.04, 0.15], at the
alternative split of E5 it is 0.60 [0.39, 0.79] (Clopper--Pearson 95\% intervals;
\cref{tab:manual-e5-split}). \Cref{rem:calibrated-bad} states when the boundary replications
calibrate an observed $\pmed$, and separates the power of the rule from its miscalibration.
\end{remark}

For the two-critic test we use the regime of \citet{bounliphone2016test}-type comparisons with
learned witnesses.

\begin{assumption}[Two-critic regime]\label{ass:rel}
(R1) Conditionally on $\Gsig$ (the training halves of the three samples and the training
randomness of both critics), the three evaluation samples are mutually independent and i.i.d.\
from $\Pnew$, $\QG$ and $\Ppast$. (R2) $n_e,m_e,k_e\to\infty$ proportionally. With
$\sigma_N^2=\sigma_X^2/n_e+\sigma_Y^2/m_e+\sigma_Z^2/k_e$, $\sigma_X^2=\Var_{\Pnew}(\fhat_G-\fhat_P)$,
$\sigma_Y^2=\Var_{\QG}\fhat_G$ and $\sigma_Z^2=\Var_{\Ppast}\fhat_P$, the quantity
$(n_e+m_e+k_e)\sigma_N^2$ is bounded away from zero and the $(2+\epsilon)$-th moments of the
three summands are bounded in probability. (R3) The bootstrap resamples $X^{\evl}$ once per
replicate for both critics (paired), and $Y^{\evl}$, $Z^{\evl}$ independently.
\end{assumption}

\begin{proposition}[Paired bootstrap, conditional validity; sketch]\label{prop:paired-boot}
Under \cref{ass:rel}, conditionally on $\Gsig$: (a) $(\Dtwo-\delta^{\mathrm{2c}})/\sigma_N\dto N(0,1)$;
(b) the bootstrap law of $\Dtwo{}^*-\Dtwo$ estimates the law of $\Dtwo-\delta^{\mathrm{2c}}$
uniformly consistently; (c) the test that rejects when
$p^*=(1+\#\{b:\Dtwo_b{}^*-\Dtwo\ge\Dtwo\})/(1+B)\le\alpha$ has
$\Prob(\text{reject}\mid\Gsig)=1-\Phi(z_{1-\alpha}-\delta^{\mathrm{2c}}/\sigma_N)+o_\Prob(1)$ as
$B\to\infty$, which is at most $\alpha+o_\Prob(1)$ on $\{\delta^{\mathrm{2c}}\le0\}$; (d) the
percentile interval has asymptotic coverage $0.95$ for $\delta^{\mathrm{2c}}$.
\end{proposition}
\noindent\textit{Proof.} Sketch in \cref{app:proofs-relative}. The centred p-value follows
\citet{hall1991two}.

\begin{proposition}[The two-critic test is anti-conservative at the boundary]
\label{prop:boundary-size}
Assume \cref{ass:rel} along a sequence of experiments, let $B\to\infty$, $\alpha<1/2$ and
$z=z_{1-\alpha}$.
\begin{enumerate}
  \item[(a)] $\Prob(\text{reject})=\E[1-\Phi(z-\delta^{\mathrm{2c}}/\sigma_N)]+o(1)$.
  \item[(b)] Suppose in addition that $\QG=\Ppast$, $m_t=k_t$ and $m_e=k_e$; that
  $\fhat_G=\Psi(X^{\trn},Y^{\trn},\xi_G)$ and $\fhat_P=\Psi(X^{\trn},Z^{\trn},\xi_P)$ for one
  training map $\Psi$ with $\xi_G,\xi_P$ i.i.d.\ and independent of the data; and that the joint
  law of the data and the index sets is invariant under exchanging $(Y,\xi_G)$ and $(Z,\xi_P)$.
  Then $(\delta^{\mathrm{2c}},\sigma_N)\eqd(-\delta^{\mathrm{2c}},\sigma_N)$ and
  \[
    \Prob(\text{reject})=\E\,\bar g\big(\abs{\delta^{\mathrm{2c}}}/\sigma_N\big)+o(1),\qquad
    \bar g(x)=1-\tfrac12\{\Phi(z-x)+\Phi(z+x)\}.
  \]
  $\bar g$ is continuous and strictly increasing on $[0,\infty)$, with $\bar g(0)=\alpha$ and
  $\bar g(\infty)=1/2$. Hence $\Prob(\text{reject})\ge\alpha+o(1)$, and if
  $\Prob(\abs{\delta^{\mathrm{2c}}}>\varepsilon\sigma_N)\ge\eta$ for some $\varepsilon,\eta>0$ and
  all large samples, then $\liminf\Prob(\text{reject})\ge\alpha+\eta\{\bar g(\varepsilon)-\alpha\}>\alpha$.
  \item[(c)] If the training halves are held fixed while the evaluation sizes grow, then
  $\Prob(\text{reject})\to\Prob(\delta^{\mathrm{2c}}>0)+\alpha\Prob(\delta^{\mathrm{2c}}=0)$. Under
  the conditions of (b) this limit equals $1/2$ when $\delta^{\mathrm{2c}}$ has no atom at zero.
\end{enumerate}
\end{proposition}
\noindent\textit{Proof.} In \cref{app:proofs-relative}.

At the boundary, $\delta^{\mathrm{2c}}$ is the difference of two exchangeable replicates of a
training-dependent quantity. It fluctuates at the scale of critic training, and nothing in
WGAN-GP training makes it $o_\Prob(\sigma_N)$. The proof uses only the conditional rejection
probability and the symmetry, so the conclusion also holds for the HAC version of the two-critic
test. The anchored test avoids the problem: with one critic the target is exactly zero at the
boundary (\cref{prop:anchored-target}(a)) and the permutation is exact (\cref{prop:anchored}(a)).
E2 part e compares the two tests at the boundary (\cref{sec:sim-relative}). With $n=500$ windows
per sample the two-critic test rejects at 0.040 and 0.065 at nominal 5\% in its two boundary
designs, and the anchored test at 0.065 and 0.055 (\cref{tab:e2-relative-boundary}): in that
design the spread of $\delta^{\mathrm{2c}}$ induced by training is small relative to $\sigma_N$,
and the excess allowed by part (b) is not visible. The anchored test is preferred because its level
is exact for every critic and every sample size, not because the two-critic test failed there.

\begin{remark}[Guarantee of the two-critic test; the MMD version]\label{rem:rel-guarantee}
Under \cref{ass:rel}, a two-critic rejection means that the critic family, given equal resources,
separates $G$ from $\Rnew$ by more than it separates $\Rpast$ from $\Rnew$. By
\cref{prop:unbiased}, $D_G$ and $D_P$ equal $\Lip(\ghat_G)\W(\QG,\Pnew)$ and
$\Lip(\ghat_P)\W(\Ppast,\Pnew)$ minus unknown non-negative slacks (in standardised coordinates),
so with respect to $H_0^W$ the test is neither exact nor conservative. It is fair only to the
extent that the two critic problems are equally hard (same architecture, budget and sample sizes),
and even then \cref{prop:boundary-size} shows over-rejection at $\QG=\Ppast$. The MMD version of
\citet{bounliphone2016test} uses a fixed kernel and a genuine metric null, but its normal
approximation needs a non-degenerate projection, which fails when all three laws coincide. For
serially dependent evaluation windows the i.i.d.\ bootstrap of $X^{\evl}$ must be replaced by a
block bootstrap \citep{kunsch1989jackknife,politis1994stationary,lahiri2003resampling} or by a HAC
variance, as implemented.
\end{remark}

% ---------------------------------------------------------------------------------------------
\subsection{Serial and cross-sectional dependence}\label{sec:theory-dependence}

The real windows of one ticker form a stationary but dependent sequence: volatility clustering
makes window features such as realised variance persistent. Windows of different tickers on the
same date are dependent through common factors, whereas synthetic windows from univariate
generators are independent across tickers. Under $H_0$ the pooled evaluation vector is then not
exchangeable and \cref{prop:exact-level} does not apply.

\begin{remark}[Size distortion of the naive permutation test (heuristic)]\label{rem:naive-perm}
Let $a_k=\fhat(X_{(k)})$ be the real evaluation values in time order, with variance $\sigma^2$ and
long-run variance $\LRV_a=\sigma^2(1+2\sum_{h\ge1}\rho_h)$, and let the synthetic values be
i.i.d.\ with variance $\sigma^2$, with $n_e=m_e$. The permutation distribution uses the variance
$2\sigma^2/n_e$, while the true variance is $(\LRV_a+\sigma^2)/n_e$, so the rejection rate tends
to $1-\Phi\big(z_{1-\alpha}\{2/(2+2\sum_h\rho_h)\}^{1/2}\big)$. For $\rho_h=\phi^h$ the nominal 5\%
test rejects about 12\% of the time at $\phi=0.5$ and 17\% at $\phi=0.67$. Contemporaneous
covariances between tickers inflate the true variance in the same way.
\end{remark}

For a fixed function $f$, let $a_k=f(X_k)$ be the real evaluation values, each with the date rank
$\tau(k)$ of its first day, and let $\mathcal S_\tau=\sum_{k:\tau(k)=\tau}(a_k-\bar a)$ be the
date sums. The Driscoll--Kraay variance \eqref{eq:dk-var} is the Bartlett-kernel long-run variance
of $(\mathcal S_\tau)$, with lags counted in windows and the truncation lag $H$ chosen by the
\citet{andrews1991heteroskedasticity} AR(1) plug-in; with one window per date it is the
\citet{newey1987simple} estimator. The synthetic values $b_k$ are treated in the same way, and
$\zhac=(\bar a-\bar b)/\{\hat V_{\mathrm{DK}}(\bar a)+\hat V_{\mathrm{DK}}(\bar b)\}^{1/2}$
\eqref{eq:hac-z}.

\begin{assumption}[Weak dependence, fixed critic $f$]\label{ass:mixing}
Each sample has a fixed number of windows per date ($N$ for a panel, one for a single series).
(D1) The date-level sums $u_\tau=\sum_{k:\tau(k)=\tau}f(X_k)$ of the real sample and
$v_\tau=\sum_{k:\tau(k)=\tau}f(Y_k)$ of the synthetic sample form strictly stationary sequences,
independent of each other, that are strongly mixing with coefficients satisfying
$\sum_h\alpha(h)^{\epsilon/(2+\epsilon)}<\infty$, and $u_\tau$, $v_\tau$ have finite
$(2+\epsilon)$-th moments for some $\epsilon>0$ (for Lipschitz $f$ this follows from finite
$(2+\epsilon)$-th moments of returns). (D2) The long-run variances $\LRV_u=\sum_h\Cov(u_0,u_h)$
and $\LRV_v$ are positive. (D3) The Driscoll--Kraay estimators are consistent:
$\hat V_{\mathrm{DK}}(\bar a)/\Var(\bar a)\pto1$, and likewise for $\bar b$; sufficient conditions,
including data-dependent lags, are in \citet{andrews1991heteroskedasticity},
\citet{newey1987simple}, \citet{driscoll1998consistent} and \citet{dejong_2000_consistency}.
(D4) The numbers of evaluation dates of the two samples tend to infinity proportionally.
\end{assumption}

With the time-block split the evaluation dates of a sample form a union of at most $K$ contiguous
blocks; because $K$ is fixed while the blocks grow, (D1) for the underlying date sequence carries
the central limit theorem and (D3) over to the union. The edge terms are of order $KH^2$ divided by
the number of dates, which vanishes when $H$ grows like the cube root of that number, the
\citet{andrews1991heteroskedasticity} rate (\cref{app:proofs-dependence}).

\begin{proposition}[Asymptotic validity of the HAC test for a fixed critic]\label{prop:hac}
Under \cref{ass:mixing} and $H_0':Pf=Qf$, $\zhac\dto N(0,1)$, so the test that rejects when
$\zhac>z_{1-\alpha}$ has asymptotic level $\alpha$. $H_0$ implies $H_0'$ for every $f$, so the test
is asymptotically valid for $H_0$ whatever the critic; the dependence across windows and across
tickers on the same date is a nuisance that the date sums absorb.
\end{proposition}
\noindent\textit{Proof.} In \cref{app:proofs-dependence}.

\Cref{prop:hac} is stated for a fixed $f$, that is, for fixed training data and growing
evaluation data. Applying it with $f=\fhat$ conditionally on $\Gsig$ requires the evaluation
sequences to satisfy (D1) conditionally on $\Gsig$, in particular to be independent of $\Gsig$. If
the evaluation data depend on the training data, $\E[\bar a-\bar b\mid\Gsig]$ can be non-zero
under $H_0$, and no variance estimator corrects a bias in the mean (\cref{prop:ticker-bias}).
Versions for critics that change with the sample size would need central limit theorems uniform
over the critic class; we do not pursue them.

\begin{fact}[Mixing of the models; \citealp{carrasco2002mixing}]\label{fact:mixing}
Under $a+b<1$, with innovations that have a positive density, GARCH(1,1) is $\beta$-mixing, hence
$\alpha$-mixing, with a geometric rate; stochastic-volatility models inherit $\beta$-mixing from the
volatility process. Window-level coefficients satisfy $\alpha_W(h)\le\alpha_r((h-1)L+1)$, because a
gap of $h$ windows is a gap of at least $(h-1)L+1$ days. Finite $(2+\epsilon)$-th moments of
\GARCHt{} returns require $\nu>2+\epsilon$ and a moment condition on $(a,b)$
\citep{francq_2019_garch}.
\end{fact}

The conditional independence of training and evaluation data fails in a time series. Write
$\beta(\mathcal A,\mathcal B)=\E\sup_{C\in\mathcal B}\abs{\Prob(C\mid\mathcal A)-\Prob(C)}$ for the
absolute-regularity coefficient of two $\sigma$-fields \citep{bradley_2005_basic}.

\begin{lemma}[Decoupling]\label{lem:decoupling}
Let $U$ (training data, $\sigma$-field $\mathcal A$) and $V$ (evaluation data, $\sigma$-field
$\mathcal B$) take values in Polish spaces, and let $\tilde V$ be a copy of $V$ independent of $U$.
For every measurable $E$, $\abs{\Prob((U,V)\in E)-\Prob((U,\tilde V)\in E)}\le\beta(\mathcal A,\mathcal B)$.
\end{lemma}
\noindent\textit{Proof.} In \cref{app:proofs-dependence}.

If training and evaluation data were contiguous blocks separated by a gap of $g$ days, then
$\beta(\mathcal A,\mathcal B)\le\beta_r(g+1)$, the $\beta$-mixing coefficient of the (panel) return
process at lag $g+1$, which decays geometrically for the models of \cref{fact:mixing}. Every level
statement whose remaining hypotheses concern only the evaluation data, \cref{prop:hac} in
particular, would then hold up to an additive $\beta_r(g+1)$, because $\tilde V$ keeps the internal
dependence of $V$. The implemented split is different (\cref{sec:method-inference}): each sample's
date range is cut into $K=10$ blocks of equal length, whole blocks are assigned at random to the
two halves, the blocks interleave, and there is no gap. At every boundary between a training and
an evaluation block the relevant coefficient is the one-day coefficient, which is not small under
persistent volatility, so \cref{lem:decoupling} gives no useful bound. At the E3 sizes a \TEST{}
block holds about three window dates per ticker, so most evaluation windows lie within one or two
windows of a training block, while the volatility half-life at persistence $0.98$ is about 34
days. What the split does guarantee is that training and evaluation dates are disjoint, which
removes the common-date bias of \cref{prop:ticker-bias} below. It does not remove the bias that a
persistent common factor carries from training dates to adjacent evaluation dates
(\cref{cor:block-bias}). The residual dependence across block boundaries is controlled only
empirically: by the simulated placebo of E1, which runs the exact real-data pipeline
(\cref{sec:emp-placebo}), and by E2 part g (\cref{sec:sim-dependent}).

\begin{proposition}[Bias under ticker splits on common dates]\label{prop:ticker-bias}
Let the real windows be $X_{i\tau}=\psi(F_\tau,\varepsilon_{i\tau})$ for tickers $i$ and dates
$\tau\in\mathcal T$, where the factor path $(F_\tau)$ is strictly stationary with marginal $\mu_F$,
the $\varepsilon_{i\tau}$ are i.i.d.\ and independent of $(F_\tau)$, and $\psi$ is measurable. Let
$P^{(v)}$ be the law of $\psi(v,\varepsilon)$, so that every real window has law
$\bar P=\int P^{(v)}\mu_F(dv)$. Let the synthetic windows be mutually independent, independent of
the real data, with law $\bar Q=\bar P$ (the pooled null holds). Split the tickers into training
and evaluation sets independently of the data, each ticker contributing one window per date, and
let $\fhat$ be measurable with respect to $\Gsig\vee\sigma(F)$, for example $\Gsig$-measurable,
where $\Gsig$ is generated by the training windows (real and synthetic), the split and $\xi$.
\begin{enumerate}
  \item[(a)] Conditionally on $(\Gsig,F)$ the real evaluation windows are independent with
  $X_{i\tau}\sim P^{(F_\tau)}$, and $\E[T\mid\Gsig,F]=(\hat P_F-\bar Q)\fhat$, where
  $\hat P_F=\abs{\mathcal T}^{-1}\sum_{\tau\in\mathcal T}P^{(F_\tau)}$ is the realised mixture.
  \item[(b)] $\E T=\E[(\hat P_F-\bar Q)\fhat]$. This vanishes if $\fhat$ is independent of $F$;
  in general it is a covariance between the critic and the realised factor path and need not
  vanish.
  \item[(c)] (Idealised critic.) If $\fhat=\log(d\hat P_F/d\bar Q)$, the Bayes logit of the
  population version of the training problem given $F$ (the limit when the number of training
  tickers per date grows, so that the training half reveals $F$), then
  $\E[T\mid\Gsig,F]=\Jeff(\hat P_F,\bar Q)\in(0,\infty]$ whenever $\hat P_F\neq\bar Q$. For Gaussian
  scale mixtures, $\psi(v,\varepsilon)=v\varepsilon$ with $\varepsilon\sim N(0,I_L)$ and $v>0$, this
  holds almost surely as soon as $\mu_F$ has no atoms.
\end{enumerate}
\end{proposition}
\noindent\textit{Proof.} In \cref{app:proofs-dependence}.

With many training tickers per date, the training half reveals the realised factor path. A critic
that separates real from synthetic training windows then learns that windows typical of the
realised dates are real; the real evaluation windows share those dates and the synthetic ones do
not. The result is a bias, $\E[T\mid\Gsig]\neq0$ and typically positive, under the pooled null.
Driscoll--Kraay and other HAC variances estimate the dispersion of $T$ around $\E[T\mid\Gsig]$ and
cannot remove it; the permutation p-value, which also ignores the cross-sectional dependence, is
worse. A placebo that compares real tickers with real tickers on the same dates cannot reveal the
bias, because both samples share the realised path. The simulated placebo E1(a) runs the ticker
split with permutation p-values as a comparison row; its rejection rate at nominal 5\% is 1.000
[0.981, 1.000], and with the Driscoll--Kraay p-value still 0.990 [0.964, 0.997], as the proposition
predicts for a bias in the mean. The time-block split with the Driscoll--Kraay p-value rejects at
0.085 [0.054, 0.132] (Wilson 95\% intervals, 200 replications; \cref{tab:e1-size-sim}).

The time-block split makes training and evaluation dates disjoint, but the proof of
\cref{prop:ticker-bias}(a)--(b) does not use that the two halves share their dates. It uses only
that, given $(\Gsig,F)$, the real evaluation windows are independent with laws $P^{(F_\tau)}$ and
the synthetic ones are independent of $(\Gsig,F)$ with law $\bar Q$.

\begin{corollary}[Time-block split under a persistent factor]\label{cor:block-bias}
In the model of \cref{prop:ticker-bias}, split the dates instead of the tickers: $\mathcal T$ is
cut into disjoint sets $\mathcal T_{\trn}$ and $\mathcal T_{\evl}$, chosen independently of the
data, and every real and synthetic ticker contributes its windows on $\mathcal T_{\trn}$ to the
training half and those on $\mathcal T_{\evl}$ to the evaluation half. Let $\fhat$ be
$\Gsig$-measurable and $\hat P_{F,\evl}=\abs{\mathcal T_{\evl}}^{-1}\sum_{\tau\in\mathcal
T_{\evl}}P^{(F_\tau)}$. Then $\E[T\mid\Gsig,F]=(\hat P_{F,\evl}-\bar Q)\fhat$ and
$\E T=\E[(\hat P_{F,\evl}-\bar Q)\fhat]$. This vanishes if $(F_\tau)_{\tau\in\mathcal T_{\evl}}$ is
independent of $(F_\tau)_{\tau\in\mathcal T_{\trn}}$. In general it is a covariance between the
critic, which depends on the factor path on the training dates, and the factor path on the
evaluation dates, and it need not vanish when $(F_\tau)$ is persistent across adjacent blocks.
\end{corollary}
\begin{proof}
The real evaluation windows are functions of $(F_\tau,\varepsilon_{i\tau})_{\tau\in\mathcal
T_{\evl}}$. $\Gsig$ depends on the real data only through $(F_\tau,\varepsilon_{i\tau})_{\tau\in
\mathcal T_{\trn}}$, so the $\varepsilon_{i\tau}$ with $\tau\in\mathcal T_{\evl}$ are independent of
$(\Gsig,F)$, and the proof of \cref{prop:ticker-bias}(a)--(b) applies with $\mathcal T_{\evl}$ in
place of $\mathcal T$, under the same integrability. If $(F_\tau)_{\tau\in\mathcal T_{\evl}}$ is
independent of $(F_\tau)_{\tau\in\mathcal T_{\trn}}$, it is independent of $\Gsig$, and
$\E[P^{(F_\tau)}\fhat\mid\Gsig]=\bar P\fhat=\bar Q\fhat$ for $\tau\in\mathcal T_{\evl}$ by
stationarity.
\end{proof}

With a persistent factor the realised window law on the evaluation dates resembles the one on the
adjacent training dates, which the critic has learned to score as real. The corollary is stated,
like \cref{prop:ticker-bias}, for mutually independent synthetic windows, the situation of
univariate generators that simulate tickers independently; when the synthetic panel has a common
factor of its own, the same computation gives $(\hat P_{F,\evl}-\hat Q_{F',\evl})\fhat$ with the
realised mixture $\hat Q_{F',\evl}$ of the synthetic factor path, which need not vanish either.
\Cref{rem:small-T} compares this mechanism with the E1 measurements.

\begin{remark}[Few evaluation dates]\label{rem:small-T}
For features driven by common factors, precision is governed by the number of distinct evaluation
dates, not by tickers times windows. At $L=20$ the \TEST{} period gives 30 window dates per ticker,
about 15 of them in the evaluation half. HAC estimates from so few dates are imprecise and biased
downwards, and normal critical values then over-reject. Fixed-$b$ critical values matched to the
Bartlett kernel and bandwidth \citep{kiefer_2005_new}, or a cluster variance over time blocks with
$t_{G-1}$ critical values \citep{bester_2011_inference}, are the standard small-sample
corrections; the two pairings should not be mixed. The pipeline reports the normal-reference
Driscoll--Kraay p-value; fixed-$b$ p-values are computed as comparison rows in E1, which measures
the size of the complete pipeline (\cref{sec:emp-placebo}). There, at nominal 5\%, the \WC{} test
rejects at 0.085 with the normal reference and at 0.050 with the fixed-$b$ reference for equal
variances when the synthetic panel shares the factor structure of the real one, but at 0.270 and
0.245 when the synthetic tickers are independent (\cref{tab:e1-size-sim}). The second case is not
a small-sample variance effect alone: the mean of $\zhac$ is 0.74 there, against 0.11 in the first
case (\cref{tab:e1-sim-z,sec:sim-dependent}). \Cref{cor:block-bias} suggests a mechanism, which
is model-based and heuristic. The conditional mean $(\hat P_{F,\evl}-\bar Q)\fhat$ of $T$ is a
function of the realised factor path, and so is the Driscoll--Kraay standard error, through the
date-level variance of the real critic values. When the synthetic tickers are independent, their
date sums carry no common factor, the standard error is driven by the real path alone, and a
conditional bias that is small on average can still shift $\zhac$ if it is dependent on that
standard error. A synthetic factor of its own adds variation to the standard error that is
independent of the real path, which weakens this dependence.
\Cref{tab:manual-e1-sim-stat} is consistent with this reading. The mean of $T$ for the \WC{} test
is small and nearly the same in the two cells (0.026 and 0.029, with Monte Carlo standard errors
0.015 and 0.014), but with independent synthetic tickers $T$ and its standard error have
Pearson correlation $-0.64$, against 0.22 in the joint cell, so that positive values of $T$ come with
small standard errors; the \LRC{} test shows the same sign pattern, and the rejections concentrate
in the replications with a small standard error (\cref{tab:e1-sim-se}). The E1 panel lies outside the
model of the corollary, because its idiosyncratic and synthetic GARCH components are also
serially dependent within a ticker, and the corollary does not quantify the effect.
\end{remark}

\begin{table}[tbp]
  \centering
  \caption{Statistic and Driscoll--Kraay standard error of the absolute learned tests in the E1(a)
    placebo (time-block split, 200 replications per cell). joint: the synthetic panel is an
    independent draw of the whole factor panel; indep.: independent synthetic tickers. Mean of
    $T$ with its Monte Carlo standard error, standard deviation of $T$, share of replications with
    $T>0$, median Driscoll--Kraay standard error of $T$, correlation of $T$ with that standard
    error over the replications, and mean of $\zhac$ (as in \cref{tab:e1-sim-z}). Derived from
    \texttt{results/e1/e1\_sim\_raw.csv} by \texttt{paper/scripts/e1\_sim\_stat.py}.}
  \label{tab:manual-e1-sim-stat}
  \small
  \fittable{manual_e1_sim_stat}
\end{table}

\begin{remark}[Wild bootstrap for the kernel tests]\label{rem:wild}
MMD and \SigMMD{} have no train/evaluation split. Under dependence they use the wild bootstrap of
\citet{chwialkowski2014wild} for the MMD V-statistic, with a Gaussian AR(1) multiplier process on
the date grid that is shared by all windows with the same date rank, in both samples, and a
plug-in correlation length (\cref{app:impl}). \citet{chwialkowski2014wild} prove validity for
weakly dependent stationary processes with a given correlation length; the date-shared multipliers
and the plug-in length are our choices and are not covered by their theorem. We give no proof;
their size is measured in E1 and in E2 part g.
\end{remark}

% ---------------------------------------------------------------------------------------------
\subsection{What the verdict controls}\label{sec:theory-verdict}

\begin{definition}[Power check]\label{def:power-check}
A power check is a pre-specified pair $(\Cdesign,\vartheta^*)$. $\Cdesign$ is a reference
simulation design that reproduces the empirical one: the sizes of $\Rpast$, $\Rnew$ and $G$, the
window length, the split, the critic settings, the inference, the seed combination and, where
possible, the dependence. $\vartheta^*$ is a pre-specified deviation of the generator. The
boundary generator $G^0$ is an independent draw with the law of the past windows. We use two
designs. D1 is a per-ticker \GARCHt{} panel fitted on \TRAIN: $\Rpast^*$ is a \TRAIN-length path
started from stationarity, $\Rnew^*$ a \TEST-length path started from the fitted end-of-\TRAIN{}
state, $G^0$ a \TEST-length path started from stationarity, and the alternative $G^*$ is $1.2$
times an independent \TEST-length path started from the end-of-\TRAIN{} state. In D2 the past
windows are the real \TRAIN{} realised volatilities times fresh Gaussian noise, $\Rnew$ is the
real \TEST{} sample, $G^0$ has the pooled law of the past windows, and $G^*=1.2\,G^0$. In D2 the
alternative differs from the boundary only by the volatility factor; in D1 it also differs by its
initial state, so the D1 power is the power against a compound deviation, not against the
volatility factor alone. The power $\beta(\vartheta^*)$ is the probability that the verdict rule
rejects under $G^*$, estimated from $R$ replications by $\hat\beta$, with one-sided
Clopper--Pearson lower bound $\betaL$ \citep{clopper_1934_use}. The same design with $G^0$ in place
of $G^*$ gives the boundary size $\abd$, estimated by $\hat\alpha_{\mathrm{bd}}$. Since $\abd$ can
exceed $\alpha$, $\hat\beta$ mixes the discrimination between $G^*$ and $G^0$ with the
miscalibration of the rule.
The size-adjusted power $\beta_{\mathrm{adj}}(\vartheta^*)=\Prob(p^{\mathrm{alt}}\le c_\alpha)$,
where $p^{\mathrm{alt}}$ is the rule's $\pmed$ under $G^*$ and $c_\alpha$ the $\alpha$-quantile of
its law $p^{\mathrm{bd}}$ under $G^0$, measures the discrimination alone. It is estimated with
$\hat c_\alpha$, the $\lfloor\alpha R\rfloor$-th smallest of the $R$ boundary values, at which the
empirical boundary size is $\lfloor\alpha R\rfloor/R\le\alpha$. We recommend reporting
$\hat\beta_{\mathrm{adj}}$ together with $\hat\alpha_{\mathrm{bd}}$ and its Clopper--Pearson
interval.
\end{definition}

\begin{definition}[Verdict]\label{def:verdict}
Let $\alpha=0.05$, $\bmin=0.8$ and let $\pmed$ be the anchored p-value combined over the synthetic
seeds (\cref{rem:combine-seeds}); as implemented, each single-seed p-value is the Driscoll--Kraay
p-value $p_s=1-\Phi(\zhac)$ of $\Danc$. The verdict \eqref{eq:verdict} is $V=\BAD$ if
$\pmed\le\alpha$; $V=\GOOD$ if $\pmed>\alpha$ and $\hat\beta\ge\bmin$ under every design (D1 and D2); and
$V=\INCONCL$ otherwise. The stationary bootstrap, for which the anchored test is invalid
(\cref{rem:anchored-null}), receives \INVALID. A stricter variant replaces $\hat\beta$ by $\betaL$.
\end{definition}

\begin{proposition}[What the verdict controls]\label{prop:verdict}
(a) If each single-seed anchored p-value $p_s$ satisfies $\Prob(p_s\le u)\le u$ for all
$u\in(0,1)$, then $\Prob(V=\BAD)\le\alpha$. For the permutation version of the anchored p-value
this premise holds in finite samples at $\QG=\Ppast$ under \cref{ass:anc} (\cref{prop:anchored}(a)).
The rule as implemented uses the Driscoll--Kraay p-values (\cref{def:verdict}), for which the
premise holds only in the limit, also when the windows are independent: at $\QG=\Ppast$ under the
conditions of \cref{prop:anchored}(c), where $\delta=0$ by \cref{prop:anchored-target}(a),
$\limsup\Prob(V=\BAD)\le\alpha$. On real data, where the past is non-stationary
(\cref{rem:nonstat-past}), the premise does not hold; there the boundary size $\abd$ of the rule
is estimated (\cref{def:power-check}), and a \BAD{} verdict has a real-data error interpretation
only when the Clopper--Pearson interval of $\hat\alpha_{\mathrm{bd}}$ covers $\alpha$ or lies below
it.
(b) If the rejection probability of the rule at the true configuration is at least
$\beta(\vartheta^*)$, then $\Prob(V=\GOOD)\le\Prob(\pmed>\alpha)\le1-\beta(\vartheta^*)$, which is
at most $1-\bmin$ up to the Monte Carlo error of $\hat\beta$.
\end{proposition}
\begin{proof}
(a) $\{V=\BAD\}=\{\pmed\le\alpha\}$; apply \cref{rem:combine-seeds}. Under
\cref{prop:anchored}(c), its Markov bound $\Prob(\pmed\le\alpha)\le(2/S)\sum_s\Prob(p_s\le\alpha/2)$
holds for every sample size, and $\limsup\Prob(p_s\le\alpha/2)\le\alpha/2$ for each seed follows
from the conditional limit by dominated convergence.
(b) $\{V=\GOOD\}\subset\{\pmed>\alpha\}$.
\end{proof}

That the Driscoll--Kraay version is valid only in the limit is visible at the sizes used. The
single-seed Driscoll--Kraay anchored p-value rejects at the boundary in 0.09 [0.04, 0.16] of the D1
replications (Clopper--Pearson; \cref{tab:e3-power}) and in 0.080 [0.050, 0.126] of the
replications of the joint cell of E1 (Wilson; \cref{tab:e1-size-sim}), two stationary simulated
panels; in D1 the seed combination absorbs the excess, with a boundary size of the rule of 0.00
[0.00, 0.04]. The \BAD{} condition at the end of (a) is a reporting diagnostic, not a change to the
pre-specified rule; \cref{rem:calibrated-bad} gives the calibration that replaces the nominal level
when it fails.

Part (b) is only as strong as its premise, which cannot be verified: the true deviation must be at
least as detectable as $\vartheta^*$ under $\Cdesign$. \GOOD{} therefore means ``no detectable
deviation of the pre-specified size and type'', never ``equal to the truth''. Our $\vartheta^*$ is
an upward volatility error, which in D2 points away from the calmer \TEST{} period. Power against
downward volatility errors or errors in other features is not covered (a factor 0.8 was not
run), and under a non-stationary past a non-rejection of the anchored null can reflect a generator
bias aligned with the regime change (\cref{rem:anchored-null}). The D2 boundary size, in turn,
is that of a boundary generator with Gaussian windows whose scales are drawn independently across
tickers and windows (\cref{sec:method-decision}). D1 is a design-stage
calculation, like a sample-size calculation. D2 is a power analysis conditional on the realised
$\Rnew$ and on the \TRAIN{} volatility path: it uses the evaluation data, though not the generators,
so \GOOD{} against \INCONCL{} depends on $\Rnew$ through $\pmed$ and through $\hat\beta$. Part (b)
is unaffected, because $\{V=\GOOD\}\subset\{\pmed>\alpha\}$. With $R=100$ replications the standard
error of $\hat\beta$ near $0.8$ is $0.04$, which motivates the Clopper--Pearson variant.

\begin{remark}[Calibrating \BAD{} against the boundary replications]\label{rem:calibrated-bad}
Let $p^{\mathrm{bd}}_1,\dots,p^{\mathrm{bd}}_R$ be the rule's $\pmed$ in the $R$ boundary
replications of a design and $p^{\mathrm{obs}}$ the observed $\pmed$ of a generator. The Monte
Carlo p-value
\[
  p_{\mathrm{MC}}=\frac{1+\#\{r:p^{\mathrm{bd}}_r\le p^{\mathrm{obs}}\}}{R+1}
\]
satisfies $\Prob(p_{\mathrm{MC}}\le u)\le u$ for every $u\in(0,1)$ if, under the null,
$(p^{\mathrm{obs}},p^{\mathrm{bd}}_1,\dots,p^{\mathrm{bd}}_R)$ is exchangeable: apply
\cref{lem:rank} to the negated values. This is the Monte Carlo test principle
\citep{dufour_2006_monte}, with ties resolved conservatively, and its validity rests entirely on
the exchangeability. D2 keeps the real $\Rnew$, as the observed value does, but provides
exchangeability only approximately. Its past windows keep the real \TRAIN{} volatility path but
not the real within-window law (Gaussian noise, so no conditional tails, within-window clustering
or leverage) nor the same-date return correlation across tickers; its boundary generator $G^0$
has the pooled \TRAIN{} window law but not the structure of a fitted model; and on real data the
standard error of $\Danc$ is 0.59 to 0.79 times its D2 value (\cref{app:emp-derived}). The Monte
Carlo p-values of \cref{sec:emp-verdicts} are therefore approximate calibrations, more
informative than the nominal level when $\hat\alpha_{\mathrm{bd}}>\alpha$, but not exact. The
size-adjusted power of \cref{def:power-check} separates discrimination from miscalibration
(\cref{tab:manual-theory-calibration}). At the main split the D2 power of the rule is 0.45 at the
nominal threshold, with boundary size 0.08, and still 0.45 at $\hat c_\alpha=0.025$, where the
boundary size is 0.05; in all 100 replications the alternative $\pmed$ lies below the boundary
$\pmed$ of the same replication. At the alternative split of E5 the D2 power and boundary size are
both 0.60 at the nominal threshold, and the size-adjusted power is still 0.60 at
$\hat c_\alpha=3.8\times10^{-5}$ (boundary size 0.04). There the rule still separates the boundary
generator from one with 20\% too much volatility, but the Clopper--Pearson interval of
$\hat\alpha_{\mathrm{bd}}$ lies above $\alpha$, so a nominal \BAD{} verdict would have no
real-data error interpretation.
\end{remark}

\begin{table}[tbp]
  \centering
  \caption{Calibration of the verdict rule in the power-check designs: nominal boundary size and
    power at $\alpha=0.05$, estimated $\alpha$-quantile $\hat c_\alpha$ of the boundary $\pmed$
    (the $\lfloor\alpha R\rfloor$-th smallest of the $R$ boundary values), boundary size and
    size-adjusted power at $\hat c_\alpha$ (\cref{def:power-check}), and the number of replications
    in which the alternative $\pmed$ is below the boundary $\pmed$ of the same replication. Main
    split: E3 (\cref{tab:e3-power}); alternative split: E5 (\cref{tab:manual-e5-split}). $R$
    replications, $S$ synthetic seeds. Derived from \texttt{results/e3/e3\_power\_raw.csv},
    \texttt{e3\_realvol\_raw.csv} and \texttt{results/e5/e5\_split\_power\_raw.csv} by
    \texttt{paper/scripts/theory\_calibration.py}.}
  \label{tab:manual-theory-calibration}
  \small
  \fittable{manual_theory_calibration}
\end{table}

\begin{remark}[Equivalence-test alternative]\label{rem:equivalence}
A \GOOD{} verdict with a formal error guarantee would follow from a non-inferiority test
\citep{schuirmann_1987_comparison}: test $\delta\ge\varepsilon_0$ against $\delta<\varepsilon_0$ and
declare \GOOD{} if the upper $(1-\alpha)$ confidence bound for $\delta$ lies below a margin
$\varepsilon_0$, for instance the expected anchored discrepancy at $\vartheta^*$ in the power-check
design. The guarantee would again be conditional on the trained critic, and the margin transfers
across configurations only to the extent that critic training behaves similarly (heuristic). We do
not use it for the verdict.
\end{remark}

\begin{remark}[Per-ticker tests and multiple testing]\label{rem:bh}
The per-ticker tests of E3 ($L=5$) give p-values $p_1,\dots,p_N$, calibrated against a pooled
\GARCHt{} null (\cref{sec:emp-perticker}). The Benjamini--Hochberg procedure at level $q$ controls
the false discovery rate of \BAD{} verdicts at $qN_0/N$ under independence or positive regression
dependence \citep{benjamini1995controlling,benjamini2001control}. Per-ticker statistics are
plausibly positively dependent through the market factor, which we do not prove, so the
Benjamini--Yekutieli correction is reported as a robustness check and Storey's $\hat\pi_0$
\citep{storey_2002_direct} as a summary. \GOOD{} is a non-rejection, so no false-discovery
statement applies to it, and with about 60 evaluation windows per ticker few per-ticker \GOOD{}
verdicts can be expected.
\end{remark}

% ---------------------------------------------------------------------------------------------
\subsection{Window distances and portfolio utility}\label{sec:theory-markowitz}

The Markowitz experiment E4 asks whether test results predict the utility of portfolios built
from synthetic data. The theory of \cref{sec:mk-theory} (proofs in \cref{app:proofs-markowitz})
predicts a link that is real but not monotone. For the unconstrained global minimum-variance
portfolio $\wgmv(\hat\Sigma)=\hat\Sigma^{-1}\one/(\one'\hat\Sigma^{-1}\one)$, evaluated with the
\TEST{} covariance $\Sigma$, the variance loss is a squared Mahalanobis distance to the ex-post
optimum (\cref{lem:quad-loss}), the weights depend on $\hat\Sigma$ only up to scale, and
\begin{equation}\label{eq:kantorovich-preview}
  1\le\frac{V(\wgmv(\hat\Sigma))}{V(\wgmv(\Sigma))}\le\frac{(1+\kappa)^2}{4\kappa},
  \qquad \kappa=\frac{\lambda_{\max}(\Sigma^{-1/2}\hat\Sigma\Sigma^{-1/2})}
  {\lambda_{\min}(\Sigma^{-1/2}\hat\Sigma\Sigma^{-1/2})},
\end{equation}
a sharp bound that is quadratic in the relative covariance error (\cref{prop:gmv-sensitivity}).
In the large-sample limit the GMV utility of a generator therefore depends on it only through its
covariance matrix, up to scale (\cref{cor:gmv-drivers}); tails, volatility clustering and jumps,
which the critic detects, matter only through that matrix. The two criteria can disagree in both
directions. Here $\kappa$ is the ratio of the extreme eigenvalues of $\Sigma^{-1}\hat\Sigma$,
written $\cond(\Sigma^{-1}\hat\Sigma)$ in \cref{sec:mk-theory}. A generator whose covariance is
$c\Sigma$ with $c\neq1$ is rejected by a consistent
critic test and loses nothing for GMV (\cref{ex:scale-blind}); a covariance error in a
low-variance spread direction attains the bound \eqref{eq:kantorovich-preview} while $\W$ between
the window laws tends to zero (\cref{ex:spread}); univariate tests cannot see correlations
(\cref{ex:corr-blind}); and a drift error of one asset equal to one \TRAIN{} standard error
dominates mean-variance utility, yet by \cref{prop:oracle-bound}(b) no level-5\% two-sample test
that uses the 615 \TEST{} days of that asset detects it with power above $0.17$
(\cref{ex:drift}). Finally, under convex
losses no portfolio rule trained on synthetic data beats the best rule based on the real past
(\cref{prop:no-free-lunch}); synthetic data can only help as a regulariser. These results motivate
the design of E4 (\cref{sec:mk-link}): an association between critic distance and GMV loss is
plausible through covariance accuracy, little association is expected for mean-variance
portfolios, which are driven by drift, and univariate verdicts carry no information about
correlations. Because the real correlations themselves changed between \TRAIN{} and \TEST{}
(\cref{tab:e4-correlation}), fidelity to the past and usefulness for the future are different
targets.

% ---------------------------------------------------------------------------------------------
\begin{table}[tbp]
  \centering
  \caption{Status of the theoretical results. ``Proved'': full proof in \cref{app:proofs} or in
  the text; ``sketch'': key steps only; ``heuristic'': argument, not a theorem.}
  \label{tab:theory-status}
  \small
  \begin{tabular}{@{}p{0.43\linewidth}p{0.17\linewidth}p{0.34\linewidth}@{}}
    \toprule
    Result & Status & Key assumptions \\
    \midrule
    \Cref{prop:window-ident,prop:markov-pairs}: what windows identify & proved & stationarity; both processes Markov \\
    \Cref{prop:unbiased}: $\E[T\mid\Gsig]=\Dcrit(\fhat)\le\Lip(\ghat)\W(P_A,Q_A)$ & proved & evaluation data independent of training \\
    \Cref{prop:exact-level}: exact level of the permutation test & proved & conditional exchangeability (A1$'$) \\
    \Cref{prop:consistency,cor:w1-limit}: consistency; $T\to\W$ & proved & (C1)--(C5), incl.\ hard Lipschitz constraint \\
    \Cref{prop:bayes-logit,cor:jeffreys}: Bayes logit, Jeffreys & proved & perfect discriminator \\
    \Cref{prop:oracle-bound}: oracle bound & proved & densities; level at the null point \\
    \Cref{prop:local-power}: local power factor & proved (fixed $f$); sketch (learned) & differentiability in quadratic mean \\
    \Cref{rem:lr-vs-w1}: \LRC{} vs \WC{} & heuristic & -- \\
    \Cref{prop:anchored}(a): anchored test exact at $\QG=\Ppast$ & proved & \cref{ass:anc} \\
    \Cref{prop:anchored}(b)--(c): weaker null; HAC & proved (standard results) & i.i.d.\ or \cref{ass:mixing}, given $\Gsig$ \\
    \Cref{prop:anchored-target}: meaning of $\delta$ & proved & idealised critics in (b), (c) \\
    \Cref{rem:nonstat-past}: non-stationary past & heuristic & -- \\
    \Cref{prop:paired-boot}: paired bootstrap & sketch & \cref{ass:rel} \\
    \Cref{prop:boundary-size}: two-critic size $>\alpha$ at boundary & proved (given \cref{prop:paired-boot}) & \cref{ass:rel}; exchangeable training \\
    \Cref{rem:naive-perm}: naive permutation distortion & heuristic & -- \\
    \Cref{prop:hac}: HAC test & proved (standard results) & \cref{ass:mixing}; fixed critic \\
    \Cref{lem:decoupling}: decoupling & proved & no useful bound without a gap \\
    \Cref{prop:ticker-bias}: ticker-split bias & proved & factor model \\
    \Cref{cor:block-bias}: time-block bias under a persistent factor & proved & factor model; independent synthetic windows \\
    \Cref{rem:small-T}: mechanism of the E1 oversize & heuristic & -- \\
    \Cref{prop:verdict}: verdict & proved under premises & valid $p_s$ (asymptotic only for the Driscoll--Kraay $p_s$ used); unverifiable power premise \\
    \Cref{rem:calibrated-bad}: Monte Carlo calibration of \BAD{} & proved (Monte Carlo test) & exchangeable observed and boundary $\pmed$ (approximate in D2) \\
    \Cref{sec:mk-theory}: GMV and utility results & proved; parts heuristic & positive definite matrices; convex loss \\
    \bottomrule
  \end{tabular}
\end{table}

\section{Simulation evidence}\label{sec:sim}
% =============================================================================================
% Target length: 5 pages including 2 figures and 3 tables (rest in \cref{app:sim}).
% Experiments: E2 (experiments/e2_ground_truth.py: parts main_a..main_d, e, f, g) and E2-dep
% (experiments/e2_dependence.py: families mix, block). Independent windows (one burned-in path
% per window) => permutation p-values exact, except part g. Numbers only from the generated
% tables; quote rates with their Wilson intervals. Plan: OUTLINE.md, Section 5.
%
% Planned subsections (labels):
%  \subsection{Design}                                             \label{sec:sim-design}
%     base laws (daily vol 0.02, mean 3e-4): t(5); t(4); GARCH(1,1)-t a = 0.08, b = 0.90, nu = 5;
%     Merton (lambda_J = 0.05, jumps N(-0.01, 0.04^2)); GBM. Families (a) volatility level,
%     (b) tails, (c) clustering at matched variance, (d) jumps (window mixture); n in
%     {500, 2000} windows of L = 20; 200 replications under H_0, 100 otherwise; common random
%     numbers across levels; oracle exact for (a), (b), (d), Monte Carlo GARCH density (512
%     draws) for (c). E2-dep: marginal held exactly fixed, dependence changed (mix, block);
%     \todo{E2-dep sizes as run: n = 2000 only, 100/50 replications (compute budget)}.
%  \subsection{Size}                                               \label{sec:sim-size}
%     all window tests hold their level with independent windows (tab:e2-size in appendix,
%     summarised here); KS oversized in family (c) (within-window dependence).
%  \subsection{Power against marginal and dependence deviations}   \label{sec:sim-power}
%     fig:e2-power + tab:e2-power: clustering at matched variance -> window tests have power
%     (critic close to MMD and Sig-MMD, above C2ST and LR-critic); marginal deviations (vol
%     level, tails, jumps) -> pooled KS far more powerful than every window test; distance to
%     the oracle; interpret with prop:local-power (factor <= 1/2) and prop:oracle-bound.
%     fig:e2dep-power: dependence-only deviations, KS powerless by construction (and invalid).
%     State plainly where simpler tests win.
%  \subsection{The relative test at its boundary}                  \label{sec:sim-relative}
%     tab:e2-relative-boundary (+ note): anchored size at Law(G) = Law(R_past) versus two-critic
%     over-rejection (prop:boundary-size, prop:anchored); interior cells (G = P_new better,
%     GARCH 0.90 / GBM worse); resample generator: grouped split (size) vs naive split
%     (leakage demonstration, not a size). n = 500 only (budget).
%  \subsection{Learned versus exact log-likelihood ratio}          \label{sec:sim-llr}
%     tab:e2-llr (+ note), fig:e2-llr-scatter (appendix): Spearman / calibration slope of the
%     learned logit vs log p/q for GBM vs Merton and GARCH 0.98 vs 0.90, three training budgets.
%  \subsection{Dependent windows}                                  \label{sec:sim-dependent}
%     tab:e2-dependent-size: consecutive windows of one GARCH path; random-split permutation
%     oversized; time-block + HAC and wild bootstrap closer to nominal; motivates the real-data
%     protocol (link to rem:naive-perm, prop:hac).
% Tables/figures (file -> label; wrapper template in OUTLINE.md):
%   figures/e2_power.pdf -> fig:e2-power (main)        tables/e2_power.tex (+e2_power_note) -> tab:e2-power (main)
%   figures/e2dep_power.pdf -> fig:e2dep-power (main)  tables/e2_relative_boundary.tex (+note) -> tab:e2-relative-boundary (main)
%   tables/e2_dependent_size.tex -> tab:e2-dependent-size (main)
%   tables/e2_llr.tex (+e2_llr_note) -> tab:e2-llr (main or appendix)
%   appendix: tab:e2-size, tab:e2-power-ci, tab:e2dep-size, tab:e2dep-power, fig:e2-llr-scatter
% =============================================================================================
\todo{Section to be written (OUTLINE.md, Section 5). Anchor citations:
\citet{bollerslev1986generalized}, \citet{merton1976option}, \citet{lehmann2022testing}.}

\section{Data and generators}\label{sec:data}
% =============================================================================================
% Target length: 2.5 pages including 2 tables and 1 figure (fit details in \cref{app:data}).
% Sources: docs/DESIGN.md Section 2 and Change log; src/synthval/data.py, models/*.py;
% experiments/e3_sp500.py parts fit and stylized. Plan: OUTLINE.md, Section 6.
%
% Planned subsections (labels):
%  \subsection{The S\&P 500 panel}                                 \label{sec:data-panel}
%     Bloomberg PX_LAST of the 503 current constituents, 2018-01-02 to 2025-06-16; complete
%     universe = no missing price (internal gaps <= 3 days forward-filled, never back-filled):
%     481 tickers; two data errors excluded by a stated rule (max |log-return| > 1 or >= 20
%     zero-volume days: KDP, unadjusted special dividend; AMCR, stale pre-listing quotes) ->
%     479; main panel = 100 tickers drawn with a fixed seed; E4 universes = 30 x 30 from the 479.
%     TRAIN 2018-01-02 to 2022-12-30 (1258 days), TEST 2023-01-03 to 2025-06-16 (615 days);
%     windows: 62 per ticker in TRAIN (6200), 30 in TEST (3000) at L = 20. Survivorship bias
%     (current constituents) stated as a limitation.
%  \subsection{Generators}                                         \label{sec:data-generators}
%     univariate (per ticker, fitted on TRAIN): GBM (MLE), Merton (MLE, truncated Poisson sum),
%     GARCH(1,1)-t (MLE, variance targeting, persistence cap 0.995), Heston (Andersen QE scheme,
%     c <= sqrt(1.5) to avoid zero-variance days, SMM on variance, tail quantile ratio,
%     ACF(r^2) at 1/5/10/20, leverage), stationary bootstrap (Politis-White block length);
%     simulation over TEST from the fitted end state; likelihood available for GBM, Merton,
%     GARCH-t only. Multivariate (E4): MV Gaussian, MV Student-t (ECME), CCC-GARCH-t (Gaussian
%     copula), MV stationary bootstrap, factor GARCH-t. Fit diagnostics: tab:e3-fit-diagnostics;
%     Heston moment match (simulated kurtosis below the data): tab:e3-heston-moments (appendix).
%  \subsection{Stylized facts}                                     \label{sec:data-stylized}
%     tab:e3-stylized and fig:e3-stylized: real TEST vs each generator (panel median and IQR of
%     moments, ACF of |r|, tails, aggregational Gaussianity, leverage).
% Tables/figures:
%   tables/e3_fit_diagnostics.tex -> tab:e3-fit-diagnostics (main)
%   tables/e3_stylized.tex -> tab:e3-stylized (main);  figures/e3_stylized.pdf -> fig:e3-stylized (main)
%   tables/e3_heston_moments.tex -> tab:e3-heston-moments (appendix)
% =============================================================================================
\todo{Section to be written (OUTLINE.md, Section 6). Anchor citations:
\citet{cont2001empirical}, \citet{heston1993closed}, \citet{politis1994stationary},
\citet{brown1992survivorship}.}

\section{Empirical results: S\&P 500}\label{sec:emp}
% =============================================================================================
% Target length: 5.5 pages including 6 tables and 3 figures. Experiments: E1
% (experiments/e1_placebo.py) and E3 (experiments/e3_sp500.py). Every real-data test uses the
% dependent protocol (time blocks + Driscoll-Kraay HAC; wild bootstrap for kernel tests; KS
% flagged invalid). Report the pre-specified power check BEFORE the real-data results, as it was
% run. Plan: OUTLINE.md, Section 7.
%
% Planned subsections (labels):
%  \subsection{Size of the real-data pipeline}                     \label{sec:emp-placebo}
%     E1(a) simulated panel placebo (factor GARCH-t fitted on TRAIN; cells joint / indep;
%     200 replications; exactly the E3 sizes): tab:e1-size-sim, fig:e1-pvalue-qq. The naive
%     protocol (ticker split + permutation) is severely oversized (\todo{rate}, theory pilot
%     about 45% at 5%); time-block + HAC closer to nominal but still above (\todo{rates});
%     fixed-b comparison rows. E1(b) real-ticker halves (secondary, 50 splits, rates only:
%     tab:e1-size-real in appendix). E1(c) R_new vs R_past: non-stationarity
%     (tab:e1-nonstat).
%  \subsection{Pre-specified power check}                          \label{sec:emp-power}
%     tab:e3-power: designs D1 (per-ticker GARCH-t, stationary R_past*) and D2 (real TRAIN
%     volatility path); power and boundary size of the seed-combined rule (100 replications x
%     5 seeds as run); a single-seed anchored p-value is not calibrated on real data.
%  \subsection{Absolute and relative tests by generator}           \label{sec:emp-tests}
%     tab:e3-pooled-tests, fig:e3-pvalues: pooled absolute tests per model (5 seeds, twice the
%     median p) and anchored / two-critic relative tests; typical finding to discuss:
%     distinguishable from the future but not worse than history.
%  \subsection{What the critic detects}                            \label{sec:emp-witness}
%     fig:e3-critic-witness: critic values against window log realised volatility and max |r|.
%  \subsection{Per-ticker tests}                                   \label{sec:emp-perticker}
%     tab:e3-perticker: L = 5, critic (HAC) and MMD (wild), p-values calibrated against a pooled
%     GARCH-t null (nominal p-values oversized: tab:e3-perticker-null in appendix), BH primary,
%     BY and Storey pi_0 as robustness; per-ticker GOOD needs per-ticker power (rarely).
%  \subsection{Agreement with likelihood benchmarks}               \label{sec:emp-likelihood}
%     tab:e3-likelihood (predictive log-likelihood, Berkowitz, PIT-KS for GBM / Merton /
%     GARCH-t), tab:e3-likelihood-ag (appendix), tab:e3-agreement (Kendall tau of rankings,
%     Cohen kappa of per-ticker verdicts); disagreement is expected when models are misspecified
%     in different directions (rem:kl-vs-w1).
%  \subsection{Verdicts}                                           \label{sec:emp-verdicts}
%     tab:e3-verdicts: GOOD / BAD / INCONCLUSIVE per model (INVALID for the bootstrap), with
%     the absolute test alongside and the Clopper-Pearson variant.
% Tables/figures (file -> label):
%   e1_size_sim.tex -> tab:e1-size-sim (main); e1_pvalue_qq.pdf -> fig:e1-pvalue-qq (main)
%   e1_nonstat.tex -> tab:e1-nonstat (main or appendix); e1_size_real.tex -> tab:e1-size-real (app.)
%   e3_power.tex -> tab:e3-power (main); e3_pooled_tests.tex -> tab:e3-pooled-tests (main)
%   e3_pvalues.pdf -> fig:e3-pvalues (main); e3_critic_witness.pdf -> fig:e3-critic-witness (main)
%   e3_perticker.tex -> tab:e3-perticker (main); e3_perticker_null.tex -> tab:e3-perticker-null (app.)
%   e3_likelihood.tex -> tab:e3-likelihood (main); e3_likelihood_ag.tex -> tab:e3-likelihood-ag (app.)
%   e3_agreement.tex -> tab:e3-agreement (main); e3_verdicts.tex -> tab:e3-verdicts (main)
% =============================================================================================
\todo{Section to be written (OUTLINE.md, Section 7). Anchor citations:
\citet{diebold1998evaluating}, \citet{berkowitz2001testing}, \citet{amisano2007comparing},
\citet{benjamini1995controlling}.}

\section{Markowitz utility of synthetic data}\label{sec:markowitz}
% =============================================================================================
% Owner: markowitz writer. Theory: docs/theory.md Section 9 (proofs in app:proofs-markowitz).
% Experiment: E4 (experiments/e4_markowitz.py), as run: 30 universes x 30 assets, B = 20
% replicas, 2 test seeds, 60 placebo replications per generator (results/e4/meta*.json).
% Main-text tables: e4_utility, e4_correlation (generated); manual_e4_tests_main,
% manual_e4_link_main (condensed, written by paper/scripts/e4_tests_main.py and
% e4_link_main.py from results/e4/*.csv). Full generated test and link tables: app:mk.
% Derived numbers not in a generated table (counts over universes, gross exposures, the sign
% count of the link tables, the 400 distinct tickers of results/e4/universes.csv, the reference
% start date 2020-07-24 = TRAIN day 643) were computed from results/e4/*.csv and the data
% loader; app:mk repeats them. Section 8.1 keeps the statements (lem:quad-loss ... ex:drift) that
% theory.tex (sec:theory-markowitz), appendix_proofs.tex and the introduction reference.
% Reliability table manual_e4_reliability (tab:manual-e4-reliability): seed-to-seed Spearman
% correlations of the four test quantities, written by paper/scripts/e4_reliability.py from
% results/e4/e4_tests.csv; its header comment records the critic's 60 to 63 training windows per
% sample (n_x - n_eval_x) and the 52097 trainable parameters of the E4 ConvCritic. The split of
% the 18/192 link count (rule: pooled view, universe-bootstrap interval excludes 0; within views,
% permutation p < 0.05; generator level, exact p < 0.05, never met; counting every printed interval
% instead would give 23) (4 positive: T vs pooled-moment GMV and GMV long-only, pooled and
% within-universe views; 14 negative: 8 z_hac, 6 wild-standardised MMD) was counted from
% results/e4/e4_link.csv. Zero risk-free rate: src/synthval/portfolio.py (rf = 0.0 default,
% not overridden in experiments/e4_markowitz.py).
% =============================================================================================

A two-sample test asks whether synthetic windows can be told apart from real ones; a user who
builds portfolios from synthetic data asks whether the portfolios perform well on real data. This
section asks whether the two answers agree. \Cref{sec:mk-theory} states what theory predicts,
\cref{sec:mk-design} describes the experiment E4, \cref{sec:mk-utility} reports the utility of
portfolios trained on synthetic data and \cref{sec:mk-link} relates it to multivariate \WC{} and
MMD tests. The answer is negative: the generator that the tests reject most often loses almost
nothing for the minimum-variance portfolio, the generator with the largest losses is rejected,
after size adjustment, less often than the two \iid{} generators, and the one quantity that
tracks the minimum-variance loss is the distance between the synthetic and the past average
correlation, with a negative sign, because the real correlation fell between \TRAIN{} and
\TEST{}. Three qualifications apply. The evidence comes from one evaluation period, in which
correlations fell, and carries no inference over dates; the portfolios use each generator only
through the unconditional moments of long simulated samples; and within a generator the
multivariate critic statistic is too noisy across synthetic seeds to show clearly whether it
tracks the loss.

% ---------------------------------------------------------------------------------------------
\subsection{What theory predicts}\label{sec:mk-theory}

\Cref{sec:theory-markowitz} summarised what the theory implies for E4; this subsection states the
results. Let $r_t\in\R^N$ be the daily returns of $N$ assets, with real \TRAIN{} sample $R_{\mathrm{tr}}$
($\Ttr\times N$) and real \TEST{} sample $R_{\mathrm{te}}$ ($\Tte\times N$), and let $(\mu,\Sigma)$
be the \TEST{} moments: population moments, or the realised moments
$(\hat\mu_{\mathrm{te}},S_{\mathrm{te}})$, for which every statement below is an identity for the
measured quantities. Portfolios lie in $\mathcal W=\{w\in\R^N:\one'w=1\}$, intersected with
$\{w\ge0\}$ for long-only rules. Given estimates $(\hat\mu,\hat\Sigma)$, the global
minimum-variance (GMV) portfolio is $\wgmv(\hat\Sigma)=\argmin_{w\in\mathcal W}w'\hat\Sigma w$
($=\hat\Sigma^{-1}\one/(\one'\hat\Sigma^{-1}\one)$ without the sign constraint) and the mean-variance
portfolio with risk aversion $\gamma$ is
$\wmv(\hat\mu,\hat\Sigma)=\argmax_{w\in\mathcal W}\{w'\hat\mu-\tfrac\gamma2w'\hat\Sigma w\}$
\citep{markowitz1952portfolio}. Trained on real data (TRTR), these maps are applied to the moments
of $R_{\mathrm{tr}}$ and give $\hat w_{\mathrm{real}}$; trained on synthetic data (TSTR), to the
moments of a synthetic sample $Y^{(b)}$ of length $\Ttr$ from a generator $G$ calibrated on
$R_{\mathrm{tr}}$, and give $\hat w^{(b)}_{\mathrm{syn}}$, $b=1,\dots,B$. With $V(w)=w'\Sigma w$,
$w^*=\wgmv(\Sigma)$ and $V^*=V(w^*)$, the losses are
\begin{gather}
  \mathcal L_{\mathrm{gmv}}(w)=V(w)-V^*,\label{eq:gmv-loss}\\
  \mathcal L_{\mathrm{mv}}(w)=\CE(w^*_{\mathrm{mv}})-\CE(w),\qquad
  \CE(w)=w'\mu-\tfrac\gamma2\,w'\Sigma w,\qquad w^*_{\mathrm{mv}}=\wmv(\mu,\Sigma).\label{eq:ce}
\end{gather}
E4 reports the variance ratio $V(\hat w_{\mathrm{syn}})/V(\hat w_{\mathrm{real}})$, which exceeds
one exactly when $\mathcal L_{\mathrm{gmv}}(\hat w_{\mathrm{syn}})>\mathcal L_{\mathrm{gmv}}(\hat w_{\mathrm{real}})$,
and the certainty-equivalent (CE) loss
$\CE(\hat w_{\mathrm{real}})-\CE(\hat w_{\mathrm{syn}})=\mathcal L_{\mathrm{mv}}(\hat w_{\mathrm{syn}})-\mathcal L_{\mathrm{mv}}(\hat w_{\mathrm{real}})$.
Proofs are in \cref{app:proofs-markowitz}.

\begin{lemma}[Quadratic loss identities]\label{lem:quad-loss}
(a) If $\mathcal W=\{w:\one'w=1\}$, then $V(w)-V^*=(w-w^*)'\Sigma(w-w^*)$ for every $w\in\mathcal W$.
(b) If $\mathcal W$ is closed and convex and $w^*$ minimises $V$ over $\mathcal W$, then
$V(w)-V^*\ge(w-w^*)'\Sigma(w-w^*)$ for every $w\in\mathcal W$.
(c) If $\mathcal W=\{w:\one'w=1\}$, then
$\mathcal L_{\mathrm{mv}}(w)=\tfrac\gamma2(w-w^*_{\mathrm{mv}})'\Sigma(w-w^*_{\mathrm{mv}})$; under
further convex constraints the equality becomes ``$\ge$''.
\end{lemma}

With $\Sigma=S_{\mathrm{te}}$, part (a) decomposes the measured GMV gap exactly into a difference
of squared Mahalanobis distances to the ex-post GMV portfolio,
$\norm{\hat w_{\mathrm{syn}}-w^*}_{S_{\mathrm{te}}}^2-\norm{\hat w_{\mathrm{real}}-w^*}_{S_{\mathrm{te}}}^2$,
where $\norm{x}_S^2=x'Sx$.

\begin{proposition}[Sensitivity of the GMV portfolio to the covariance]\label{prop:gmv-sensitivity}
Let $\mathcal W=\{w:\one'w=1\}$, let $\Sigma$ and $\hat\Sigma$ be symmetric positive definite,
$A=\Sigma^{-1/2}\hat\Sigma\Sigma^{-1/2}$, $\kappa=\lambda_{\max}(A)/\lambda_{\min}(A)=\cond(\Sigma^{-1}\hat\Sigma)$
and $M=\Sigma^{-1}-\Sigma^{-1}\one\one'\Sigma^{-1}/(\one'\Sigma^{-1}\one)$, and let
$\norm{\cdot}_{\mathrm{op}}$ be the spectral norm.
\begin{enumerate}
  \item[(a)] (Scale invariance.) $\wgmv(c\hat\Sigma)=\wgmv(\hat\Sigma)$ for every $c>0$; in
  particular $V(\wgmv(c\Sigma))=V^*$.
  \item[(b)] (Kantorovich bound.) $1\le V(\wgmv(\hat\Sigma))/V^*\le(1+\kappa)^2/(4\kappa)$, and the
  upper bound is sharp.
  \item[(c)] (Quadratic in the relative error.) If $\varepsilon_0=\norm{\Sigma^{-1/2}(\hat\Sigma-\Sigma)\Sigma^{-1/2}}_{\mathrm{op}}$
  is below one, then $V(\wgmv(\hat\Sigma))/V^*-1\le\varepsilon_0^2/(1-\varepsilon_0^2)$.
  \item[(d)] (First order.) As $E\to0$, $\wgmv(\Sigma+E)-w^*=-MEw^*+O(\norm{E}^2)$ and
  $V(\wgmv(\Sigma+E))-V^*=(Ew^*)'M(Ew^*)+O(\norm{E}^3)$, where $M\succeq0$ has kernel spanned by
  $\one$: the quadratic term vanishes if and only if $w^*$ is also the GMV portfolio of $\Sigma+E$.
\end{enumerate}
\end{proposition}

\begin{corollary}[What drives the GMV utility of a generator]\label{cor:gmv-drivers}
Let $\hat\Sigma_{\mathrm{syn}}$ be the sample covariance of a synthetic sample of length $T_s$ from
a stationary ergodic generator with covariance $\Sigma_G$, and $\kappa_G=\cond(\Sigma^{-1}\Sigma_G)$.
\begin{enumerate}
  \item[(i)] As $T_s\to\infty$,
  $V(\wgmv(\hat\Sigma_{\mathrm{syn}}))/V^*\to V(\wgmv(\Sigma_G))/V^*\le(1+\kappa_G)^2/(4\kappa_G)$
  almost surely: in the large-sample limit the GMV utility of a generator depends on it only
  through $\Sigma_G$, and only up to scale.
  \item[(ii)] (Heuristic.) At finite $T_s$,
  $\E\,\mathcal L_{\mathrm{gmv}}\approx\{(\Sigma_G-\Sigma)w^*\}'M\{(\Sigma_G-\Sigma)w^*\}+T_s^{-1}\Tr(M\LRV_G)$,
  with $\LRV_G$ the long-run covariance of $r_t(r_t'w^*)$ under $G$: tails and volatility
  clustering enter only through the second, estimation-noise term.
  \item[(iii)] (Heuristic.) If $\Sigma_G\approx\hat\Sigma_{\mathrm{tr}}$, the \TRAIN{} sample
  covariance, and $T_s=\Ttr$, then
  $\hat\Sigma_{\mathrm{syn}}-\Sigma=(\hat\Sigma_{\mathrm{tr}}-\Sigma)+(\hat\Sigma_{\mathrm{syn}}-\Sigma_G)$
  with two conditionally uncorrelated terms of comparable size: one synthetic sample roughly
  doubles the estimation-noise part of the TRTR loss, and pooling $B$ samples divides the added
  part by $B$. The same holds, with more force, for the mean in mean-variance portfolios.
\end{enumerate}
\end{corollary}

\begin{proposition}[No free lunch from synthetic data]\label{prop:no-free-lunch}
Let $Y=g(R_{\mathrm{tr}},U)$ with $U$ independent of $(R_{\mathrm{tr}},\Theta)$, where $\Theta$ is
any \TEST{} quantity (such as $\Sigma$, $(\mu,\Sigma)$ or $R_{\mathrm{te}}$). Let $\pi$ be a
portfolio rule with values in a convex set $\mathcal W$ and $\ell(w;\Theta)$ be convex in $w$.
Then $\bar w=\E[\pi(Y)\mid R_{\mathrm{tr}}]$ lies in $\mathcal W$, uses the data only through
$R_{\mathrm{tr}}$, and $\ell(\bar w;\Theta)\le\E[\ell(\pi(Y);\Theta)\mid R_{\mathrm{tr}},\Theta]$.
\end{proposition}

Since $V$ and $-\CE$ are convex, no TSTR rule beats, on average over the synthetic randomness,
the best rule based on the real past alone, and averaging the weights over replicas does at least
as well as averaging the losses of replica-wise portfolios. Synthetic data can beat a
\emph{particular} TRTR rule, such as the sample-covariance GMV portfolio, only by acting as a
regulariser, as covariance shrinkage does \citep{ledoit2004well,jagannathan2003risk}. Test
distance and utility can disagree in both directions:

\begin{example}[Rejected by the critic, perfect for GMV]\label{ex:scale-blind}
Let returns be \iid{} $N(0,\Sigma)$ and $G$ \iid{} $N(0,c\Sigma)$ with $c\neq1$. Windows satisfy
$Y\eqd\sqrt c\,X$, so $\W(P,Q)=\abs{\sqrt c-1}\,\E\norm{X}>0$ and a consistent critic test rejects
with probability tending to one (\cref{prop:consistency}); by \cref{prop:gmv-sensitivity}(a) the
GMV loss of $\Sigma_G=c\Sigma$ is zero.
\end{example}

\begin{example}[Nearly invisible in $\W$, large GMV loss]\label{ex:spread}
Let $N=2$, $v_1=(\cos\phi,\sin\phi)'$, $v_2=(-\sin\phi,\cos\phi)'$, $\Sigma=v_1v_1'+\eta^2v_2v_2'$
and $\Sigma_G=v_1v_1'+s\eta^2v_2v_2'$ with $0<s,\eta<1$ and
$\tan\phi=(1-\eta\sqrt s)/(1+\eta\sqrt s)$. Then $\kappa_G=1/s$ and
$V(\wgmv(\Sigma_G))/V^*=(1+s)^2/(4s)$, the Kantorovich bound, which equals $1.5625$ for $s=1/4$
whatever $\eta$. For \iid{} Gaussian windows of length $L$, the quadratic Wasserstein distance
gives $\W\le W_2=\sqrt L\,\eta(1-\sqrt s)\to0$ as $\eta\to0$. The error lies in the low-variance
spread direction $v_2$, which GMV weights heavily (for $\eta=0.1$, $w^*=(-1.48,2.48)'$ and
$\wgmv(\Sigma_G)=(-4.49,5.49)'$) and which $\W$ discounts by $\sqrt{\lambda_2}=\eta$.
\end{example}

\begin{example}[Univariate tests cannot see correlation]\label{ex:corr-blind}
A generator with exactly the univariate window laws of the data and independent assets satisfies
every per-ticker null of E3, so each univariate test rejects it with probability at most $\alpha$;
its covariance is $\diag(\Sigma)$, and its GMV portfolio can be arbitrarily bad.
\end{example}

\begin{example}[Drift errors are undetectable yet decisive]\label{ex:drift}
Let returns be \iid{} Gaussian with known standard deviation $\sigma$, and let the synthetic drift
of one asset be off by $\delta_\mu$. By \cref{prop:oracle-bound}(b), no level-$\alpha$ two-sample
test that uses $\Tte$ real days has power above
$1-\Phi(z_{1-\alpha}-\abs{\delta_\mu}\sqrt{\Tte}/\sigma)$, however many synthetic draws it uses.
For $\abs{\delta_\mu}=\sigma/\sqrt{\Ttr}$, one standard error of the
\TRAIN{} mean, and $\Tte/\Ttr=615/1258$, the bound is $1-\Phi(1.645-0.70)\approx0.17$ at
$\alpha=0.05$. Errors of this size dominate mean-variance utility
\citep{michaud1989markowitz,chopra1993effect}: with $\hat\Sigma=\Sigma$ and the budget constraint
only, \cref{lem:quad-loss}(c) gives the loss $(2\gamma)^{-1}(\hat\mu-\mu)'M(\hat\mu-\mu)$.
\end{example}

For E4 this predicts that GMV losses track the accuracy of the covariance up to scale, not the
rejection rate; that mean-variance losses are dominated by mean estimation noise, which no test
detects and pooling replicas reduces; and that the covariance that matters is that of \TEST{}, so
that fidelity to the past and usefulness for the future differ when the covariance changes.

% ---------------------------------------------------------------------------------------------
\subsection{Design of E4}\label{sec:mk-design}

\paragraph{Universes and generators.}
E4 draws 30 universes of $N=30$ assets, each without replacement and with its own seed, from the
479 names of the complete universe (\cref{sec:data-panel}); universes may share tickers, and
together they contain 400 distinct names. Five
multivariate generators are fitted to the real \TRAIN{} returns of each universe: MV Gaussian
(\iid{}, sample mean and covariance); MV Student-$t$ (\iid{}, location, scatter and degrees of
freedom by ECME; \citealp{liu1995ml}); CCC-GARCH-$t$ (\GARCHt{} margins with variance targeting
and persistence cap 0.995, Gaussian copula with constant correlation estimated from the normal
scores of the standardised residuals; \citealp{bollerslev1990modelling,joe2005asymptotic}); MV
bootstrap (stationary bootstrap of whole \TRAIN{} days with blocks shared by all assets;
\citealp{politis1994stationary}); and factor GARCH-$t$ (single index equal to the equal-weight
return of the same 30 assets, OLS loadings, \GARCHt{} factor and independent \GARCHt{} residuals).
Fitting details are in \cref{app:impl} and diagnostics in \cref{app:mk-generators}.

\paragraph{Utility.}
For each universe and generator, $B=20$ synthetic samples of $\Ttr=1258$ days are simulated from a
stationary start (burn-in 250 days). The portfolios therefore use a generator only through the
moments of long simulated samples, essentially its unconditional mean and covariance: the
conditional covariance forecasts that CCC-GARCH-$t$ and factor GARCH-$t$ could deliver from their
end-of-\TRAIN{} state are deliberately not used, and a variant that would use them (for instance
\TEST{}-length samples started from that state, as in the tests below) was not run. Comparisons
between dynamic and \iid{} generators below concern their unconditional covariances, not their
forecasts. On each sample, and on $R_{\mathrm{tr}}$, the mean is the sample mean and the
covariance the sample covariance or the Ledoit--Wolf estimator with identity target
\citep{ledoit2004well}. Six rules are built: GMV, GMV long-only, long-only mean-variance with
$\gamma\in\{1,5,20\}$ and $1/N$ \citep{demiguel2009optimal}. The weights are held fixed with daily
rebalancing and evaluated on the 615 \TEST{} days through daily simple returns, against the same
rule trained on $R_{\mathrm{tr}}$. The losses are the variance ratio, the Sharpe loss
$\SR(\hat w_{\mathrm{real}})-\SR(\hat w_{\mathrm{syn}})$ (annualised, with a zero risk-free
rate) and the CE loss with the $\gamma$ of the rule ($\gamma=5$ for GMV), in percent per
year ($252\times100$ times the daily value); ratios above one and positive losses mean that TSTR
is worse. Every portfolio is fully invested, so a positive risk-free rate would leave the weights,
the variance ratios and the CE losses unchanged; it would lower the Sharpe ratios and change the
Sharpe losses between portfolios of different volatility. The returns are price returns
(\cref{sec:data-panel}), so the \TRAIN{} means and the levels of the certainty equivalents and
Sharpe ratios omit the dividend yield, and the selection of the universe on membership at the end
of \TEST{} biases the \TRAIN{} means upward. The losses are aggregated over
replicas in three ways: \emph{replica-wise}, the mean loss of the $B$ single-replica portfolios
(the DESIGN measure, which contains the added noise of \cref{cor:gmv-drivers}(iii));
\emph{pooled}, the loss of one portfolio built from the moments of the $B$ stacked replicas, which
isolates misspecification; and \emph{averaged weights} (\cref{prop:no-free-lunch}). The $1/N$
gaps are zero by construction. No inference over dates is provided. All losses are point
estimates on the same 615 \TEST{} days, which every universe shares, so a statement that holds in
all 30 universes describes one realisation seen through 30 overlapping portfolios, not 30
independent confirmations; a test over \TEST{} days, such as the HAC test of Sharpe-ratio
differences of \citet{ledoit2008robust} or its analogue for variances, was not run.

\paragraph{Tests.}
For each universe and generator, one synthetic sample of \TEST{} length is simulated from the end
of \TRAIN{} (GARCH filters continue from their \TRAIN{} state). It and the real \TEST{} sample are
divided by the per-ticker \TRAIN{} standard deviation and cut into non-overlapping windows of
shape $(N,L)=(30,5)$, 123 per sample. The tested object thus differs from the one used for
utility: one \TEST{}-length sample started from the end-of-\TRAIN{} state, against $B$
\TRAIN{}-length samples from a stationary start; for the \iid{} generators the two differ only in
length, for the others also in the starting state. The \WC{} test uses the convolutional critic of
\cref{sec:method-critic}, the time-block split and $\phac$. Each critic is trained on the 60 to 63
windows per sample of its training half, each of dimension $30\times5=150$, far fewer than the
52\,097 trainable parameters of the network. No architecture or training ablation was run for
this multivariate critic (the ablation of \cref{sec:robust} concerns the univariate critic), and
\cref{sec:mk-link} shows that its statistic is noisy across synthetic seeds. MMD uses the wild
bootstrap with a bandwidth fixed per universe (median heuristic of the real \TEST{} windows), so
that all generators of a universe face the same kernel, and they also share the split, critic
initialisation and mini-batches (the critic's input standardisation still uses the pooled
training half, which depends on the generator). Two synthetic seeds were run (DESIGN: ten;
compute budget); rejection shares across universes are reported for the first seed and for $\pmed$
\eqref{eq:pmed}. A reference row uses the last 615 days of real \TRAIN{} (2020-07-24 to
2022-12-30) in place of a synthetic sample. Because $\phac$ and $\pwild$ need not be calibrated
with 123 windows and persistent volatility, 60 placebo replications per generator (DESIGN: 200)
apply the same call to two independent draws of the fitted generator; they give the size and
size-adjusted Monte Carlo p-values (\cref{app:mk-placebo}). E4 uses the absolute tests only: no
anchored relative test, power check or verdict is computed for the multivariate generators.

\paragraph{Link.}
Each (universe, generator) pair is one point, with the seed means of $T$, $\zhac$, $\TMMD$ and the
wild-standardised MMD statistic, and a benchmark that is not a test: the absolute difference
$\abs{\Delta\,\text{avg.\ corr.}}$ between the mean pairwise correlation of the synthetic replicas
and that of real \TRAIN{}. Spearman correlations with the losses are computed pooled over the 150
points, within universes (across the five generators), within generators (across the 30
universes, which removes generator-level offsets) and between the five generator means
(\cref{app:mk-link}). Universes share every \TEST{} date and overlap in tickers, so intervals that
treat them as independent are optimistic; conclusions are conditional on the five generators.
The reproducibility of each test quantity is measured by the Spearman correlation between its
values for the two synthetic seeds (\cref{tab:manual-e4-reliability}).

% ---------------------------------------------------------------------------------------------
\subsection{Utility of portfolios trained on synthetic data}\label{sec:mk-utility}

\begin{table}[tbp]
  \centering
  \caption{Utility losses of portfolios trained on synthetic data relative to the same rule
  trained on real \TRAIN{} (E4). Means over 30 universes of 30 assets of the replica-wise losses
  ($B=20$ synthetic \TRAIN{}-length samples per universe and generator), evaluated on the 615
  \TEST{} days. Var.\ ratio: \TEST{} variance of the TSTR portfolio over that of the TRTR
  portfolio; Sharpe loss: annualised Sharpe ratio of TRTR minus TSTR; CE loss: certainty
  equivalent of TRTR minus TSTR, with the $\gamma$ of the rule ($\gamma=5$ for GMV), percent per
  year. Ratios above one and positive losses: TSTR worse. Sample: sample covariance; LW:
  Ledoit--Wolf (identity target), used for both TSTR and TRTR. Point estimates; no intervals.
  Other aggregations: \cref{tab:e4-utility-agg}.}
  \label{tab:e4-utility}
  \small
  \fittable{e4_utility}
\end{table}

\begin{table}[tbp]
  \centering
  \caption{Correlation structure of real and synthetic samples (E4). Avg.\ corr.: mean pairwise
  return correlation, mean over the 30 universes (sd across universes); synthetic values are
  means over the 20 \TRAIN{}-length replicas. Log-vol corr.: mean pairwise correlation of
  rolling 20-day log-volatilities. $\kappa$: median over universes of
  $\cond(S_{\mathrm{te}}^{-1}\hat\Sigma)$, with $\hat\Sigma$ the real \TRAIN{} covariance or that
  of a synthetic replica (median over replicas first); Kant.\ bound: the corresponding median of
  $(1+\kappa)^2/(4\kappa)$, the bound of \cref{prop:gmv-sensitivity}(b). At cap: GARCH components
  whose persistence reached the cap 0.995, out of all components (CCC-GARCH-$t$: 30 margins per
  universe; factor GARCH-$t$: factor and 30 residuals).}
  \label{tab:e4-correlation}
  \small
  \fittable{e4_correlation}
\end{table}

\Cref{tab:e4-utility} reports the replica-wise losses and \cref{tab:e4-correlation} the
correlation structure; the TRTR benchmark and the other aggregations are in
\cref{app:mk-aggregation}. Ledoit--Wolf moments change the losses little; we discuss sample
moments.

\paragraph{Minimum-variance portfolios.}
Four generators give GMV portfolios whose \TEST{} variance is within 6\% of the TRTR variance
(replica-wise ratios 0.954 for MV Student-$t$, 1.021 for CCC-GARCH-$t$, 1.038 for MV Gaussian,
1.054 for the MV bootstrap); factor GARCH-$t$ has 1.436. With pooled moments the ratios are 0.896,
0.869, 1.004, 0.993 and 1.466 in the same order (\cref{tab:e4-utility-agg}): pooling reduces the
noise but does not help factor GARCH-$t$, whose loss is misspecification. With pooled moments,
CCC-GARCH-$t$ and MV Student-$t$ build better GMV portfolios than the real data they were fitted to
in all 30 universes, and factor GARCH-$t$ has the largest replica-wise ratio of the five generators
in all 30 (counts from the E4 output files; as noted in \cref{sec:mk-design}, these are 30 views
of one \TEST{} period).

\paragraph{The correlation regime shift.}
The mean pairwise correlation fell from 0.419 in \TRAIN{} to 0.268 in \TEST{}, in each of the 30
universes, and the correlation of log-volatilities from 0.603 to 0.302 (\cref{tab:e4-correlation});
\TRAIN{} contains the market-wide crash of 2020. Generators calibrated on \TRAIN{} inherit
the past correlation to different degrees: MV Gaussian exactly (0.419), the MV bootstrap and factor
GARCH-$t$ nearly (0.406, 0.398), MV Student-$t$ less (0.359; plausibly because the $t$ scatter
estimate downweights days with a large Mahalanobis distance, among them the market-wide days of
2020) and CCC-GARCH-$t$ least (0.310), a documented attenuation of the model, whose volatilities
co-move only through the correlation of the shocks (log-volatility correlation 0.127). By
fidelity to the past CCC-GARCH-$t$ is the worst generator; by usefulness for the future, measured
by the pooled-moment GMV variance, it is the best. Lower correlations act like shrinkage of the
covariance towards its diagonal, which is the direction in which the real covariance moved: its
GMV weights are less leveraged (mean gross exposure $\sum_i\abs{w_i}$ of the replica-wise
portfolios 1.706 against 1.986 for TRTR), and its pooled-moment variance ratio is the lowest in 19
of the 30 universes (MV Student-$t$ in 10). This does not contradict \cref{prop:no-free-lunch}:
the attenuated covariance is a regulariser that happened to point in the direction of the regime
change, an advantage that a different future could reverse. Whether the real past alone, treated
in the same spirit, achieves the same gain was not examined: shrinkage towards a
constant-correlation target \citep{ledoit_2004_honey}, a covariance estimated on the most recent
part of \TRAIN{} (such as its last 615 days, the reference row of \cref{tab:e4-tests}) and EWMA or
DCC forecasts \citep{engle2002dynamic} were not run. They would be needed to show that synthetic
data add anything beyond a shrinkage of the correlations.

Factor GARCH-$t$ illustrates the mechanism of \cref{ex:spread}. It matches the average correlation
and the co-movement of volatilities (0.374) better than CCC-GARCH-$t$, yet it has the largest GMV
loss and the largest condition number (31.16, Kantorovich bound 8.299). Two features of its
construction plausibly contribute, and E4 does not separate them (\cref{app:mk-generators}).
First, its covariance is one factor plus a diagonal, so the residual covariances, which carry the
industry structure, are zero. Second, the factor is the equal-weight return of the same 30
assets, so in the data the OLS residuals add up to zero every day and have negative covariances
on average; the simulated residuals are independent, violate this identity and add idiosyncratic
variance to the equal-weight direction. Both errors change the low-variance spread directions on
which GMV portfolios load, while the average correlation is right. An external factor, such as
the equal-weight return of all 479 names or an index, or factors from principal components would
remove the adding-up artefact and leave only the omitted industry structure; this variant was not
run. The Kantorovich bound held in every replica but does not rank the generators: CCC-GARCH-$t$
has a larger $\kappa$ (24.38) than the real \TRAIN{} covariance (18.65) and the lowest pooled
variance ratio.

\paragraph{Mean-variance portfolios.}
For every generator and mean-variance rule the replica-wise TSTR portfolio has a larger \TEST{}
variance and a lower CE than the TRTR portfolio: variance ratios from 1.117 to 1.597, CE losses
from 0.32 to 5.03 percent per year. These losses are mostly noise in the synthetic means: pooling
the replicas lowers the CE loss in 14 of the 15 pairs of generator and rule (for MV Gaussian with
$\gamma=1$ from 3.57 to $-1.36$), as \cref{cor:gmv-drivers}(iii) predicts, and averaged weights do
at least as well as the replica-wise portfolios in every universe, generator and rule, as
\cref{prop:no-free-lunch} requires. Even trained on real data, no mean-variance rule beats $1/N$
in Sharpe ratio (0.72 for $\gamma=5$ against 0.83, with a zero risk-free rate;
\cref{tab:e4-trtr}), as in \citet{demiguel2009optimal}.

% ---------------------------------------------------------------------------------------------
\subsection{Do test results predict utility?}\label{sec:mk-link}

\begin{table}[tbp]
  \centering
  \caption{Rejection shares of the multivariate tests across the 30 universes at $\alpha=0.05$
  (E4): real \TEST{} against one synthetic \TEST{}-length sample per seed, windows
  $(N,L)=(30,5)$, 123 per sample. \WC{}: time-block split, Driscoll--Kraay p-value; MMD: wild
  bootstrap, bandwidth fixed per universe. First: first synthetic seed; $\pmed$:
  $\min\{1,2\med_sp_s\}$ over two seeds. Placebo size: rejection rate when both samples are
  independent draws from the fitted generator (60 replications per generator,
  \cref{tab:e4-placebo}). Size-adjusted: p-values calibrated against the placebo scores of the
  same generator and test. Reference row: the last 615 days of real \TRAIN{} instead of a
  synthetic sample; it has no placebo. Rates only; Wilson 95\% intervals, which treat universes as
  independent, are in \cref{tab:e4-tests-full}.}
  \label{tab:e4-tests}
  \small
  \fittable{manual_e4_tests_main}
\end{table}

\paragraph{What the tests say.}
With nominal p-values and the first seed, the \WC{} test rejects each generator in 0.63 to 1.00
of the universes and MMD in 0.33 to 0.97, and the \WC{} test rejects the recent real past in 0.67
of them (\cref{tab:e4-tests}). The placebo shows that the
nominal $\phac$ is oversized for generators with persistent volatility: 0.27 at 5\% for
CCC-GARCH-$t$ and factor GARCH-$t$ and 0.13 for the MV bootstrap, against 0.05 and 0.07 for the
\iid{} generators (\cref{fig:e4-placebo-size}; the cause is discussed in \cref{app:mk-placebo}).
The reference row cannot be size-adjusted. After size adjustment the \WC{} test still rejects MV
Gaussian in all universes and MV Student-$t$ in 0.87 of them (first seed), but CCC-GARCH-$t$, the
MV bootstrap and factor GARCH-$t$ only in 0.33, 0.33 and 0.40 (0.07, 0.10 and 0.27 with $\pmed$).
MMD also separates the \iid{} generators from the dynamic ones and rejects the latter less often.
Both tests find the generators without volatility clustering easiest to distinguish from real
\TEST{} windows.

\paragraph{What the tests do not say.}
This ordering is unrelated to utility. MV Gaussian is rejected in every universe after size
adjustment (first seed) and has the largest critic statistic of the five generators in 17 of the
30 universes, yet its replica-wise GMV variance ratio lies between 1.007 and 1.060 in every
universe and its pooled-moment ratio is 1.004 on average. No witness analysis was run for the
multivariate critic, so what it detects is not known. If it behaves like the univariate critic of
\cref{sec:emp-witness}, which is close to a realised-volatility statistic, it responds mainly to
the level of window volatility and, for a generator without clustering, to the missing dispersion
of that level across windows. A common volatility level acts as a scale, to which the GMV
portfolio is blind (\cref{ex:scale-blind}), and the dispersion enters the GMV loss only through
estimation noise (\cref{cor:gmv-drivers}(ii)). Factor GARCH-$t$ is rejected less often than
either \iid{} generator after size adjustment, and has the largest GMV loss in every universe.
Mean-variance losses come from mean errors that no test can see (\cref{ex:drift}).

\begin{table}[tbp]
  \centering
  \caption{Reproducibility of the E4 test quantities across synthetic seeds: Spearman
  correlation between the values of the first and the second synthetic seed. Within generator:
  across the 30 universes of one generator (the variation used by the within-generator view of
  \cref{tab:e4-link}), and the mean of the five values. Pooled: over the 150 (universe,
  generator) points. Reference: last 615 days of real \TRAIN{}, the same data in both seeds, so
  that only the randomness of the test differs (split, initialisation and mini-batches of the
  critic; wild-bootstrap multipliers of MMD). \WC{}: time-block split, critic trained on 60 to 63
  windows of dimension 150 per sample; MMD: bandwidth fixed per universe.}
  \label{tab:manual-e4-reliability}
  \small
  \fittable{manual_e4_reliability}
\end{table}

\paragraph{How reproducible the test quantities are.}
A correlation between a test quantity and a loss across universes is informative only if the
quantity itself is reproducible across universes. \Cref{tab:manual-e4-reliability} shows that,
within a generator, the critic statistic is not. Across the 30 universes of a generator, the
values of $T$ for the two synthetic seeds have Spearman correlation between $-0.03$ and 0.30 (mean
0.14), and those of $\zhac$ between 0.10 and 0.51. For MMD$^2_u$ the values are 0.22 to 0.83 (mean
0.44). With the data held fixed (reference row), $T$ still has seed-to-seed correlation 0.28,
while MMD$^2_u$ is identical in both seeds (1.00): the randomness of the test alone, a split and a
network trained on few windows, makes $T$ barely reproducible. The link uses the mean of the two
seeds, which halves the variance of the seed noise but leaves it large relative to the variation
of $T$ across the universes of a generator (\cref{app:mk-link} quantifies the attenuation).

\begin{table}[tbp]
  \centering
  \caption{Spearman correlations between test quantities and utility losses (E4; sample
  moments). Panel A: replica-wise portfolios (the DESIGN measure); panel B: portfolios from the
  moments of the pooled replicas (misspecification only). Test quantities are seed means (two
  seeds); $\abs{\Delta\,\text{avg.\ corr.}}$: absolute difference between the synthetic and the
  real \TRAIN{} mean pairwise correlation (not a test). Pooled: over the 150 (universe,
  generator) points, universe-bootstrap 95\% percentile interval (2000 draws). Within-universe and
  within-generator: mean Spearman correlation across the 5 generators of a universe or the 30
  universes of a generator, two-sided permutation p-values (2000 permutations). Generator level:
  Spearman correlation of the 5 generator means, exact two-sided p-value over the 120 orderings.
  Universes share dates, so intervals are optimistic. All losses:
  \cref{tab:e4-link-full,tab:e4-link-pooled}.}
  \label{tab:e4-link}
  \small
  \fittable{manual_e4_link_main}
\end{table}

\paragraph{Correlations between test quantities and losses.}
\Cref{tab:e4-link} and \cref{fig:e4-utility-vs-distance} make the comparison systematic. For the
replica-wise GMV variance ratio, the \WC{} statistic $T$ has a pooled correlation of 0.10
[$-0.04$, 0.25], a within-generator correlation of 0.01 (permutation $p=0.853$) and a
generator-level correlation of 0.50 (exact $p=0.450$, $n=5$). Given the low reproducibility of
$T$, the within-generator value is weak evidence of no association: it is compatible with seed
noise alone and does not exclude a moderate association between the loss and a noise-free critic
distance. The MMD quantities, whose statistics are on average more reproducible, are no better:
both have within-generator correlations with the GMV variance ratio of $-0.08$ for replica-wise
and $-0.02$ for pooled-moment portfolios, and the wild-standardised MMD statistic is even
correlated negatively with the loss ($-0.21$ pooled). For the CE loss with $\gamma=5$ no test
quantity has a pooled correlation larger than 0.16 in absolute value. With pooled-moment
portfolios (panel B), $T$ orders the generators roughly as their GMV losses (pooled 0.30 [0.19,
0.42], within universe 0.39 with $p<0.001$, generator level 0.70 with exact $p=0.233$). Its
within-generator correlation, $-0.01$ ($p=0.895$), is weak evidence for the same reason, and the
between-generator association rests on five generators, of which only factor GARCH-$t$ has a
large loss. The evidence that test results do not predict utility thus comes from the comparison
of generators (MV Gaussian rejected everywhere yet with GMV variance ratios close to one, 1.038
replica-wise and 1.004 pooled; factor GARCH-$t$ rejected less often yet losing most) and from MMD, not from the within-generator correlations of the critic.
Over the full tables, 18 of the 192 combinations of test quantity, loss, aggregation and view have
a pooled interval excluding zero or a permutation p-value below 0.05. If all associations were null and
the comparisons independent, about eight would be expected by chance (\cref{app:mk-link}); the
comparisons share their points, so the count is descriptive. The four positive ones are all $T$
against pooled-moment GMV losses, in the pooled and within-universe views, which mostly rank the
generators; the other 14 are negative and all involve a studentised statistic ($\zhac$ or the
wild-standardised MMD statistic).

\begin{figure}[tbp]
  \centering
  \inputfigure{e4_utility_vs_distance}
  \caption{Utility losses against test statistics (E4), one point per (universe, generator), 30
  universes $\times$ 5 generators. Top: GMV variance ratio (TSTR over TRTR); bottom: mean-variance
  ($\gamma=5$) CE loss in percent per year; replica-wise portfolios, sample moments. Left: \WC{}
  statistic $T$; right: $\TMMD$; seed means over two seeds. Insets: pooled Spearman $\rho$ and
  within-generator $\bar\rho$ with universe-bootstrap 95\% percentile intervals. Dashed lines: no
  loss.}
  \label{fig:e4-utility-vs-distance}
\end{figure}

\paragraph{The correlation shift.}
The one quantity with a consistent association is not a test statistic.
$\abs{\Delta\,\text{avg.\ corr.}}$ is correlated with the GMV variance ratio at $-0.29$ [$-0.44$,
$-0.14$] pooled and $-0.36$ ($p<0.001$) within generator for replica-wise portfolios, and at
$-0.64$ [$-0.70$, $-0.57$] and $-0.41$ ($p<0.001$) for pooled-moment portfolios. The further a
generator's correlation is from the past, the smaller its GMV loss. The synthetic correlation lies
below the \TRAIN{} one in 122 of the 150 points (in all 30 universes for CCC-GARCH-$t$ and MV
Student-$t$), and the real correlation fell, so departing from the past meant moving towards the
future; a generator that reproduced the past correlation perfectly was penalised for it. The
association holds within generators, so it does not come from generator-level offsets, and it is
specific to GMV: for the CE loss with $\gamma=5$ it is $-0.03$ pooled. It is also a property of
this evaluation period, in which the correlation fell in each of the 30 universes; under rising
correlations the sign would plausibly reverse. The alternative split of \cref{sec:robust}, with
\TEST{} starting in 2022, was not run for E4.

\paragraph{Summary.}
Closeness in distribution does not imply usefulness for portfolio construction. The GMV portfolio
depends only on the covariance matrix up to scale, so features that a window test detects can be
irrelevant to it; mean-variance utility is driven by mean errors that no test sees at these sample
sizes; and when cross-sectional dependence shifts between calibration and evaluation, the
generator closest to the past is not the most useful. These findings are conditional on the
design. They come from one evaluation period, in which correlations fell, as point estimates
without inference over dates, and from five generators; in a period of rising correlations the
negative association between the correlation distance and the GMV loss could reverse. The tests
examine a \TEST{}-length sample started from the end-of-\TRAIN{} state, whereas the portfolios
use each generator only through the unconditional moments of long stationary samples, so the
conditional forecasts of the GARCH-type generators are not exploited. Within a generator, the multivariate critic statistic is so noisy
across seeds that its correlation with the loss is weak evidence either way, and the evidence
rests on the comparison of generators and on MMD. None of this is a failure of the test, which
answers, with an approximate size adjustment, whether synthetic windows are distinguishable from
real \TEST{} windows. Whether synthetic data are good for building portfolios has to be checked
on that task, by comparing TSTR with TRTR portfolios on held-out real data.

\section{Robustness}\label{sec:robust}
% =============================================================================================
% E5 (experiments/e5_robustness.py; parts grid, split, simpower as run; gridsize, gridpower,
% gridrealvol not run) and the post-hoc architecture ablation (experiments/e5_architecture.py ->
% results/e5arch/ -> tables e5_arch_sim.tex, e5_arch_real.tex).
% Main text: tab:manual-e5-sensitivity, fig:e5-sensitivity, tab:manual-e5-split (condensed tables
% written by paper/scripts/e5_tables.py; the generated e5_sensitivity.tex and e5_split.tex are too
% wide for the main text and are in app:robust as tab:e5-sensitivity, tab:e5-split).
% Numbers: paper/tables/e5_*.tex, manual_e5_*.tex, results/e5/*.csv, results/e5arch/*.csv,
% results/e3/e3_{power,realvol}_{raw,summary}.csv. Derived counts (orderings, McNemar tests):
% recomputed from the raw CSVs. Paired alternative/boundary outcomes of the rule in D1/D2 at both
% splits (minimum boundary pmed, alternatives below it): paper/scripts/e5_split_discrimination.py
% -> paper/tables/manual_e5_split_paired.tex (input in app:robust as tab:manual-e5-split-paired).
% Architecture pilot: docs/DESIGN.md change log 2026-11-16 ("four E2 designs at n = 500"; the pilot
% itself is not in the repository).
% =============================================================================================

E5 reruns the real-data decision rule of \cref{sec:method-decision} under one-at-a-time changes of
its tuning choices and under an alternative \TRAIN/\TEST{} split, and examines the optimiser
settings of the critic on simulated data; a post-hoc ablation replaces the critic network. Every
real-data run uses the 100 tickers, fitted models, five synthetic seeds and test randomness of
E3, so the default configuration reproduces E3 exactly: its 50 p-values (absolute \WC{} and
anchored tests, five generators, five seeds) are identical to those of E3. Three planned parts
were not run for lack of compute (\cref{app:robust}): the size of the absolute test under each
configuration, and the power and boundary size of the rule under each configuration in designs
D1 and D2. A non-default configuration therefore has no power estimate; its verdict can be \BAD{}
or \INCONCL{} but not \GOOD, and \INCONCL{} there only means that the anchored test did not
reject. The size of the default pipeline is measured in E1 (\cref{tab:e1-size-sim}): at nominal
5\% the absolute \WC{} and anchored tests reject a true null at 0.085 and 0.080 when the synthetic
panel is a joint draw, but at 0.270 and 0.155 when the synthetic tickers are independent, as they
are for the univariate generators examined here.

\paragraph{Sensitivity grid.}
The grid changes one setting at a time (\cref{tab:manual-e5-sensitivity,fig:e5-sensitivity}): the
window length $L\in\{5,10,60\}$ (default 20), the width and depth of the critic ($32\times2$ and
$256\times4$; default $128\times3$), the number of training steps (500 and 3000; default 1000),
the learning rate ($10^{-4}$; default $10^{-3}$), the gradient-penalty weight ($\lgp=1$; default
10), the number of time blocks ($K=4$; default 10), the training fraction (0.3 and 0.7; default
0.5) and the HAC lag rule (the rule of \citet{newey1994automatic} instead of the plug-in of
\citet{andrews1991heteroskedasticity}). Each change applies to the absolute and to the anchored
test. Three results hold in every configuration. First, the absolute \WC{} test rejects at 5\% in
69 of the 70 configuration--generator cells; the exception is the stationary bootstrap with $K=4$
($\pmed=0.095$). Second, the anchored test never rejects \GARCHt{}: its $\pmed$ ranges from 0.100
($L=5$) to 1.000. Third, the extremes of the ranking of the likelihood models do not change. In
every configuration GBM has the smallest anchored $\pmed$ and \GARCHt{} the largest (shared with
Heston at 1.000 when the training fraction is 0.3), and the median absolute statistic decreases
from GBM to Merton, Heston and \GARCHt{}. The anchored order of Merton and Heston is not stable:
Heston has the smaller $\pmed$ with $L=60$ (0.070 against 0.110) and with the $256\times4$ network
(0.028 against 0.029). Both statistics thus rank GBM below Merton and Merton below \GARCHt, the
order of the predictive log-score (\cref{tab:e3-agreement}).

The \BAD{} verdicts are less stable. The anchored test rejects GBM in 13 of the 14
configurations, Merton in 11 and Heston in 9. All nine changes against the default go from
rejection to non-rejection; no configuration rejects a generator that the default does not
reject. They occur where the anchored test has fewer evaluation dates or a weaker critic: longer
windows ($L=60$ leaves 10 windows per ticker in \TEST; Merton $\pmed=0.110$, Heston 0.070), a
30\% training share (GBM 0.170, Merton 0.671, Heston 1.000), four time blocks (Merton 0.236,
Heston 0.379) and, for Heston only, the small network (0.102) and the lower learning rate
(0.079). Shorter windows, the larger network, fewer or more steps, $\lgp=1$, a 70\% training
share and the Newey--West rule change no decision, but the narrowest decision of the grid is
among them: with 500 steps Heston has $\pmed=0.0499$, printed as 0.050 in
\cref{tab:manual-e5-sensitivity} and just below the level. At the default, Merton ($\pmed=0.005$) and
Heston (0.017) are closer to the level than GBM ($<0.001$). \Cref{fig:e5-sensitivity} also shows
why the seeds are combined: the single-seed p-values of one generator in one configuration
spread over several orders of magnitude (\cref{rem:combine-seeds}).

\begin{table}[tbp]
  \centering
  \caption{Sensitivity of the E3 decisions to one-at-a-time changes of the pipeline (E5 grid;
    100 tickers, \TEST{} 2023-01-03 to 2025-06-16, five synthetic seeds, $\alpha=0.05$).}
  \label{tab:manual-e5-sensitivity}
  \small
  \fittable{manual_e5_sensitivity}
  \tablenote{Generator columns: $\pmed=\min\{1,2\med_sp_s\}$ over the five seeds of the
    Driscoll--Kraay p-values (time-block split) of the anchored relative test, followed by the
    verdict (B = \BAD, I = \INCONCL). $^\dagger$: the decision at 5\% differs from the default row.
    $^\ddagger$: relative test invalid for the stationary bootstrap (no verdict). Last column:
    number of the five generators that the absolute \WC{} test ($\pmed$ over the same seeds)
    rejects at 5\%. The power check was not rerun for any configuration, so a non-default verdict
    is \BAD{} or \INCONCL{}; the default verdicts coincide with \cref{tab:e3-verdicts}, whose power
    comes from \cref{tab:e3-power}. All p-values: \cref{tab:e5-sensitivity} in \cref{app:robust}.}
\end{table}

\begin{figure}[tbp]
  \centering
  \inputfigure{e5_sensitivity}
  \caption{p-values of the E5 grid (100 tickers, \TEST{} 2023-01-03 to 2025-06-16, five synthetic
    seeds). Left: anchored relative test; right: absolute \WC{} test; one row per configuration.
    Open markers: the five single-seed Driscoll--Kraay p-values; filled markers:
    $\pmed=\min\{1,2\med_sp_s\}$. Dashed line: $\alpha=0.05$. p-values below $10^{-4}$ are drawn at
    the left edge. The stationary bootstrap is faint in the left panel because its relative test
    is invalid.}
  \label{fig:e5-sensitivity}
\end{figure}

\paragraph{Alternative split.}
With \TRAIN{} returns from 2018-01-03 to 2021-12-31 and \TEST{} 2022-01-03 to 2025-06-16 (50 and 43 windows
per ticker), the five generators are refitted and the full E3 battery is rerun with five seeds;
the power and boundary size of the rule are estimated at the sizes of this split with 25
replications per design (\cref{tab:manual-e5-split}). The absolute tests reject fewer
generators, although \TEST{} is longer than at the main split. The \WC{} test rejects only GBM
($\pmed<0.001$; Merton 0.055, Heston 0.076, \GARCHt{} 1.000), against all five generators at the
main split, and no valid window test rejects \GARCHt{} or the bootstrap. The anchored test rejects
no generator (GBM 0.091, Merton 0.704, \GARCHt{} 0.516, Heston 0.966). The power of the rule is
1.00 [0.87, 1.00] in D1 but 0.60 [0.41, 0.77] in D2, below $\bmin$, so every verdict is
\INCONCL{}. None of the three \BAD{} verdicts of the main split is reproduced.

The boundary sizes show that \INCONCL{} is the right outcome, for a reason that the rule does not
encode. In D1 the rule never rejects at the boundary (0.00 [0.00, 0.14], Clopper--Pearson). In D2
its boundary size is 0.60 [0.39, 0.79]: a generator that reproduces the pooled \TRAIN{} window law
exactly is declared worse than history in 15 of 25 replications. These are the same 15
replications in which the alternative $G^*=1.2\,G^0$ is rejected, and the alternative is rejected
in no other replication (\cref{app:robust}). At the nominal threshold the D2 power of the rule is
therefore not interpretable. The critic still separates the two laws: in each of the 15 rejecting
replications the alternative has $\pmed\le7.8\times10^{-6}$, below the smallest boundary $\pmed$ of
the 25 replications ($3.8\times10^{-5}$), so a critical value just below that smallest boundary
value, itself estimated from these replications, would reject the boundary in none of the 25 and
the alternative in the same 15 (\cref{tab:manual-e5-split-paired} in \cref{app:robust}). What fails at this split is the calibration of the nominal
threshold, not discrimination; whether the nominal test fires depends on the draw of the
replication. The median anchored statistic at the D2 boundary is 0.231, against 0.012 at the main
split, so the anchored target at the pooled-\TRAIN{} boundary is far from zero, the case that
\cref{rem:nonstat-past} anticipates for a heterogeneous past. Three consequences follow. First, a
\BAD{} verdict at this split, read at the nominal threshold, would have said nothing about the
generator. Second, the rule \eqref{eq:verdict} checks power but not boundary size, and
its \BAD{} branch can be interpreted only where the estimated boundary size is close to $\alpha$,
as at the main split (D1 0.00 [0.00, 0.04], D2 0.08 [0.04, 0.15]; \cref{tab:e3-power}). A rule
that turns \BAD{} into \INCONCL{} whenever the Clopper--Pearson interval of the boundary size lies
above $\alpha$ would encode this; it is not the rule that produced our tables. Third, under
non-stationarity a verdict is a statement about one pair of periods: the past that serves as the
benchmark changes with the split, and so do the verdicts and the absolute rejections.

\begin{table}[tbp]
  \centering
  \caption{Alternative split (E5): \TRAIN{} 2018-01-03 to 2021-12-31, \TEST{} 2022-01-03 to
    2025-06-16, same 100 tickers, generators refitted; five synthetic seeds; rule power and
    boundary size from 25 replications per design; $\alpha=0.05$.}
  \label{tab:manual-e5-split}
  \small
  \fittable{manual_e5_split}
  \tablenote{Entries are $\pmed$ over the five seeds of the Driscoll--Kraay p-values (time-block
    split) of the absolute \WC{} test and of the anchored relative test; bold: absolute rejection
    at 5\%. $^\ddagger$: relative test invalid for the stationary bootstrap (verdict n/a). Rule
    power (synthetic volatility $\times1.2$) and boundary size: share of replications (number in
    parentheses) in which $\pmed\le0.05$ over five seeds, at the sizes of each split, with Wilson
    (power) and Clopper--Pearson (boundary size) 95\% intervals; main split as in
    \cref{tab:e3-power}. All tests of the battery: \cref{tab:e5-split} in \cref{app:robust}.}
\end{table}

\paragraph{Optimiser settings.}
On simulated independent windows with exact permutation p-values ($n=500$ windows per sample, 50
replications per cell; \cref{tab:e5-simpower}), the default critic has power 0.40 [0.28, 0.54]
against a 10\% increase in the volatility of Student-$t_5$ returns and 0.30 [0.19, 0.44] against
GARCH persistence 0.90 instead of 0.98. With learning rate $10^{-4}$ and the default 1000 steps,
the power against the volatility shift rises to 0.78 [0.65, 0.87], while against clustering it
is 0.34 [0.22, 0.48], not detectably different. Power is not monotone in the training budget: at
rate $10^{-4}$ it is highest at 1000 steps, and 500 steps under-train the critic. The level does
not depend on any of these settings (\cref{prop:exact-level}); the power does. The defaults were
fixed before any E2 run and were not re-tuned on these results, so the \WC{} power of
\cref{sec:sim} is that of the default and may understate what a differently trained critic
achieves against volatility-level deviations. On real data the lower rate did not reject more;
it turned Heston from \BAD{} to \INCONCL{} ($\pmed=0.079$). Its rule power in D1 and D2 was not
estimated, so this change can be attributed neither to size nor to power, and whether the lower
rate would raise the D2 power of the rule is unknown.

\paragraph{Critic architecture.}
The network matters at least as much as the optimiser, but the evidence for the alternative
network is weaker than for the default. A time-pooled convolutional critic, which averages one
learned function of nine-day patches over the window (\cref{app:robust}), was chosen after the
real-data results were known and after a pilot on four E2 designs at $n=500$ had shown it to be
more powerful than the perceptron; the pilot is not part of the reported experiments. It is
therefore a post-hoc ablation, and the pre-specified verdicts remain those of the perceptron. In
replications paired with those of E2 and E2-dep (\cref{tab:e5-arch}), it keeps the level with
independent windows (0.06 and 0.05 in 200 null replications) and raises the power in every
alternative design, significantly in all but the block design of E2-dep: from 0.10 to 0.98
against lighter tails ($t_{10}$ against $t_4$), from 0.10 to 0.92 when every synthetic window is
GBM instead of Merton, and from 0.36 to 0.88 against persistence 0.90. Because the network was
selected on E2 designs of this kind, possibly on some of these very cells, the E2 gains are
optimistic. The two E2-dep cells ($n=2000$), outside the pilot as recorded, are the closest to
held-out evidence: 0.90 and 0.96 against 0.46 and 0.86 for the perceptron (exact McNemar
$p<0.001$ and $p=0.18$). In the designs on which the pooled critic was selected, the weakness of
the \WC{} test against marginal deviations in \cref{sec:sim} is thus largely a property of the
multilayer perceptron. The size of the pooled critic under the dependent protocol (E1(a)) was not
measured. On the S\&P~500 panel it gives the same five decisions (\cref{tab:e5-arch-real}). Its
rule power and boundary size were estimated only in the first 20 replications of D1 and D2, where
they equal those of the perceptron on the same replications: power 1.00 and 0.45 [0.26, 0.66],
boundary size 0.00 and, in D2, 0.20 [0.06, 0.44] (Clopper--Pearson), against 0.08 [0.04, 0.15]
over all 100 replications of the perceptron. These replications are not representative of E3 at
the D2 boundary and cannot separate the two networks. Whether a more powerful critic, or the
lower learning rate, would lift the D2 power above $\bmin$ is therefore open.

\paragraph{What is robust.}
At the main split, across the thirteen changes of the grid, the absolute \WC{} test rejects every
parametric generator (and the bootstrap in every configuration but the one with four time blocks,
$\pmed=0.095$), the anchored test does not reject \GARCHt{}, and both statistics rank GBM below
Merton and Merton below \GARCHt. At the alternative split the absolute \WC{} test rejects only GBM
and the anchored test rejects no generator. \GARCHt{} is therefore never \BAD{}; it is never
\GOOD{} either, because its D2 power is below $\bmin$ at both splits and was not estimated in the
grid. GBM is \BAD{} in all but one configuration of the grid. Four conclusions are not robust. The
\BAD{} verdicts of Merton and Heston disappear under three and five of the thirteen changes, and
that of Heston with 500 steps rests on $\pmed=0.0499$. No \BAD{} verdict survives the change of
split, and at the alternative split the nominal threshold of the rule is not calibrated at the
boundary in the design with the real volatility path, although the critic still separates the
alternative from the boundary. The absolute rejections of Merton, Heston, \GARCHt{} and the
bootstrap also disappear at that split. The power of the \WC{} test depends on the optimiser
settings, not monotonically, and on the network. The rule power on real data, on which a \GOOD{}
verdict depends, has been estimated only for the default optimiser, with the perceptron and, in
20 replications, with the pooled critic. The exact level for independent windows holds for every
setting. On real data, calibration at the relative boundary is a
property of the chosen periods and has to be checked for each split, as the boundary size does.
Whether a generator is good enough for a task such as portfolio construction has to be checked
on that task (\cref{sec:markowitz}).

\section{Conclusion}\label{sec:concl}
% =============================================================================================
% Owner: front writer. Plan: OUTLINE.md, Section 10. Every number is taken from a generated or
% condensed table that the result sections input: tab:e3-power (manual_e3_power_main),
% tab:e3-pooled-tests, tab:e3-verdicts, tab:e1-size-sim (manual_e1_size_sim_main),
% tab:e2-relative-boundary, tab:e2-power, tab:e2dep-power, tab:e2-llr, tab:e3-perticker-null,
% tab:e4-correlation, tab:e4-utility, tab:e4-utility-agg, tab:e4-tests, tab:e4-link,
% tab:manual-e5-sensitivity, tab:manual-e5-split, tab:e5-simpower, tab:e5-arch (e5_arch_sim),
% tab:e5-arch-real (manual_e5_arch_real), tab:manual-theory-calibration (size-adjusted power),
% tab:manual-e1-sim-stat (correlation of T with its HAC s.e.), or from the text of
% sec:emp-witness (0.65) and sec:emp-likelihood (1/6). The configuration counts of the E5 grid
% (13, 11 and 9 of 14 BAD; 69 of 70 absolute rejections) are counted from
% manual_e5_sensitivity.tex. Two derived counts were recomputed with uv run python from the raw
% results (revision v1, 2026-11-16): (i) alternative split, design D2
% (results/e5/e5_split_power_raw.csv, pmed = min(1, 2 median) over the five seeds per
% replication): the alternative and the boundary are rejected in the same 15 of 25 replications;
% every rejected alternative has pmed <= 7.8e-6, below the smallest of the 25 boundary pmed
% (3.78e-5). (ii) E4 (results/e4/e4_tests.csv, test == critic): Spearman
% correlation of T between the two test seeds across the 30 universes of a generator, -0.03
% (MV bootstrap) to 0.30 (CCC-GARCH-t; 0.3046), as in tab:manual-e4-reliability. The architecture ablation (results/e5arch) is post hoc;
% the pre-specified verdicts are those of the MLP critic. The CNN was chosen after a pilot on four
% E2 designs at n = 500 (docs/DESIGN.md, change log 2026-11-16). Integration pass (2026-11-16):
% tab:e1-sim-se (0.480 vs 0.060), tab:e3-rv-baseline (0.69-0.94), tab:e4-utility-agg (1.038,
% 1.004); classifier checkpoints from app:impl-nets.
% =============================================================================================

The question in the title has a qualified answer. For the five univariate generators and the
S\&P~500 panel we studied, synthetic data from GBM, Merton and Heston calibrated on 2018--2022
are not good enough: their windows are detectably less like the 2023--2025 market than the
2018--2022 history they were fitted to ($\pmed<0.001$, 0.005 and 0.017 for the anchored test;
\cref{tab:e3-verdicts}). Each verdict concerns a model together with its calibration rule.
Fitted to a period that contains the 2020 crash, these generators reproduce its volatility level,
whereas the median realised volatility of the real evaluation windows is 0.65 times the \TRAIN{}
standard deviation, and the witness of the critic shows that this excess volatility is mainly
what the test detects (\cref{sec:emp-witness}). Heston's rejection is
borderline against the boundary replications of the realistic design (\cref{sec:emp-verdicts}),
and it concerns Heston with the coefficient of variation of the variance bounded by
$\sqrt{1.5}$, the bound that keeps our daily discretisation free of zero-variance days
(\cref{sec:data-generators}). \GARCHt{}, the best of the three likelihood models by every
likelihood criterion, is not detectably worse than history ($\pmed=0.432$), but the protocol
cannot call it \GOOD: under the design that keeps the real volatility path of the calibration
period, the rule detects a volatility 20\% too high only with probability 0.45 [0.36, 0.55]
(\cref{tab:e3-power}). The stationary bootstrap, which replays the past, receives no verdict:
with the split we ran, its windows duplicate windows the anchored critic was trained on, and the
split by source window that removed this leakage in simulation (\cref{tab:e2-relative-boundary})
was not applied to it. For portfolio construction the answer comes from the task, not from the
test, and here from one evaluation period in which correlations fell. The generators that gave
the best minimum-variance portfolios, CCC-GARCH-$t$ and MV Student-$t$, were those whose
correlations departed most from a past whose average correlation (0.419) was well above that of
the evaluation period (0.268; \cref{tab:e4-correlation}); under rising correlations this
association could reverse. The generator most often rejected, MV Gaussian, gave minimum-variance
portfolios with \TEST{} variance ratios close to one (1.038 replica-wise, 1.004 with pooled
moments; \cref{tab:e4-utility-agg}). Within a generator the critic statistic was too noisy across
synthetic seeds for its near-zero correlation with the loss (\cref{tab:e4-link}) to be
informative: across the 30 universes of a generator its Spearman correlation between the two test
seeds lies between $-0.03$ and 0.30 (\cref{tab:manual-e4-reliability}).

What the protocol guarantees is limited but explicit. With independent windows the permutation
p-value of the \WC{} test has exact level for every critic, and the permutation version of the
anchored relative test is exact at its boundary $\Law(G)=\Law(\Rpast)$ for every critic and every
law of the new data (\cref{prop:exact-level,prop:anchored}); the simulations confirm both
(\cref{sec:sim-size,sec:sim-relative}; the anchored test rejected at 0.065 and 0.055 at nominal
5\% at its boundary). On dependent panels the time-block split removes the common-date bias of
ticker splits (\cref{prop:ticker-bias}), whose naive version rejected a true null in 0.990 to
1.000 of the placebo replications, though not the bias that a persistent factor can carry into
adjacent blocks (\cref{cor:block-bias}). The Driscoll--Kraay p-value is asymptotically valid for a fixed critic
whose evaluation data are independent of its training data (\cref{prop:hac}), a premise that our
split without a gap does not ensure; in the placebo it rejected at 0.085 when the synthetic panel
shared the factor structure of the real one (\cref{tab:e1-size-sim}), but not when the synthetic
tickers were independent (\cref{sec:concl-limitations}). The verdicts use these Driscoll--Kraay
p-values, so a \BAD{} verdict is a rejection whose error rate at the relative boundary is
estimated, not assumed (0.08 [0.04, 0.15] in the realistic design); a \GOOD{} verdict would rest
on a power premise that cannot be verified (\cref{prop:verdict}). Within this protocol the critic
is one calibrated and interpretable tool, not the most powerful one. The pooled
Kolmogorov--Smirnov test wins against deviations of the marginal law, and kernel tests are at
least as powerful against changes of GARCH persistence at matched variance
(\cref{sec:sim-power}). The case for the critic is that it is one function: it can be evaluated
window by window, it detected dependence-only deviations that a marginal test cannot see
(\cref{fig:e2dep-power}), it ordered the three likelihood models as the predictive log-score
does (\cref{tab:e3-agreement}), an agreement that a random ordering reaches with probability
1/6, it extends to multivariate windows, and it supports a relative null that is exact at its
boundary.

\subsection{Limitations}\label{sec:concl-limitations}

\paragraph{Verdicts depend on the split.}
Under non-stationarity a verdict is a statement about one pair of periods. With the calibration
period ending in 2021 instead of 2022, every parametric generator is \INCONCL{}, and in the realistic design
the rule rejects the boundary generator in 0.60 [0.39, 0.79] of replications, as often as it
rejects the 20\% alternative (\cref{tab:manual-e5-split}). In the 15 replications in which both
are rejected, however, the alternative is far more significant: its $\pmed$ lies below every one
of the 25 boundary values, and the size-adjusted power of \cref{def:power-check} is still 0.60 at
a threshold where the boundary size is 0.04 (\cref{tab:manual-theory-calibration}). The failure
at that split is one of calibration rather than of discrimination, although the size-adjusted
power also stays below $\bmin$. The \BAD{}
verdicts of Merton and Heston also disappear under some one-at-a-time changes of the pipeline,
in 3 and 5 of the 13 changes (\cref{tab:manual-e5-sensitivity}). The decision rule
\eqref{eq:verdict} checks power but not boundary size, so it would have issued uninterpretable
\BAD{} verdicts at the alternative split had the anchored test rejected (\cref{sec:robust}). A
calibration of the generators on a recent part of the calibration period, which would separate
the model from the volatility level of its calibration window, was not run.

\paragraph{Driscoll--Kraay inference can be anti-conservative.}
The normal-reference Driscoll--Kraay p-value relies on few evaluation dates (about 15 per half in
E3). When the synthetic tickers are drawn independently, as for the univariate generators, the
learned tests over-reject in the placebo even with the time-block split (0.270 for the absolute
\WC{} test, 0.155 for the anchored test; \cref{tab:e1-size-sim}), and fixed-$b$ references do not
repair it. In the multivariate study, with 123 windows per sample and persistent volatility, the
\WC{} test rejected a true null at 0.27 for CCC-GARCH-$t$ and factor GARCH-$t$
(\cref{tab:e4-tests}); per ticker, at $L=5$, it rejected at 0.12 (\cref{tab:e3-perticker-null}).
We therefore size-adjusted or calibrated these p-values against simulated nulls, which is only as
good as the simulated null. The interleaved blocks without a gap leave dependence between
training and evaluation windows that the theory does not control (\cref{lem:decoupling}).
\Cref{cor:block-bias,rem:small-T} give a model-based and heuristic explanation of the oversize. A
critic trained on dates adjacent to the evaluation blocks carries the persistent realised factor
path into them, and with independent synthetic tickers the conditional bias of $T$ and its
Driscoll--Kraay standard error depend on the same real path. In the placebo the mean of $T$ is
nearly the same in the two cells, but with independent synthetic tickers $T$ and its standard
error are negatively correlated (Pearson $-0.64$, against 0.22 when the synthetic panel has a
factor of its own; \cref{tab:manual-e1-sim-stat}), so positive values of $T$ come with small standard errors and
$\zhac$ is shifted: with the \WC{} test, 0.480 of the replications with a standard error at or
below the median reject, against 0.060 of the others (\cref{tab:e1-sim-se}). An embargo between
training and evaluation blocks, as in $hv$-block and purged cross-validation
\citep{racine_2000_consistent,lopezdeprado_2018_advances}, a contiguous split or a critic trained
on an independent copy of the panel would test this explanation; none was run.

\paragraph{Power under realistic non-stationarity.}
In a stationary reference design the verdict rule detects a 20\% volatility error in every
replication; with the real volatility path of 2018--2022 it does so in 0.45 of them. The median
standard error of the anchored statistic at the boundary is 0.245 in that design against 0.032 in
the stationary one (\cref{tab:e3-power}), plausibly because volatility regimes common to all
tickers make the few evaluation dates the binding constraint (\cref{rem:small-T}). On this panel
and split no generator could have been \GOOD, which was known before any generator was tested in
the full run (\cref{sec:emp-power}). The power check is also one-sided in an economic sense.
Its deviation, a volatility too high, points away from the calmer evaluation period, whereas a generator error
that points towards the new regime makes the generator more like the new data than history is,
and the anchored null does not penalise it.

\paragraph{The critic is not a likelihood estimator.}
Learned scores rank windows but track the exact log-likelihood ratio poorly (Spearman
correlations from 0.024 to 0.556; \cref{tab:e2-llr}), WGAN-GP critics are poor estimators of
$\W$ \citep{mallasto2019well}, and on real data the critic behaved mostly as a realised-volatility
statistic. On the seed-0 windows of the absolute test, a fixed log realised-volatility statistic
reaches between 0.69 and 0.94 of the critic's $\zhac$ and ranks the three likelihood models in the
same order (\cref{tab:e3-rv-baseline}); learning added some precision but no new direction. We did
not run this natural non-learned baseline inside the anchored rule with five seeds and the designs
D1 and D2, so whether the verdicts need a learned critic is not shown.
Where a likelihood exists it is the sharper tool, and the critic statistic should be read as a
test statistic and a witness, not as a distance or a likelihood.

\paragraph{The gap to the oracle.}
The Neyman--Pearson oracle rejects in at least 0.98 of the replications at every independent-window
alternative, whereas every window test stays near its null rate against the smaller jump
deviations (\cref{tab:e2-power}). Part of the gap is structural: with a 50/50 split and equal
sample sizes, a test based on the held-out mean difference of a fixed critic has at most half the
oracle's local noncentrality (\cref{prop:local-power}). The rest is the cost of learning: against
a 5\% volatility increase with 2000 windows the \WC{} test rejects at 0.58, MMD at 0.93. The
\LRC{} and C2ST classifiers selected their first checkpoint, after 50 steps, in almost every
simulated fit (\cref{app:impl-nets}), so their power in \cref{sec:sim} is that of nearly
untrained classifiers; longer training without early stopping was compared only in a pilot. The
power of the critic also depends on its training. A learning rate of $10^{-4}$ raised its power against
a 10\% volatility increase from 0.40 to 0.78 (\cref{tab:e5-simpower}), and the verdicts reported
here are those of the pre-specified multilayer-perceptron critic. With a time-pooled
convolutional critic, examined after the real-data results were known, the five real-data
decisions are unchanged, and in the 20 replications of the realistic design that it shares with
the perceptron its rule power and boundary size equal those of the perceptron (0.45 and 0.20 in
that subset; \cref{tab:e5-arch-real}). In simulation with independent windows it keeps the level and raises
the power against every deviation we studied, for instance from 0.10 to 0.98 against
Student-$t_{10}$ tails and from 0.10 to 0.92 when the synthetic windows lack the jumps of the
real ones (\cref{tab:e5-arch}), which suggests that much of the power gap is architectural,
although the network was chosen after a pilot on some of these designs and its size under the
dependent protocol was not measured. We did not compare with power-optimised deep-kernel tests
\citep{liu2020learning}.

\paragraph{Scope.}
The panel consists of the S\&P~500 constituents at the end of the evaluation period. Both
periods are therefore selected on information from the end of \TEST{}, a look-ahead that goes
beyond survivorship bias, and part of the volatility gap between the generators and the real
\TEST{} returns may be a selection effect; a replication should draw the universe from the
constituents at the end of \TRAIN{} (\cref{sec:data-panel}). One market and one pair of periods
are studied. Univariate verdicts say nothing
about cross-sectional dependence (\cref{ex:corr-blind}). The Markowitz conclusions rest on five
multivariate generators, two test seeds, 60 placebo replications per generator and 30 universes
that share every evaluation date; the utility differences are point estimates without inference
over dates, the portfolios use only the unconditional moments of the generators, and several
planned replications were cut for compute (\cref{app:impl}). The
generators are classical parametric models, without asymmetric GARCH specifications, and
resampling schemes; the generative networks that motivate much of the literature were not evaluated.

\paragraph{Future work.}
Five extensions follow directly. First, the decision rule should also check calibration at the
boundary, for instance by turning \BAD{} into \INCONCL{} when the Clopper--Pearson interval of the
boundary size lies above $\alpha$, by reporting the size-adjusted power of \cref{def:power-check}
beside the nominal one, or by calibrating $\pmed$ against the boundary replications as a Monte
Carlo test (\cref{rem:calibrated-bad}), which is valid only insofar as the boundary design
reproduces the real configuration; a \GOOD{} verdict with a formal error guarantee could come from the
equivalence version of \cref{rem:equivalence}. Second, the dependent-data protocol needs a gap
between training and evaluation blocks, so that \cref{lem:decoupling} applies, together with
small-sample corrections calibrated on placebos that share the dependence of the generator under
test. Third, the rule needs more evaluation dates, from overlapping windows (a step shorter than
$L$, with Driscoll--Kraay inference) or a longer evaluation period, and verdicts should be
aggregated over several split points in a rolling-origin design, so that they describe a
generator rather than one pair of periods. Fourth, the real-data comparison should include a
fixed realised-volatility statistic in the same anchored rule with D1 and D2, generators
recalibrated on a recent window (for instance the last 615 \TRAIN{} days), the stationary
bootstrap, the historical-simulation benchmark of practitioners, tested with the split by source
window (which would also give a real-data placebo at the boundary), per-ticker nulls simulated
from each generator rather than from one pooled \GARCHt{} null, and a horizon-matched likelihood
benchmark such as the predictive log-score of 20-day returns. Fifth, power-optimised critics, architectures that pool over time and
pre-registered optimiser settings should be compared on the designs of \cref{sec:sim}, with the
size of each measured under the dependent protocol, and the protocol applied to neural generators
and to multivariate panels, where the use-based check of \cref{sec:markowitz} should accompany
every distributional verdict.


\bibliographystyle{plainnat}
\bibliography{references}

\clearpage
\appendix
\crefalias{section}{appendix}   % \cref{app:proofs} -> "Appendix A"
% =============================================================================================
% APPENDIX (input after \appendix in main.tex; sections are lettered A, B, ...).
% Target length: 15-25 pages in total. Plan: OUTLINE.md, "Appendix".
% =============================================================================================

\section{Proofs}\label{app:proofs}
% Proofs of Section 4 and of Section 8.1 statements, in the order of the main text, each headed
% "Proof of \cref{...}". Source: docs/theory.md (proof text exists for every proved statement).
% Includes lem:exchangeable (Lemma 3.2), lem:rank (Lemma 3.3), lem:uniform-lip (Lemma 4.1),
% proofs of prop:window-ident, prop:markov-pairs, prop:unbiased, prop:exact-level,
% prop:consistency, cor:w1-limit, prop:bayes-logit, prop:oracle-bound, prop:local-power (sketch),
% prop:anchored (new), prop:paired-boot (sketch), prop:boundary-size, prop:hac, lem:decoupling,
% prop:ticker-bias, prop:verdict, lem:quad-loss, prop:gmv-sensitivity, cor:gmv-drivers(i),
% prop:no-free-lunch, examples 9.5-9.8. Target 6-8 pages.
\todo{Proofs to be written (source docs/theory.md).}

\section{Implementation details}\label{app:impl}
% B.1 critic and classifier architectures, optimiser, standardisation, early stopping (C2ST,
%     LR-critic: early stopping keeps the first checkpoint at n = 500, DESIGN change log);
% B.2 time-block split (labels per sample, random assignment of blocks, no gap), DK variance,
%     Andrews AR(1) plug-in lag, grid of date ranks; fixed-b p-values (E1 comparison rows);
% B.3 wild bootstrap (AR(1) multipliers shared by date, block-length plug-in) and signature
%     features (Chen identity, depth 4);
% B.4 model fitting: Merton multistart MLE, GARCH-t variance targeting and cap, Heston QE + SMM
%     (moments, weights, profiling, bounds), stationary bootstrap block length, MV models;
% B.5 seeds, checkpoints, compute budget actually used (seeds and replication counts as run;
%     list every cut relative to DESIGN Section 4). Target 3-4 pages.
\todo{Implementation details to be written (src/synthval, experiments/README.md).}

\section{Additional simulation results}\label{app:sim}
% tab:e2-size (+ e2_size_note), tab:e2-power-ci, tab:e2dep-size (+ note), tab:e2dep-power
% (+ note), fig:e2-llr-scatter; E2-dep marginal check. Target 3-4 pages.
\todo{Additional simulation tables to be inserted.}

\section{Data details}\label{app:data}
% universe construction, gap rule, the two excluded anomalies (KDP, AMCR) with the rule,
% panel tickers, tab:e3-heston-moments. Target 1-2 pages.
\todo{Data details to be written.}

\section{Additional empirical results}\label{app:emp}
% tab:e1-size-real, tab:e1-nonstat (if not in the main text), tab:e3-perticker-null,
% tab:e3-likelihood-ag. Target 2-3 pages.
\todo{Additional empirical tables to be inserted.}

\section{Markowitz experiment: details}\label{app:mk}
% tab:e4-trtr, tab:e4-utility-agg, tab:e4-placebo (+ note), fig:e4-placebo-size,
% tab:e4-link-pooled. Target 2-3 pages.
\todo{Markowitz details to be inserted.}

\section{Robustness: details}\label{app:robust}
% tab:e5-simpower, tab:e5-calibration (if produced). Target 1-2 pages.
\todo{Robustness details to be inserted. Anchor citations: \citet{hemerik2018exact},
\citet{chwialkowski2014wild}.}


\end{document}
